\documentclass[11pt,reqno]{amsart}
\usepackage[T1]{fontenc}
\usepackage{lmodern,amsmath,amssymb,mathtools,booktabs,array,longtable}
\usepackage[margin=1.05in]{geometry}
\usepackage{microtype}
\usepackage{xcolor}
\usepackage[unicode,colorlinks=true,linkcolor=blue!45!black,citecolor=blue!45!black,urlcolor=blue!45!black]{hyperref}
\usepackage{enumitem,needspace,etoolbox}
\setcounter{tocdepth}{1}
\raggedbottom
\hypersetup{pdftitle={A determinant bound of 498 for the tenth discriminant},pdfauthor={Stephen Chung}}
\numberwithin{equation}{section}
\newtheorem{theorem}{Theorem}[section]
\newtheorem{proposition}[theorem]{Proposition}
\newtheorem{lemma}[theorem]{Lemma}
\newtheorem{corollary}[theorem]{Corollary}
\theoremstyle{remark}
\newtheorem{remark}[theorem]{Remark}
\newcommand{\OO}{\mathcal O}
\newcommand{\ZZ}{\mathbb Z}
\newcommand{\QQ}{\mathbb Q}
\newcommand{\CC}{\mathbb C}
\newcommand{\FF}{\mathbb F}
\newcommand{\NN}{\mathbb Z_{\geq0}}
\newcommand{\ord}{\operatorname{ord}}
\newcommand{\rank}{\operatorname{rank}}
\newcommand{\wt}{\operatorname{wt}}
\newcommand{\N}{\operatorname{N}}
\newcommand{\floor}[1]{\left\lfloor #1\right\rfloor}
\newcommand{\ceil}[1]{\left\lceil #1\right\rceil}
\newcommand{\pos}[1]{\left(#1\right)_+}
\newcommand{\fall}[2]{(#1)_{\underline{#2}}}
\newcommand{\rise}[2]{(#1)_{#2}}
\newcommand{\dd}{\,\mathrm d}
\newcommand{\code}[1]{\texttt{\detokenize{#1}}}
\newcommand{\arith}{\mathcal A}
\newcommand{\metric}{\mathcal U}
\newcommand{\supp}{\operatorname{supp}}
\setlength{\emergencystretch}{2em}
\title[A determinant bound for the tenth discriminant]{A determinant bound of 498 for the tenth discriminant}
\author[Stephen Chung]{Stephen Chung}
\date{Working draft, 28 October 2026}
\thanks{This work is based on an externally verified finding made by autonomous
rStation agents (\url{https://rstation.ai/}). See the section on
\hyperref[sec:ai-use]{AI use} for details.}
\keywords{Imaginary quadratic fields, class number one, interpolation determinants, Frobenius, explicit bounds}
\BeforeBeginEnvironment{lemma}{\Needspace{7\baselineskip}}
\BeforeBeginEnvironment{proposition}{\Needspace{9\baselineskip}}
\BeforeBeginEnvironment{theorem}{\Needspace{7\baselineskip}}
\makeatletter
\def\@setauthors{%
  \vskip 14pt
  \begingroup
  \centering\normalsize
  Stephen Chung\par
  \vskip 4pt
  {\small DualverseAI\par}%
  \endgroup
  \vskip 20pt
}
% Table captions are below the table; place their separation above them.
\setlength{\belowcaptionskip}{0pt}
\pretocmd{\@makecaption}{\ifnum\@tempcnta>63\relax\vskip8pt\fi}{}{}
\makeatother
\begin{document}
\begin{abstract}
We address the challenge of obtaining a small explicit upper bound for a
hypothetical tenth imaginary quadratic discriminant of class number one,
without using Heegner's methods or the known classification. We also make no
quasi-Riemann-hypothesis assumption. We prove the
computer-assisted bound $D\leq498$ for every negative fundamental discriminant
$-D$ with class number one. The proof constructs a nonzero interpolation
determinant from polynomial values and derivatives at algebraic integers of
small absolute value. Inert primes force its norm to have a large divisor,
while a spectral estimate bounds the same norm from above. A rank theorem
allows the determinant to be formed using nearly the minimum number
of points permitted by the dimensions of the matrix. This reduces the upper
bound enough to exclude $D=499$. The numerical
comparisons are evaluated with directed rounding.
\end{abstract}
\maketitle
\section{Introduction}\label{sec:intro}
Let $-D<0$ be a fundamental discriminant, let $F_D=\mathbb Q(\sqrt{-D})$,
and write $\OO_D$ for its ring of integers. Gauss's class number one problem
asks when $\OO_D$ has unique factorization, or equivalently when its ideal
class number $h(-D)$ is one. Heegner~\cite{heegner}, Baker~\cite{baker} and
Stark~\cite{stark} proved that this occurs exactly for
\begingroup
\abovedisplayskip=14pt
\belowdisplayskip=14pt
\abovedisplayshortskip=14pt
\belowdisplayshortskip=14pt
\[
 D\in\{3,4,7,8,11,19,43,67,163\}.
\]
\endgroup
Here we prove a weaker bound by a different method.

\begin{theorem}\label{thm:main}
Let $-D<0$ be a fundamental discriminant with $h(-D)=1$. Then $D\leq498$.
\end{theorem}

The challenge is to obtain as small an explicit upper bound as possible for
a hypothetical tenth discriminant, without relying on Heegner's methods or
the known classification~\cite{tao2026}. We approach this problem through the
determinant method developed in recent work on
Landau--Siegel zeros~\cite{green2026,taosketch}. In particular, the bound in
Theorem~\ref{thm:main} is obtained from determinant inequalities; we do not
exclude the remaining discriminants by a finite test of quadratic residues.

\subsection{Overview of the proof}\label{subsec:overview}
Suppose that $h(-D)=1$ with $D>498$. We will form a matrix by evaluating
polynomials in $(z,w)$ and their $z$-derivatives at pairs $(z,\bar z)$ with
$z\in\OO_D$. If this
matrix has full row rank, it has a nonzero square minor $\Delta$. Since its
entries lie in $\OO_D$, the norm $N=\Delta\bar\Delta=|\Delta|^2$ is a positive
integer. We will bound $N$ from above using the absolute values of the
entries, and show that it is divisible by an integer larger than this upper
bound. This contradiction will exclude the proposed value of $D$. The
choice of polynomials and evaluation points must therefore serve both
estimates while ensuring that the matrix has full row rank.

To see how the class number assumption enters, suppose that a prime $p$ is
not inert. There is then an ideal of norm $p$, which must be principal because
$h(-D)=1$. Its generator has norm $p$ and cannot be a rational integer. But
every non-rational element of $\OO_D$ has norm at least $D/4$, so every prime
$p<D/4$ must be inert. At such a prime, conjugation agrees with Frobenius
modulo $p$, giving $z^p\equiv\bar z\pmod{p\OO_D}$. Thus the polynomial
$z^p-w$ takes values divisible by $p$ at every point $(z,\bar z)$ with
$z\in\OO_D$. By expressing the rows of our matrix in terms of powers of
$z^p-w$, we obtain prime powers dividing $N$.

This change of basis must remain within the chosen polynomial space. We use
the monomials $z^rw^s$ with $r+Ks<R$, where $K$ and $R$ are parameters. For
$p\geq K$, replacing $z^p$ by $w$ does not increase $r+Ks$, so the triangular
space of exponents permits all the required row operations. Taking several derivatives at each point gives more columns from the same
set of lattice points. Differentiation reduces the resulting divisibility,
but by an amount that can be bounded explicitly in terms of the derivative
order. We also use inert primes above
$D/4$: the same norm argument bounds the number of non-inert primes in an
interval by the number of elements of $\OO_D$ whose norms lie there.

For the upper bound, we choose the $n$ lattice points of $\OO_D$ nearest to
the origin. Reducing $n$ reduces the radius of the disk containing them and
hence the bounds on the polynomial values. The estimates must also account
for the derivative columns at each point. These columns are far from
orthogonal, so bounding each separately loses too much at the smaller
values of $D$. Instead, we estimate their Gram matrix, retaining the saving
from these correlations before combining the contributions of the
different points.

It remains to determine how few points we can use without losing full row
rank. If the polynomial space has dimension $M$ and we take $J$ derivative
orders at each point, there are $nJ$ columns, so necessarily $nJ\geq M$.
We prove a sufficient condition close to this necessary one. Direct
reduction modulo an auxiliary inert prime $p$ loses information because
$z^p-w$ vanishes at every reduced node. We therefore rescale both the
polynomial rows and the derivative columns using their known divisibility,
obtaining an integral matrix whose reduction retains this information.
A hypothetical row relation then gives a polynomial in two variables with
prescribed local vanishing. Applying a differential operator and taking a
determinant of coefficients produces a nonzero polynomial in one variable
whose zeros, counted with multiplicity, exceed its degree. This contradiction
permits a smaller $n$ and therefore improves the upper bound for $N$.

We apply these estimates as follows. Elementary norm arguments first reduce
$D>498$ to primes congruent to $3$ modulo $8$. A uniform estimate and $101$
adjacent interval comparisons exclude all $D\geq875$, leaving sixteen
possible values below $875$. Fifteen are excluded by fixed determinant
comparisons; the last is $499$. Table~\ref{tab:opening} shows the effect of
the smaller node count there, with the polynomial space and derivative
orders fixed. After the derivative-order contributions have been combined,
the arithmetic term changes only slightly, whereas the size estimate
decreases enough to make the normalized difference positive.

\begin{table}[ht]
\centering
\begin{tabular}{lrr}
\toprule
 & Larger node count & Smaller node count\\
\midrule
Nodes $n$ & $41\,793\,600$ & $35\,394\,400$\\
Proportion $M/(nJ)$ & $0.83333\ldots$ & $0.98399\ldots$\\
Arithmetic minus size bound & $-0.12343\ldots$ & $+0.05119\ldots$\\
\bottomrule
\end{tabular}
\caption{The comparison at $D=499$, with the polynomial space and derivative
orders fixed (approximate values). The directed bounds are given in
Section~\ref{sec:completion}.}
\label{tab:opening}
\end{table}

Section~\ref{sec:comparison} introduces the determinant and the inequality
we must contradict. Sections~\ref{sec:arithmetic} and~\ref{sec:metric} develop
its divisibility and size estimates. Section~\ref{sec:rank} proves the rank
criterion, and Section~\ref{sec:completion} completes the numerical argument.
The explicit arithmetic formulas, subsidiary spectral proofs, uniform
estimates for $D\geq875$, and numerical calculations are collected in the
accompanying repository,
\url{https://github.com/dualverse-ai/tenth_discriminant}.

\section{The determinant comparison}\label{sec:comparison}

\subsection{A model: an integer that is small and highly divisible}
\label{subsec:model}

Every nonzero integer $N$ satisfies
\begin{equation}\label{eq:product-formula}
  \log|N|=\sum_p v_p(N)\log p .
\end{equation}
Our integer will be $N=|\Delta|^2$ for a nonzero algebraic integer $\Delta\in
\mathcal{O}_D$, which is the norm of $\Delta$ and hence a positive rational
integer. The strategy is to bound $\log N$ from above by Archimedean
estimates, to bound $v_p(N)$ from below at many primes, and to show that the
two bounds are incompatible when $h(-D)=1$. This subsection illustrates the comparison using elementary estimates. The rest of the paper sharpens each of its
three ingredients.

\subsubsection*{The matrix}
We work with polynomials in two independent variables $z$ and $w$. The variable
$w$ will be specialized to the complex conjugate of the value of $z$, but it is
treated as independent of $z$ throughout. For $P\in\mathbb{Z}[z,w]$ define the
Hasse derivatives in $z$, with $w$ held fixed, by
\[
  P(z+t,w)=\sum_{j\ge0}\bigl(\partial_z^{[j]}P\bigr)(z,w)\,t^j,
  \qquad
  \partial_z^{[j]}\bigl(z^rw^s\bigr)=\binom rj z^{r-j}w^s .
\]
Hasse derivatives keep integer coefficients; ordinary derivatives would
introduce a factor $j!$. Fix a positive integer $J$. Let $z_1,\dots,z_n\in\mathcal{O}_D$ be
distinct \emph{nodes}, and let $\Lambda\subset\mathbb{Z}_{\ge0}^2$ be a finite
set of exponents. The \emph{jet matrix} has rows indexed by $(r,s)\in\Lambda$,
columns indexed by pairs $(i,j)$ with $1\le i\le n$ and $0\le j<J$, and entries
\begin{equation}\label{eq:jet-matrix}
  V_{(r,s),(i,j)}
  =\bigl(\partial_z^{[j]}z^rw^s\bigr)(z_i,\overline{z_i})
  =\binom rjz_i^{\,r-j}\,\overline{z_i}^{\,s}\in\mathcal{O}_D .
\end{equation}
For a set $C$ of $|\Lambda|$ columns, write $\Delta_C$ for the corresponding
maximal minor. More generally, the row of any polynomial $P$ in the span of the
monomials $z^rw^s$, $(r,s)\in\Lambda$, is the vector of its jets
$(\partial_z^{[j]}P)(z_i,\overline{z_i})$. Integral unitriangular
changes of the spanning polynomials leave every $\Delta_C$ unchanged.

\subsubsection*{One Frobenius calculation}
The divisibility comes from the behaviour of conjugation modulo an inert prime.

\begin{lemma}\label{lem:frobenius}
Let $p$ be a prime that is inert in $F_D$, and let $z\in
\mathcal{O}_D$. Then $z^p\equiv\overline z\pmod{p\mathcal{O}_D}$.
Consequently, if $g=z^p-w$, then
\[
  g(z+t,\overline z)=pU(t)+t^p
\]
for a polynomial $U\in\mathcal{O}_D[t]$ of degree less than $p$, and for all
integers $m,j\ge0$,
\begin{equation}\label{eq:layer-divisibility}
  p^{\lfloor j/p\rfloor}\,[t^j]\,g(z+t,\overline z)^m\in
  p^m\mathcal{O}_D .
\end{equation}
\end{lemma}

\begin{proof}
The quotient $\OO_D/p\OO_D$ is the field of $p^2$ elements. Conjugation
induces its nontrivial automorphism, which is the $p$th-power map. This gives
$z^p\equiv\bar z\pmod{p\OO_D}$, including when $p=2$. Next,
\[
  (z+t)^p-\overline z
  =(z^p-\overline z)+\sum_{k=1}^{p-1}\binom pkz^{p-k}t^k+t^p,
\]
and every term except $t^p$ is divisible by $p$. Finally,
$(pU+t^p)^m=\sum_{q=0}^m\binom mq p^{m-q}U^{m-q}t^{pq}$, and a term contributes
to $t^j$ only if $q\le\lfloor j/p\rfloor$. Each such term is divisible by
$p^{m-\lfloor j/p\rfloor}$, which proves~\eqref{eq:layer-divisibility}.
\end{proof}

Geometrically, the lemma says that, $p$-adically, every point
$(z,\overline z)$ lies close to the curve $z^p=w$~\cite{taosketch}.
The polynomial $g$ is small at every node, and its $z$-derivatives are small
too, since $\partial_z g=pz^{p-1}$. This is why the jets differentiate only in
$z$: a derivative in $w$ would not preserve the divisibility.

\subsubsection*{Where the triangle comes from}
To use Lemma~\ref{lem:frobenius}, write each row exponent as $r=mp+a$ with
$0\le a<p$ and replace the monomial $z^rw^s$ by
\begin{equation}\label{eq:frobenius-basis}
  g^m z^aw^s=z^rw^s+\sum_{i=1}^m(-1)^i\binom mi z^{r-pi}w^{s+i}.
\end{equation}
By the Leibniz rule and~\eqref{eq:layer-divisibility}, every jet of
$g^mz^aw^s$ of order $j$, multiplied by $p^{\lfloor j/p\rfloor}$, lies in
$p^m\mathcal{O}_D$. The change of basis~\eqref{eq:frobenius-basis} is
unitriangular for the ordering by $z$-degree, but only if every exponent
$(r-pi,s+i)$ on the right is again in $\Lambda$. The support must therefore be
closed under the Frobenius substitution $z^p\mapsto w$.

The rectangle $\{r<R,\ s<S\}$ used by Green is not closed under this
substitution: it is stopped by the ceiling $s<S$. There the substitution can be
applied only $\min(\lfloor r/p\rfloor,S-1-s)$ times, and the $p$-adic credit of
the row is truncated accordingly~\cite[Lemma~5.1]{green2026}. The remedy is to
grade the monomials so that the substitution never raises the degree. Fix integers $K\geq2$ and $S\geq1$. Give $z$
weight $1$ and $w$ weight $K$, and take
\begin{equation}\label{eq:triangle}
  \Lambda=\{(r,s)\in\mathbb{Z}_{\ge0}^2:\ r+Ks<R\},\qquad R=KS,\qquad
  M=|\Lambda|=\frac{R(S+1)}{2}.
\end{equation}
The substitution $z^p\mapsto w$ changes the weight by $K-p$, so it never leaves
$\Lambda$ when $p\ge K$. Every row then receives its full credit
$\lfloor r/p\rfloor$. Primes $p<K$ receive truncated credit, and that boundary
is treated in Section~\ref{sec:arithmetic}.

Scaling column $(i,j)$ by $p^{\lfloor j/p\rfloor}$ multiplies $\Delta_C$ by
$p^{\sum_{(i,j)\in C}\lfloor j/p\rfloor}$. Since $p\mathcal{O}_D$ is prime and
$v_p(|\Delta_C|^2)=2v_{p\mathcal{O}_D}(\Delta_C)$, we obtain, for every inert
prime $p\geq K$,
\begin{equation}\label{eq:model-valuation}
  v_p\bigl(|\Delta_C|^2\bigr)\ \ge\
  2\sum_{(r,s)\in\Lambda}\Bigl\lfloor\frac rp\Bigr\rfloor
  \;-\;2\sum_{(i,j)\in C}\Bigl\lfloor\frac jp\Bigr\rfloor .
\end{equation}
The row contribution, including its factor $2$, is about $W/p$, where
$W=R^2S/3\approx2\sum_\Lambda r$. The second is a column loss that is present
only when derivative orders reach $p$. In our application $J$ is much larger
than $p$, so the loss is real. Section~\ref{sec:arithmetic} shows how it
combines with the factorial sizes of the Hasse coefficients.

\subsubsection*{Sizes}
If every node satisfies $|z_i|\le\rho$ with $\rho\ge1$, then each entry
of row $(r,s)$ has absolute value at most $\binom rj\rho^{r+s-j}$. Hadamard's
inequality therefore gives
\begin{equation}\label{eq:model-hadamard}
  \log|\Delta_C|^2\ \le\ \sum_{(r,s)\in\Lambda}2(r+s)\log\rho
  \;+\;\sum_{(r,s)\in\Lambda}\log\sum_{(i,j)\in C}\binom rj^2\rho^{-2j}.
\end{equation}
The first term is $(1+1/K+O(1/S))\,W\log\rho$. Combining
\eqref{eq:product-formula}, \eqref{eq:model-valuation} and
\eqref{eq:model-hadamard}, any nonzero $\Delta_C$ satisfies
\begin{equation}\label{eq:model-inequality}
  \sum_{\substack{p\geq K\\ p\text{ inert}}}2\Bigl(\sum_{\Lambda}
  \Bigl\lfloor\frac rp\Bigr\rfloor-\sum_{C}\Bigl\lfloor\frac jp\Bigr\rfloor
  \Bigr)\log p
  \le\ \sum_{\Lambda}2(r+s)\log\rho
  +\sum_{\Lambda}\log\sum_{C}\binom rj^2\rho^{-2j}.
\end{equation}
After division by $W$, the left side is roughly the sum of $(\log p)/p$ over
inert primes $K\leq p<R$, less the column loss. Under $h(-D)=1$ this inert sum is
bounded below by the corresponding sum over all primes, minus a loss controlled
by counting values of the norm form, as in Section~\ref{sec:arithmetic}. The
right side is roughly $(1+1/K)\log\rho$ plus the derivative term. The theorem
follows when the left side is larger.

This model already shows where the difficulty lies, and it points to the three
refinements. First, the arithmetic credit must be computed exactly, including
small primes, prime powers, and the loss from non-inert primes, bounded through
the norm form (Section~\ref{sec:arithmetic}). Second, Hadamard's inequality treats the
entries of each node block as unrelated and discards their cancellation. The
spectral estimates of Section~\ref{sec:metric} recover part of it. Because the
nonzero minor is not identified, those estimates must hold for every column set
$C$. Third, $\rho$ is governed by the node count. The $n$ nearest lattice points
lie in a disk of radius about $\sqrt{n\sqrt D/(2\pi)}$
(Lemma~\ref{lem:node-disk}), and $n\ge M/J$ is forced.

\subsection{Arithmetic of class number one}\label{subsec:arith}

Every element of $\mathcal{O}_D$ has the form $\alpha=(x+y\sqrt{-D})/2$ with
$x,y\in\mathbb{Z}$ and $x\equiv yD\pmod 2$, and its norm is
\begin{equation}\label{eq:norm}
  N(\alpha)=|\alpha|^2=\frac{x^2+Dy^2}{4}.
\end{equation}
The elementary reductions below follow from one observation about norms.

\begin{lemma}[Norm obstruction]\label{lem:norm-obstruction}
If $\alpha\in\mathcal{O}_D\setminus\mathbb{Z}$, then $N(\alpha)\ge D/4$. If
$\alpha\in\mathbb{Z}$, then $N(\alpha)$ is a square.
\end{lemma}

\begin{proof}
If $\alpha\notin\mathbb{Z}$ then $y\ne0$ in~\eqref{eq:norm}, so $N(\alpha)\ge
Dy^2/4\ge D/4$. If $\alpha\in\mathbb{Z}$ then $N(\alpha)=\alpha^2$.
\end{proof}

\begin{lemma}\label{lem:small-inert}
Suppose $h(-D)=1$. Then every prime $p<D/4$ is inert in $F_D$.
\end{lemma}

\begin{proof}
If $p$ is split or ramified, there is a prime ideal of norm $p$. It is
principal, generated by some $\alpha$ with $N(\alpha)=p$. Since $p$ is not a
square, $\alpha\notin\mathbb{Z}$, and Lemma~\ref{lem:norm-obstruction} gives
$p\ge D/4$.
\end{proof}

\begin{lemma}[Reduction]\label{lem:reduction}
Suppose $h(-D)=1$ and $D>16$. Then $D$ is a prime and $D\equiv3\pmod 8$.
\end{lemma}

\begin{proof}
The ramified primes are the prime divisors of $D$, so by
Lemma~\ref{lem:small-inert} every prime divisor of $D$ is at least $D/4$. If
$D$ is even, then $2$ ramifies and $2<D/4$, which is impossible. Hence $D$ is
odd, and since $-D$ is a fundamental discriminant, $D$ is squarefree with
$D\equiv3\pmod4$. If $D$ were composite, its least prime factor $q$ would
satisfy $q\le\sqrt D<D/4$, again impossible. So $D$ is prime. Finally, if
$D\equiv7\pmod8$ then $-D\equiv1\pmod 8$ and $2$ splits, contradicting
Lemma~\ref{lem:small-inert}. Therefore $D\equiv3\pmod8$.
\end{proof}

For the proof of Theorem~\ref{thm:main} we assume
$D>498$ from now on, so $D$ is a prime congruent to $3$ modulo $8$, and both $2$
and every odd prime below $D/4$ are inert. Since $D\equiv3\pmod4$, the
integer cutoff $p<D/4$ is equivalently $p\leq(D-3)/4$. We write
$\omega=(1+\sqrt{-D})/2$, so that $\OO_D=\ZZ[\omega]$. In these coordinates
$\N(x_0+y\omega)=(x_0+y/2)^2+Dy^2/4$; it is the same norm form as
\eqref{eq:norm}, with $x=2x_0+y$.

\subsection{Nearest nodes}\label{subsec:geometry}

Our nodes are the lattice points of $\mathcal{O}_D\subset\mathbb{C}$ nearest to
the origin. Let $z_1,z_2,\dots$ be the elements of $\mathcal{O}_D$
listed with $|z_1|\le|z_2|\le\cdots$, with ties broken arbitrarily.
By~\eqref{eq:norm}, $\mathcal{O}_D$ is the union of the horizontal rows
\[
  \Bigl\{\tfrac{x}{2}+\tfrac{y}{2}\sqrt{-D}:\ x\equiv yD\ (\mathrm{mod}\ 2)\Bigr\},
  \qquad y\in\mathbb{Z}.
\]
The rows are spaced $\sqrt D/2$ apart vertically, and within each row the
points form a translate of $\mathbb{Z}$. The covolume of $\mathcal{O}_D$ is
therefore $\sqrt D/2$.

\begin{lemma}[Nearest-node disk]\label{lem:node-disk}
Put $\delta=\sqrt{D+4}/4$. For every $i\ge1$,
\begin{equation}\label{eq:node-disk}
  |z_i|\ \le\ \rho_i:=\sqrt{\frac{i\sqrt D}{2\pi}}+\delta .
\end{equation}
\end{lemma}

\begin{proof}
Every point of the plane is within vertical distance $\sqrt D/4$ of a
lattice row, and within horizontal distance $1/2$ of a point on that row.
Thus the lattice has covering radius at most $\delta$.
For $z\in\mathcal{O}_D$, let $V(z)$ be its Voronoi cell, the set of
points of $\mathbb{C}$ at least as close to $z$ as to any other lattice
point. These cells are translates of one another, they cover $\mathbb{C}$, and
they overlap only on their boundaries. Each has area $\sqrt D/2$. Every point of $V(z)$ lies within $\delta$ of $z$.

Let $\rho>\delta$. If $|\zeta|\le\rho-\delta$ and $\zeta\in V(z)$, then
$|z|\le|\zeta|+\delta\le\rho$. Hence the closed disk of radius
$\rho-\delta$ is covered by the cells of lattice points of modulus at most
$\rho$, and comparing areas gives
\[
  \#\{z\in\mathcal{O}_D:\ |z|\le\rho\}\ \ge\
  \frac{2\pi(\rho-\delta)^2}{\sqrt D}.
\]
With $\rho=\rho_i$ the right side equals $i$, so at least $i$ lattice points
have modulus at most $\rho_i$. This gives $|z_i|\le\rho_i$.
\end{proof}

We use Lemma~\ref{lem:node-disk} in two ways. It supplies the radius
$\rho=\rho_n$ in the Archimedean estimates. It also controls residues, through
the following consequence.

\begin{corollary}[Distinct residues]\label{cor:distinct-residues}
Let $p$ be a prime with $2\rho_n<p$. Then $z_1,\dots,z_n$ have
distinct residues modulo $p\mathcal{O}_D$.
\end{corollary}

\begin{proof}
A nonzero element $p\beta$ of $p\mathcal{O}_D$ has modulus
$p\sqrt{N(\beta)}\ge p$. The difference of two distinct nodes is nonzero and
has modulus at most $2\rho_n<p$.
\end{proof}

At $D=499$, with $n=35\,394\,400$, Lemma~\ref{lem:node-disk} gives
$2\rho_n<22\,491.5$. This is far below every candidate auxiliary prime in
$(160\,000,170\,000]$, so the residue hypothesis of the rank theorem holds
whichever of those primes is inert.


\subsection{The inequality to be proved}
\label{subsec:central}
We now give the comparison in a form that separates the algebra from the numerical estimates. The notation introduced here will be used throughout the paper. Write
\begin{equation}\label{eq:normalization}
 W=\frac{R^2S}{3},\qquad
 W_1=2\sum_{(r,s)\in\Lambda}r,\qquad
 W_2=2\sum_{(r,s)\in\Lambda}s.
\end{equation}
The factor $W$ is a convenient normalization: the main logarithmic terms in both estimates have size $W\log D$. Direct summation gives
\begin{equation}\label{eq:radius-coefficient}
 \frac{W_1+W_2}{W}
 =1+\frac1K+\frac{3(K-1)}{2KS}+\frac{K-5}{2KS^2}.
\end{equation}

Choose $M$ columns $C$ for which the square submatrix $V_C$ is nonsingular, and put $\Delta_C=\det V_C$. The rank argument will supply such a choice; it will not specify the columns. The arithmetic estimates therefore have to hold for every possible choice $C$. They take the form
\begin{equation}\label{eq:arithmetic-interface}
 W\arith-2\sum_{(i,j)\in C} b(j)\ \leq\ \log|\Delta_C|^2.
\end{equation}
Here $\arith$ is a column-independent lower bound assembled from local divisibilities, and $b(j)$ is the loss from using a derivative of order $j$. Both depend on the chosen parameters. Section~\ref{sec:arithmetic} describes their construction and states the local estimate used below.

The main part of $b(j)$ is the first-prime-power sum
$\sum_q\lfloor j/q\rfloor\log q$. Legendre's formula shows that subtracting
it from $\log(j!)$ leaves higher-prime-power savings, with additional
contributions from non-inert primes. This remaining function is
nearly linear in $j$; we subtract a slope and bound its finite remainder. Set
\begingroup\postdisplaypenalty=10000
\begin{equation}\label{eq:F-H}
 F(j)=\log(j!)-b(j),\qquad H(j)=j\sigma-F(j),
\end{equation}
\endgroup
where $\sigma$ is a slope chosen in the arithmetic calculation. Let $c$ be any bound such that
\begin{equation}\label{eq:column-remainder}
 2\sum_{(i,j)\in C}H(j)\leq Mc
 \quad\hbox{for every }|C|=M.
\end{equation}
There are only $n$ columns of any one order $j$, so this is a finite selection problem with a capacity constraint. It need not be bounded by the much cruder $2M\max_jH(j)$.

Let $\rho$ contain all the nodes, put $u_i=z_i/\rho$, and define
\begin{equation}\label{eq:residual-matrix}
 B_{(r,s),(i,j)}(\ell)
 =\frac{\fall rj}{R^j}e^{\ell j}u_i^{r-j}\bar u_i^{s},
 \qquad \ell=\log(R/\rho)-\sigma.
\end{equation}
Here $\fall rj=r(r-1)\cdots(r-j+1)$ is the falling factorial, with
$\fall r0=1$; the entry is zero when $j>r$. The exact relation to the integral matrix is
\begin{equation}\label{eq:diagonal-scaling}
 B_C(\ell)
 =\operatorname{diag}_{\Lambda}(\rho^{-r-s})\,
 V_C\,
 \operatorname{diag}_{C}\bigl(j!(\rho/R)^je^{\ell j}\bigr).
\end{equation}
Taking squared absolute determinants, substituting \eqref{eq:F-H}, and then using \eqref{eq:arithmetic-interface} gives
\[
 W\arith\leq (W_1+W_2)\log\rho
       +\log|\det B_C(\ell)|^2+2\sum_C H(j).
\]
This calculation explains the normalization: the factorials are cancelled exactly, rather than estimated twice.

\begin{proposition}[Determinant comparison]\label{prop:comparison}
Suppose that $V$ has full row rank, that \eqref{eq:arithmetic-interface} and \eqref{eq:column-remainder} hold, and that
\[
 \frac1W\log|\det B_C(\ell)|^2\leq U_{\rm var}
 \qquad (|C|=M)
\]
holds uniformly over all column selections. Then
\begin{equation}\label{eq:master}
 \arith\leq \frac{W_1+W_2}{W}\log\rho
                  +U_{\rm var}+\frac{Mc}{W}.
\end{equation}
In particular, a strictly larger arithmetic lower bound is impossible.
\end{proposition}
\begin{proof}
Choose a nonsingular minor and apply the preceding identities. All bounds are uniform in the choice, so no information about which minor is nonsingular is needed.
\end{proof}

For the numerical comparisons we use the slightly weaker bound obtained
by adding $2M\log M/W$ to the right side of~\eqref{eq:master}.
This positive term is less than $3\times10^{-7}$ at the parameters for $499$.
Thus define
\begin{equation}\label{eq:metric-total}
 \metric=\frac{W_1+W_2}{W}\log\rho+U_{\rm var}
             +\frac{Mc+2M\log M}{W},
 \qquad G=\arith-\metric.
\end{equation}
A positive lower bound for $G$ rules out the proposed field. There are three separate tasks: prove full rank, establish the local arithmetic estimates, and bound the same minor at the complex embedding. The auxiliary-prime scalings used to prove rank in Section~\ref{sec:rank} do not enter \eqref{eq:diagonal-scaling}; after nonvanishing has been established, all size and divisibility estimates concern the original integral minor $\Delta_C$.

It is already possible to see why the number of nodes matters. The disk estimate gives $\rho\asymp n^{1/2}D^{1/4}$ once its boundary error is small. Reducing $n$ therefore lowers the common term in \eqref{eq:metric-total} by approximately
\[
 \frac12\left(1+\frac1K\right)\log\frac{n_{\rm old}}{n_{\rm new}}.
\]
At the two node counts in Table~\ref{tab:499}, this is about $0.085$. The remaining terms also change: there are fewer copies of each derivative order, and the proportion $\theta=M/(nJ)$ of available columns that must be selected increases. We call this proportion the \emph{occupancy}. Section~\ref{sec:metric} exploits these restrictions. In the actual calculation the combined gain is enough to change the sign of $G$.


\section{Divisibility of the minor}\label{sec:arithmetic}
We write $q$ for a general prime, reserving $p$ for the auxiliary prime in
Section~\ref{sec:rank}. The basic Frobenius argument in Section~\ref{subsec:model} gives a row
contribution from each inert prime and a loss depending on the derivative
orders selected. We now strengthen the row contribution, first by using the
conjugate variable and then by retaining the small primes. Finally we explain
how class number one controls the losses from non-inert primes.

Recall that for an inert prime $q\geq K$, the integral row basis is
$P_{r,s}=z^{r\bmod q}(z^q-w)^{\lfloor r/q\rfloor}w^s$.
Writing $v_q$ for the valuation normalized by $v_q(q)=1$, the calculation
in Lemma~\ref{lem:frobenius} gives
\begin{equation}\label{eq:basic-valuation}
 v_q(\partial_z^{[j]}P_{r,s}(z_i,\bar z_i))
 \geq\max(0,\lfloor r/q\rfloor-\lfloor j/q\rfloor).
\end{equation}
Each determinant term uses every row and every selected column once. Hence
\begin{equation}\label{eq:basic-minor-valuation}
 v_q(\Delta_C)\geq\sum_\Lambda\lfloor r/q\rfloor
                  -\sum_C\lfloor j/q\rfloor.
\end{equation}
The right side can be negative, in which case it is merely a weak bound.
Its usefulness emerges after summing over primes: the column loss combines
with the factorials in the Hasse derivatives, as in \eqref{eq:F-H}.

\subsection{Additional divisibility in the conjugate variable}
Differentiation in $z$ leaves polynomials in $w$ untouched. Such polynomials can contribute additional divisibility at no derivative cost. We use the following elementary tower. Define
\[
 T_1(w)=w^{q^2}-w,\qquad
 \nu_a=\frac{q^{2a}-1}{q^2-1},\qquad
 T_{a+1}(w)=T_a(w)^{q^2}-q^{\nu_a(q^2-1)}T_a(w).
\]
These are monic polynomials of degrees $q^{2a}$. Inductively,
\begin{equation}\label{eq:tower-value}
 T_a(\bar z_i)\in q^{\nu_a}\OO_D.
\end{equation}
Indeed, if $T_a(\bar z_i)=q^{\nu_a}b$, then the next value equals
$q^{q^2\nu_a}(b^{q^2}-b)$, which is divisible by
$q^{q^2\nu_a+1}=q^{\nu_{a+1}}$.

Write $s=\sum_{a\geq0}d_aq^{2a}$ in base $q^2$, and put
\[
 A_s(w)=w^{d_0}\prod_{a\geq1}T_a(w)^{d_a}.
\]
It is monic of degree $s$, and \eqref{eq:tower-value} gives
\[
 v_q(A_s(\bar z_i))\geq
 \sum_{a\geq1}d_a\nu_a
 =\sum_{a\geq1}\floor{s/q^{2a}}.
\]
Replacing $w^s$ by $A_s(w)$ in the previous basis preserves triangularity. We obtain the stronger row estimate
\begin{equation}\label{eq:tower-valuation}
 v_q\bigl(\partial_z^{[j]}(z^{r\bmod q}(z^q-w)^{\floor{r/q}}A_s(w))
                 (z_i,\bar z_i)\bigr)
 \geq\max(0,\floor{r/q}-\floor{j/q})
         +\sum_{a\geq1}\floor{s/q^{2a}}.
\end{equation}
This is a local statement for any nodes in the order; it has no rank hypothesis.

The row contributions are finite sums over the triangle. Their exact
floor-sum formulas are given in the supplementary notes, Section~\ref{supp:rows}.

\subsection{Small primes and non-inert primes}
For an inert prime $q<K$, replacing $z^q$ by $w$ increases weighted degree. The number of replacements in a row is therefore bounded by
\begin{equation}\label{eq:small-h}
 h_q(r,s)=\min\left(\floor{r/q},
          \floor{(R-1-r-Ks)/(K-q)}\right).
\end{equation}
There is still a substantial contribution: summing this function over the triangle yields
\[
 2\sum_{\Lambda}h_q(r,s)=W_2=\frac{R(S^2-1)}3.
\]
One way to verify the identity is to sum the sizes of its superlevel sets. The condition $h_q\geq a$ is equivalent, after writing $r=qa+r'$, to
$r'+Ks<R-Ka$, and its number of solutions is $K(S-a)(S-a+1)/2$. Summing for $1\leq a<S$ gives the displayed value.

The small-prime estimate used in the computation improves this elementary bound by combining $z^q-w$, $w^q-z$, and $z^{q^2}-z$. The first factor can lose a power of $q$ only after spending at least $q$ orders of differentiation; the latter two can lose one after spending one order. Optimizing where each type is used is a finite triangular row calculation. The resulting finite lower bound is recorded in the supplementary notes, Section~\ref{app:arithmetic}, together with its cut parameter. In particular, $q=2$ is included. Working in the maximal order is essential here, since $2$ is inert when $D\equiv3\pmod8$.

The arithmetic calculation also uses primes above $(D-3)/4$, where inertness is no longer guaranteed. Class number one still supplies useful information. If $q$ splits or ramifies, an ideal of norm $q$ is principal, so
\begin{equation}\label{eq:prime-norm}
 q=\N(x_0+y\omega)=\left(x_0+\frac y2\right)^2+\frac D4y^2,
 \qquad y\ne0.
\end{equation}
At a fixed norm cutoff there are only finitely many such pairs. One can therefore start with an all-prime estimate and subtract an upper bound for the contribution of possible non-inert primes, using these norm representations. This deduction uses principality; it does not test the quadratic character at a particular discriminant.

Counting all norm values as possible primes overestimates the loss. Unique factorization allows some of those values to be removed in advance. Products of two or three nonrational factors have composite norm, and the corresponding primitive products can be counted with their multiplicities controlled. The product counts used for this subtraction are recorded in the supplementary arithmetic notes. At the finite comparisons in Section~\ref{sec:completion}, only the two- and three-factor corrections are used. The upper bound for possible losses is decreased by a rigorously bounded contribution of these composite values. Thus this correction \emph{increases} $\arith$; it is not an additional subtraction of prime divisibility.

For orientation, the resulting arithmetic term has the form
\begin{equation}\label{eq:arithmetic-parts}
 \arith=A_{\rm all}-L_{\rm norm}+C_2+C_3.
\end{equation}
The first term sums inert-prime row credits as a provisional all-prime bound; $L_{\rm norm}$ pays for norm values that might be non-inert primes; $C_2$ and $C_3$ restore credits paid unnecessarily for certified composite values. Split and ramified primes also have their own local conjugate-variable contributions. These, and the derivative-order refunds, are included in the arithmetic lower bound. This matters numerically near $D=499$.

\subsection{The arithmetic estimate used below}
We use the following consequence of the local estimates, in a form that
continues to hold when the number of nodes is reduced. The normalization is $\log|\Delta_C|^2/W$: all prime-sum and
norm-loss terms in $\arith$ are divided by $W$, while $b(j)$ is the column
debit in logarithmic units before this division.

\begin{proposition}[Local arithmetic estimate]\label{prop:arithmetic}
Let $-D$ be a fundamental discriminant with $h(-D)=1$ at one of the
finite comparisons below $875$ in Section~\ref{sec:completion}, with the
parameters specified there. Let $z_1,\dots,z_n$ be the $n$ nearest elements
of $\OO_D$, and let $V$ be the matrix~\eqref{eq:jet-matrix} with
$0\leq j<J$ and $\theta=M/(nJ)\leq1$. Define $\arith,\sigma,c,b(j)$ by
the local arithmetic formulas in the supplementary notes,
Section~\ref{app:arithmetic}, including the norm-product corrections
and derivative-order refunds.
Then~\eqref{eq:arithmetic-interface} and~\eqref{eq:column-remainder} hold for
every set $C$ of $M$ columns with $\Delta_C\ne0$.
\end{proposition}

Here $\Delta_C$ is a minor of the original integral Hasse matrix $V$.
The bound on $D$ will follow by comparing these local estimates with the size
estimate of Section~\ref{sec:metric}. The preceding arguments explain the
constructions behind the estimates. The supplementary notes give the finite formulas and their derivations. Every quantity
depending on $n$ is evaluated at the node count in use. This includes the
number of columns of each order, and hence the least possible sum of their
first-power floors and the largest possible remainder sum.

The higher-prime-power saving can be seen directly from Legendre's formula:
\[
 \log(j!)=\sum_q\sum_{a\geq1}\floor{j/q^a}\log q.
\]
The debit $\sum_q\floor{j/q}\log q$ from \eqref{eq:basic-minor-valuation} cancels its $a=1$ part exactly. The remaining terms, together with the split-prime refunds, form $F(j)$ in \eqref{eq:F-H}. Sorting $H(0),\ldots,H(J-1)$ in decreasing order, say as $H_{(1)},\ldots,H_{(J)}$, and writing $M=nt+r$ with $0\leq r<n$, yields the exact capacity bound
\begin{equation}\label{eq:sorted-remainder}
 \sum_C H(j)\leq n\sum_{a=1}^{t}H_{(a)}+rH_{(t+1)}.
\end{equation}
When $r=0$ the final term is omitted. The right side gives an admissible $Mc/2$. Grouping the floor functions into intervals on which their quotients are
constant gives a finite evaluation of the right side. The resulting bound
for $c$ includes the remainder at every derivative order.

\section{An upper bound for every column selection}\label{sec:metric}
We turn to the complex absolute value of the minor. An entrywise bound followed by Hadamard's inequality is enough for a large discriminant bound. At the smaller values of $D$ considered here, however, it loses too much information about the derivative columns at a single node. Those columns are far from orthogonal. Our estimate keeps them together, bounds their Gram matrix, and then combines the node contributions.

There is a complication worth making explicit. Full row rank gives a nonsingular minor but does not tell us how many of its columns come from each node, or which derivative orders are selected. Neither the first nor the last $m$ orders need maximize an $m$-column minor. The estimates below therefore allow an arbitrary subset of orders at each node, and only use the restrictions that there are at most $J$ columns per node and $M$ columns in total.

\subsection{Weights and node Gram matrices}
Let $B(\ell)$ be the normalized matrix in \eqref{eq:residual-matrix}. Choose real numbers $\lambda_{r,s}$ and form
\[
 Y=\operatorname{diag}_{\Lambda}(e^{-\lambda_{r,s}/2})B(\ell).
\]
Denote by $Y_i$ the block of $J$ columns at node $i$, and write $G_i=Y_i^*Y_i$. If $C_i$ is the set of orders selected at that node, Fischer's inequality gives
\begin{equation}\label{eq:fischer}
 \log|\det B_C(\ell)|^2
 \leq\sum_\Lambda\lambda_{r,s}
       +\sum_{i=1}^n\log\det G_i[C_i,C_i].
\end{equation}
Indeed, $Y_C^*Y_C$ is a positive semidefinite Gram matrix, and its determinant is at most the product of the determinants of its diagonal blocks. The weights are free parameters, but their total cost $\sum\lambda_{r,s}$ must be restored. Some weights have negative logarithmic cost; this signed sum will be estimated separately.

\begin{lemma}[Radial monotonicity]\label{lem:radial}
For fixed row weights and a fixed set of derivative orders, each principal minor of a node Gram matrix is a polynomial with nonnegative coefficients in $q=|u_i|^2$.
\end{lemma}
\begin{proof}
Write $u_i=\sqrt q\,e^{\mathrm i\phi}$. The phase of an entry of $Y_i$ is $e^{\mathrm i(r-s-j)\phi}$, so it is removed by diagonal unitary transformations on the rows and columns. By Cauchy--Binet, the principal Gram minor is a sum of squared minors of $Y_i$. For a row set $E$ and order set $T$ of equal size, such a squared minor has the form
\[
 e^{-\sum_E\lambda_{r,s}}
 \prod_{j\in T}\frac{e^{2\ell j}}{R^{2j}}
 \bigl(\det[\fall rj]_{E,T}\bigr)^2
 q^{\sum_E(r+s)-\sum_Tj}.
\]
If the determinant is nonzero, at least one term in its expansion is nonzero. That term matches each $j$ to a row with $r\geq j$, so the exponent of $q$ is nonnegative. All coefficients are nonnegative, proving the claim; the value at $q=0$ follows by continuity.
\end{proof}

Thus a node can be moved outward to a radius bound before estimating its selected Gram minor. This assertion concerns principal minors, not monotonicity of individual eigenvalues. It is exactly the form needed to replace the ordered nodes by a finite set of radial bounds.

\subsection{Why negative-binomial weights help}
Fix parameters $a>0$, $0<v<1$ and $\beta<0$, and let $b=aR$. We use
\begin{equation}\label{eq:row-weights}
 e^{-\lambda_{r,s}}=\frac{(b)_r}{r!}v^r e^{-\beta s},
\end{equation}
where $(b)_r=b(b+1)\cdots(b+r-1)$. These are negative-binomial weights in the $r$ variable. They are useful because their moments against falling factorials can be computed explicitly.

The triangle allows a better estimate than completing the rows to a rectangle. Write $x_+=\max(0,x)$. At a node of squared radius $q>0$, put
\[
 \delta_q=(\log q-\beta)_+,\qquad v_q=v e^{-\delta_q/K},
 \qquad w=v_qq,\qquad \mathcal B=S e^{S\delta_q}.
\]
For a fixed $r<R$, the permitted $s$ satisfy $0\leq s<S-r/K$. Hence
\begin{equation}\label{eq:triangular-tilt}
 \sum_{0\leq s<S-r/K}(qe^{-\beta})^s
 \leq S e^{(S-r/K)\delta_q}.
\end{equation}
Multiplication by $e^{-r\delta_q/K}$ replaces $v$ by $v_q$. We may then extend the $r$ sum to all nonnegative integers: the added row outer products are positive semidefinite. After removing the node angle by diagonal unitary conjugation, $G_i$ is
bounded above in the positive semidefinite order by the $J\times J$ matrix
\begin{equation}\label{eq:moment-matrix}
 \mathcal G_{jk}
 =\mathcal B\left(\frac{e^\ell}{R\sqrt q}\right)^{j+k}
 \sum_{r\geq0}\frac{(b)_r}{r!}w^r\fall rj\fall rk,
 \qquad 0\leq j,k<J.
\end{equation}
The formulas at $q=0$ are understood by continuity. In the finite radial
partition used below all upper radii are positive. This completion changes the estimate, not the integral determinant being bounded. All its principal minors majorize the corresponding ones of $G_i$.

Here is the moment calculation behind the choice of weights. Normalize $(b)_rw^r/r!$ to have total mass one. The resulting random integer $r$ satisfies
\begin{equation}\label{eq:nb-generating}
 \mathbb E[(1+x)^r(1+y)^r]
 =\left(\frac{1-w}{1-w(1+x)(1+y)}\right)^b.
\end{equation}
Factoring the denominator as
\begingroup\postdisplaypenalty=10000
\[
 (1-w)\bigl((1-\alpha x)(1-\alpha y)-\delta_0xy\bigr),
 \qquad \alpha=\frac w{1-w},\quad
 \delta_0=\frac w{(1-w)^2},
\]
\endgroup
and expanding in $x$ and $y$ gives a triangular orthogonal factorization. Let $P_l$ be the monic polynomials orthogonal for this normalized
negative-binomial measure. These are the Meixner polynomials; their formulas are
\begin{align}
 \|P_l\|^2&=l!(b)_l\frac{w^l}{(1-w)^{2l}},\label{eq:meixner-norm}\\
 \fall rj&=\sum_{l=0}^j\binom jl(b+l)_{j-l}
          \left(\frac w{1-w}\right)^{j-l}P_l(r).
 \label{eq:meixner-expansion}
\end{align}
The generating identity proves the constants directly. No limiting asymptotic formula for orthogonal polynomials is involved.

Choose a column scale $\tau>0$ and set
\begin{equation}\label{eq:column-prices}
 \nu_j=\log(j!)+\log((b)_j)+2j(\ell-\log(R\tau)).
\end{equation}
After dividing column $j$ by $e^{\nu_j/2}$, the factorization
\eqref{eq:meixner-norm}--\eqref{eq:meixner-expansion} identifies the completed matrix factor with
$\sqrt{\mathcal B(1-w)^{-b}}$ times the matrix of composition by
\begin{equation}\label{eq:affine-map}
 \phi(z)=c_0+dz,\qquad
 c_0=\frac{\tau v_q\sqrt q}{1-w},\qquad
 d=\frac{\tau\sqrt{v_q}}{1-w},
\end{equation}
on polynomials of degree less than $J$. The Hilbert space has orthonormal monomials $\sqrt{(b)_j/j!}\,z^j$, or equivalently reproducing kernel $(1-z\bar w)^{-b}$. For example, the coefficient in row $l$, column $j$ of composition is
\[
 \binom jl\sqrt{\frac{(b)_j l!}{j!(b)_l}}c_0^{j-l}d^l,
\]
which is precisely the normalized factor obtained above. This converts a large derivative Gram matrix into a problem about an affine map of the disk.

\subsection{Two bounds for singular products}
We state the two estimates used in the spectral part of the computation. We prove the first here and give the second proof in Section~\ref{supp:complementary} of the supplementary notes.

\begin{lemma}[Affine composition]\label{lem:composition}
Suppose $c_0\geq0$, $d>0$ and $c_0+d<1$. Set
\[
 U=c_0d,\quad V=d^2,\quad Z=1-c_0^2,\qquad
 h=\frac{2U}{Z-V+\sqrt{(1-(c_0+d)^2)(1-(c_0-d)^2)}},
\]
\[
 \Gamma=(Z-Uh)^{-1},\qquad \xi^2=\frac{V+Uh}{Z-Uh}.
\]
The squared singular values of composition by $c_0+dz$ on the space with kernel $(1-z\bar w)^{-b}$ are $\Gamma^b\xi^{2l}$, $l\geq0$.
\end{lemma}
\begin{proof}
One has $0\leq h<1$, $0<\xi<1$, and
$Uh^2-(Z-V)h+U=0$. The kernel of the product of the composition operator with its adjoint is
\[
 [Z-U(z+\bar w)-Vz\bar w]^{-b}.
\]
The disk automorphism $(z-h)/(1-hz)$, accompanied by the multiplier
$(1-h^2)^{b/2}/(1-hz)^b$, induces a unitary operator, as is checked by substituting in the reproducing kernel. Conjugating the displayed kernel by this unitary and using the quadratic equation for $h$ gives $\Gamma^b$ times the kernel of the dilation operator $f(z)\mapsto f(\xi^2z)$. Its eigenvalues on the monomials are $\xi^{2l}$. Taking square roots gives the singular values of the original operator.
\end{proof}

Compression to polynomials of degree less than $J$ cannot increase singular values. If $m$ orders are selected at a node, let $N_m(\tau)$ be the sum of the largest $m$ numbers among $\nu_0,\ldots,\nu_{J-1}$. Restoring the column scales and the scalar mass gives the bound
\begin{equation}\label{eq:forward-cost}
 F_q^{\rm f}(m)=N_m(\tau)
   +m\log\mathcal B+bm\log\frac\Gamma{1-w}
   +m(m-1)\log\xi
\end{equation}
for its logarithmic Gram determinant. This applies to every selection of $m$ orders.
The term $m(m-1)\log\xi$ is negative and quadratic in $m$. It records the
geometric decay of the singular values, a saving lost when all $m$ singular
values are bounded by the largest one. This explains why selecting many
columns at a node can improve the estimate. The other terms, including the
capacity cost $N_m$, still have to be paid; the final calculation accounts
for all of them together.
Conveniently, one may choose $0<\xi<1$ first and then set
\[
 \Gamma=\frac{1+\xi^2+\sqrt{(1-\xi^2)^2+4w\xi^2}}2,\qquad
 \tau=\frac{\xi(1-w)}{\Gamma\sqrt{v_q}}.
\]
These choices automatically give a contracting map in \eqref{eq:affine-map}.

Near $m=J$, it can be better to estimate the complementary singular product. This allows an expanding finite affine map as well.
\begin{lemma}[Complementary singular product]\label{lem:complementary}
Let $T_J$ be composition by $c_0+dz$ on polynomials of degree less than $J$, where $c_0\geq0$ and $d>0$. Put $k=J-m$. For $0<r_0<1$, define
\[
 G_0=\frac{1+r_0^2+\sqrt{(1-r_0^2)^2+4c_0^2r_0^2}}2,\qquad
 s_0=\frac{dr_0}{G_0}.
\]
The logarithm of the product of the largest $m$ squared singular values of $T_J$ is at most
\begin{equation}\label{eq:complementary-cost}
 J(J-1)\log d+kb\log G_0+k(k-1)\log r_0-N\log s_0,
\end{equation}
where $N=k(2J-k-1)$ for $s_0\leq1$, and $N=k(k-1)$ for $s_0\geq1$.
\end{lemma}


Adding $N_m(\tau)+m\log\mathcal B-bm\log(1-w)$ to \eqref{eq:complementary-cost} bounds the same node minor as \eqref{eq:forward-cost}. We can take the minimum over either estimate and over any finite collection of its parameters. Each choice gives an upper bound for the original selected minor, so minimizing does not presume that a single parameter is optimal at every node.

\subsection{A bound summed over all order sets}
The preceding estimates bound a selected Gram minor through singular values.
We can also bound it by summing all principal minors of the same size: each
term is nonnegative, so the sum bounds any one of them, regardless of which
derivative orders occur in the nonsingular minor. The explicit moments above
give the following estimate for this sum. Its Schur-function proof is given
in the repository notes on spectral estimates. Let
\[
 M_{jk}(b,w)=\sum_{r\geq0}\frac{(b)_r}{r!}w^r\fall rj\fall rk,\qquad
 G_{jk}=\mathcal B\sigma_0^{j+k}M_{jk}(b,w),
\]
where $b,\mathcal B,\sigma_0>0$ and $0<w<1$. Put
$c_* =\sigma_0\sqrt w/(1-w)$ and $z_*=(bc_*)^{-2}$.
If $bc_*>1$, then, for $0\leq m\leq J$, with $a_*=J-m$ and
$s_*=ma_*+m(m-1)/2$,
\begin{equation}\label{eq:schur}
 \sum_{\substack{T\subseteq\{0,\ldots,J-1\}\\|T|=m}}
 \det G[T,T]
 \leq \mathcal B^m(1-w)^{-bm}c_*^{2s_*}
       \prod_{i=0}^{m-1}i!(b)_{2a_*+i}\,(1-z_*)^{-m^2}.
\end{equation}
The empty product is one. In our application $G=\mathcal G$ from
\eqref{eq:moment-matrix}, with $\sigma_0=e^\ell/(R\sqrt q)$. Taking the
logarithm of the right side of \eqref{eq:schur} therefore gives a third
upper bound for every selected node minor.

The reason Schur functions appear can be described without using them in the rest of the proof. The Meixner factorization and Cauchy--Binet first bound a selected minor by a gamma moment determinant. If its order set is $j_0<\cdots<j_{m-1}$, its deviation from consecutive orders is a partition. The gamma determinant is bounded by the consecutive-order value times the square of the corresponding Schur dimension. Summing these dimensions over all partitions uses the Cauchy identity and produces $(1-z_*)^{-m^2}$. It is the sum over all order sets, rather than an assertion that one consecutive set is largest, that makes \eqref{eq:schur} valid. The numerical estimate takes its minimum with the two singular-product bounds.

\subsection{Combining nodes with an occupancy constraint}
Let $F_q(m)$ be any common upper bound for the logarithmic determinant of an $m$-column minor at squared radius $q$, obtained as above. For a real parameter $a_0>0$, choose $\Phi(q)\geq0$ such that
\begin{equation}\label{eq:H-occupancy}
 F_q(m)\leq a_0Rm+R^2\Phi(q)\qquad(0\leq m\leq J).
\end{equation}
Any function $\Phi$ satisfying this inequality will suffice. Summing over nodes and using $\sum_i|C_i|=M$, \eqref{eq:fischer} gives
\begin{equation}\label{eq:occupancy-exact}
 \frac1W\log|\det B_C(\ell)|^2
 \leq U_{\rm row}+\frac{3a_0(S+1)}{2S}
                   +\frac3S\sum_{i=1}^n\Phi(q_i),
 \qquad U_{\rm row}=\frac1W\sum_\Lambda\lambda_{r,s}.
\end{equation}
This elementary use of a common slope incorporates the total number of selected columns. Bounding each node independently at its worst occupancy would discard that information.

The disk estimate supplies
\[
 q_i\leq Q(i/n),\qquad
 Q(t)=\left(\frac{\sqrt t+\eta}{1+\eta}\right)^2,\qquad
 \eta=\frac{\sqrt{D+4}/4}{\sqrt{n\sqrt D/(2\pi)}}.
\]
Partition $i/n$ into $N_t$ intervals. For the $k$th interval, first move all
its nodes outward to squared radius $Q(k/N_t)$ using Lemma~\ref{lem:radial},
and then apply \eqref{eq:H-occupancy} at that radius. This needs no monotonicity
assumption on the chosen function $\Phi$. Each interval contains at most
$n/N_t+1$ nodes. Substituting $n=R(S+1)/(2\theta J)$ gives
\begin{equation}\label{eq:metric-mesh}
 U_{\rm var}=U_{\rm row}+\frac{3a_0(S+1)}{2S}
   +3\left(\frac{1+1/S}{2\theta(J/R)N_t}+\frac1S\right)
       \sum_{k=1}^{N_t}\Phi(Q(k/N_t)),
 \qquad \theta=\frac{M}{nJ}.
\end{equation}
The nonnegativity of $\Phi$ justifies using upper bounds for the node counts. The signed row sum $U_{\rm row}$ remains outside that estimate.

For the finite comparisons below we take
\begin{equation}\label{eq:spectral-parameters}
 (a,v,\beta,a_0)=\left(\frac15,\frac7{10},-\frac45,\frac3{50}\right),
 \qquad (N_t,N_z,N_m)=(128,256,512).
\end{equation}
The latter two mesh sizes are used for column-order capacity and occupancy, respectively. Whole-cell upper bounds are taken before choosing the minimum of the three estimates; endpoints alone would not bound an entire cell. The supplementary notes, Section~\ref{app:metric}, give these finite bounds. Substitution into \eqref{eq:metric-mesh}, with directed rounding, supplies $U_{\rm var}$ in the determinant comparison. At no stage is the enormous matrix itself constructed.

\section{Nonvanishing with fewer nodes}\label{sec:rank}
The last algebraic input is full row rank of the matrix $V$. A dimension count
requires $nJ\geq M$, but it does not prove that the rows are independent.
The criterion below comes close to that count and is sufficient for all the
finite comparisons between $499$ and $875$. The proof uses a divided-Frobenius basis to retain information lost by direct reduction modulo a prime.

Direct reduction modulo an inert prime loses information: the polynomial
$z^p-w$ vanishes at every reduced node. The remedy is to divide out the
powers of $p$ whose presence we already know. After this division, a hypothetical
row relation gives a polynomial in two variables with many prescribed local
zeros. A differential operator converts those conditions into too many zeros
of a nonzero polynomial in one variable. The two main points are to prove that
this last polynomial is nonzero and to compare its degree with its orders of
vanishing.

\begin{theorem}[Sufficient rank criterion]\label{thm:simple-rank}
Let $p$ be an inert prime, and suppose that the nodes $z_1,\ldots,z_n$ have
distinct residues modulo $p\OO_D$. Let $R=KS$ and
\[
 K<p<R\leq p^2,\qquad 1\leq J\leq R.
\]
If $n\geq2S$ and
\begin{equation}\label{eq:simple-rank-ineq}
 n\left(J-\frac{RJ}{p^2}-\frac Rp\right)>
 \frac{(R+p)^2}{2K},
\end{equation}
then the matrix $V$ in \eqref{eq:jet-matrix} has full row rank.
\end{theorem}

The leading terms of \eqref{eq:simple-rank-ineq} are $nJ>R^2/(2K)$, close
to the necessary inequality $nJ\geq M=R(S+1)/2$. The errors are small when
$p/R$ and $R/p^2$ are small. For our application the choice of $p$ must also
allow distinct node residues and guarantee that at least one candidate prime
is inert. These requirements will be checked together in
Section~\ref{sec:completion}.

\subsection{Dividing out the Frobenius factors}
Put $g=z^p-w$ and $d=\lfloor(R-1)/p\rfloor<p$. The integral basis change
\eqref{eq:frobenius-basis} gives the rows
\begin{equation}\label{eq:g-basis}
 g^m z^a w^s,\qquad 0\leq m\leq d,\quad 0\leq a<p,\quad a+Ks<R-mp.
\end{equation}
At node $i$ write
\[
 g(z_i+t,\bar z_i)=pU_i(t)+t^p,
 \quad
 U_i(t)=\frac{z_i^p-\bar z_i}{p}
       +\sum_{v=1}^{p-1}\frac1p\binom pv z_i^{p-v}t^v\in\OO_D[t].
\]
Direct reduction would replace $(pU_i(t)+t^p)^m$ by $t^{mp}$, discarding
$U_i$ and making layer $m$ vanish in every column of order less than $mp$.
Instead, divide each row in layer $m$ by $p^m$ and multiply each column of order
$j=qp+v$, $0\leq v<p$, by $p^q$. Call the resulting matrix $\widetilde V$.
Its entries are integral. Indeed, in the binomial expansion of
$(pU_i+t^p)^m$, the term indexed by $q'$ can contribute to order $j$ only
when $q'\leq q$; after rescaling its coefficient contains $p^{q-q'}$.
Modulo $p$, only the term $q'=q$ survives.

Write $k=\OO_D/p\OO_D$, $\xi_i=z_i\bmod p$, and
$u_i(z)=(U_i\bmod p)(z-\xi_i)$. Thus, when $q\leq m$, the reduced entry is
\begin{equation}\label{eq:scaled-entry}
 \binom mq[t^v]\bigl(u_i(\xi_i+t)^{m-q}(\xi_i+t)^a\xi_i^{ps}\bigr).
\end{equation}
When $q>m$, all terms are divisible by $p$, and the reduced entry is zero.
The basis change and all scalings are invertible in characteristic zero.
It therefore suffices to prove that $\widetilde V\bmod p$ has full row rank
over $k$. These scalings are used only for nonvanishing; the determinant
comparison continues to concern the original matrix $V$.

Suppose, to the contrary, that the reduced rows satisfy a nontrivial relation
with coefficients $c_{m,a,s}\in k$. Package that relation as
\begingroup\postdisplaypenalty=10000
\begin{equation}\label{eq:rank-F}
 F(z,u)=\sum_{m=0}^d H_m(z)u^m,\qquad
 H_m(z)=\sum_{a,s}c_{m,a,s}z^{a+ps}.
\end{equation}
\endgroup
Within each layer the exponents $a+ps$ are distinct, since $0\leq a<p$.
Consequently $F\ne0$. Write $e=\deg_uF\leq d<p$. The triangular support gives
\begin{equation}\label{eq:rank-degree}
 \deg H_m<E_m:=\frac{p(R-mp)}K,
\end{equation}
because
$a+ps=(p/K)(a+Ks)-(p/K-1)a<p(R-mp)/K$.
The importance of this bound is its constant drop:
$E_m-E_{m+1}=p^2/K>p$.

To express the local conditions on $F$, put
\[
 L_l=\min\bigl(p,\max(0,J-lp)\bigr).
\]
There are exactly $L_l$ derivative orders in the block $j=lp+v$.
In characteristic $p$, $(\xi_i+t)^{ps}\equiv\xi_i^{ps}\pmod{t^p}$.
We may therefore replace $\xi_i^{ps}$ by $(\xi_i+t)^{ps}$ when taking any
of the coefficients $[t^v]$ with $v<L_l$. Summing \eqref{eq:scaled-entry}
against the relation coefficients gives
\begin{equation}\label{eq:local-hasse-vanish}
 \ord_{z=\xi_i}\bigl((\partial_u^{[l]}F)(z,u_i(z))\bigr)\geq L_l
 \qquad(i=1,\ldots,n;\ l\geq0).
\end{equation}
Here $\partial_u^{[l]}$ denotes the Hasse derivative in $u$. Thus a relation
among the scaled rows has become a family of vanishing conditions along
polynomial curves $u=u_i(z)$.

\subsection{One differential operator for all nodes}
Differentiating $pU_i(t)=(z_i+t)^p-\bar z_i-t^p$ before reducing modulo $p$
gives
\begin{equation}\label{eq:ui-prime}
 u_i'(z)=z^{p-1}-(z-\xi_i)^{p-1}.
\end{equation}
The first term is independent of the node, and the second vanishes to order
$p-1$ there. This suggests using the common operator
\[
 \mathcal D=\partial_z+z^{p-1}\partial_u
\]
in place of differentiation along each separate curve. Its key property is
\begin{equation}\label{eq:Dp-identity}
 \mathcal D^p=-\partial_u\qquad\hbox{on }k(z)[u].
\end{equation}
Indeed, the $p$th iterate of a derivation in characteristic $p$ is a derivation.
It sends $z$ to zero and $u$ to $(p-1)!=-1$, which determines it on $k(z)[u]$.

\begin{lemma}\label{lem:D-cyclic}
The polynomials $F,\mathcal DF,\ldots,\mathcal D^eF$ are linearly independent
over $k(z)$. In particular
\begin{equation}\label{eq:WF-def}
 W_F(z)=\det\bigl([u^m]\mathcal D^jF\bigr)_{0\leq j,m\leq e}
\end{equation}
is nonzero, and
\begin{equation}\label{eq:WF-degree}
 \deg W_F\leq\sum_{m=0}^e E_m-\frac{e(e+1)}2.
\end{equation}
\end{lemma}
\begin{proof}
If these polynomials were dependent, let $r\leq e$ be the first index of a
dependence. Their span up to index $r-1$ would be stable under $\mathcal D$:
applying $\mathcal D$ to a rational coefficient differentiates that coefficient,
and the image of the last spanning vector is in the same span by assumption.
By \eqref{eq:Dp-identity} this space is also stable under $\partial_u$.
It would then contain
$F,\partial_uF,\ldots,\partial_u^eF$. Since $e<p$, these have distinct
$u$-degrees $e,e-1,\ldots,0$ and nonzero leading coefficients. They are
independent, contradicting the dimension $r\leq e$.

For the degree bound, induction using
\[
 [u^m]\mathcal DG=([u^m]G)' +(m+1)z^{p-1}[u^{m+1}]G
\]
shows that $\deg([u^m]\mathcal D^jF)\leq E_m-j$. The first term lowers
degree by one; the second has degree at most
$E_{m+1}-j+p-1<E_m-j-1$. In every determinant term each column index $m$
and row index $j$ occurs once. Adding these degree bounds proves
\eqref{eq:WF-degree}.
\end{proof}

The use of $e<p$ here is essential: it prevents the successive $u$-derivatives
from vanishing prematurely in characteristic $p$. The identity
$\mathcal D^p=-\partial_u$ is what connects this elementary fact to
differentiation in $z$.

\subsection{Counting zeros of the coefficient determinant}
Fix a node $i$. Translation from the basis $1,u,\ldots,u^e$ to
$1,u-u_i(z),\ldots,(u-u_i(z))^e$ is unitriangular. It leaves $W_F$ unchanged
and replaces its entries by
\[
 B^{(i)}_{j,l}(z)=(\partial_u^{[l]}\mathcal D^jF)(z,u_i(z)).
\]
We claim that
\begin{equation}\label{eq:Bi-order-bound}
 \ord_{\xi_i}B^{(i)}_{j,l}\geq\max(0,L_l-j).
\end{equation}
To see this, let $\mathcal D_i=\partial_z+u_i'(z)\partial_u$, which differentiates
along the curve. Equation~\eqref{eq:ui-prime} gives
$\mathcal D-\mathcal D_i=(z-\xi_i)^{p-1}\partial_u$. An induction, in which
each differentiation lowers the order of the error by at most one, gives
\[
 \mathcal D^jG\equiv\mathcal D_i^jG
       \pmod{(z-\xi_i)^{p-j}k[z,u]}\qquad(0\leq j<p).
\]
The Hasse derivatives in $u$ commute with both operators: their coefficients
depend only on $z$, and ordinary and Hasse $u$-derivatives commute.
After specializing $u=u_i(z)$, the corresponding entry is therefore congruent
to the $j$th ordinary derivative of
$(\partial_u^{[l]}F)(z,u_i(z))$, modulo $(z-\xi_i)^{p-j}$.
Its order is at least $\max(0,L_l-j)$ by \eqref{eq:local-hasse-vanish}.
Since $L_l\leq p$, the error has at least that order as well. This proves the claim.

We can now finish the proof by factoring powers of $z-\xi_i$ out of rows
or columns. Put $Q=\lfloor J/p\rfloor$. If $e<Q$, then $L_l=p$ for all
$0\leq l\leq e$. Every entry in row $j$ of $B^{(i)}$ has order at least
$p-j$, so
\[
 \ord_{\xi_i}W_F\geq\sum_{j=0}^e(p-j)
 =(e+1)(p-e/2)>(e+1)p/2.
\]
The residues are distinct and $n\geq2S=2R/K$. Summing this bound over the nodes
therefore gives more than $(e+1)pR/K$ zeros counted with multiplicity,
whereas \eqref{eq:WF-degree} bounds the degree by at most that number.
This contradicts $W_F\ne0$.

If $e\geq Q$, every entry in column $l$ has order at least
$\max(0,L_l-e)$. Factoring column by column and using
$\sum_{l=0}^Q L_l=J$ gives
\begin{equation}\label{eq:rank-local-total}
 \ord_{\xi_i}W_F\geq\sum_{l=0}^Q(L_l-e)
 =J-(Q+1)e\geq J-\frac{RJ}{p^2}-\frac Rp.
\end{equation}
For the degree, write $R=dp+s_0$, where $1\leq s_0\leq p$. Then
\[
 2p\sum_{m=0}^d(R-mp)
 =(R-s_0+p)(R+s_0)
 =R^2+pR+s_0(p-s_0)<(R+p)^2.
\]
All $E_m$ are positive. Hence
\[
 \deg W_F\leq\sum_{m=0}^e E_m\leq\sum_{m=0}^d E_m
 <\frac{(R+p)^2}{2K}.
\]
Summing \eqref{eq:rank-local-total} over the nodes and applying
\eqref{eq:simple-rank-ineq} again makes the total order of vanishing exceed
the degree. Both cases are impossible. The reduced matrix has full row rank,
and so does $V$, proving Theorem~\ref{thm:simple-rank}.

\section{The numerical comparisons and the bound 498}\label{sec:completion}
We now assemble the proof. There are two parts to the numerical argument. A uniform estimate, supplemented by adjacent
interval comparisons, excludes every $D\geq875$. Below $875$ the elementary
reduction leaves sixteen possible values, each of which is excluded by a
separate determinant comparison. We describe the coverage first, then give
the parameters and the decisive calculation at $499$.

\subsection{The range \texorpdfstring{$D\geq875$}{D >= 875}}
For this range we use the resultant rank criterion of
Proposition~\ref{prop:resultant-rank} and the radial size estimate of
Proposition~\ref{prop:radial-size}. Appendix~\ref{app:global} applies these
estimates on an unbounded interval and then on finitely many adjoining
intervals below it. Set $D_0=5\,760\,000$, $X=(D-3)/4$, and
$X_0=(D_0-3)/4$. On the unbounded interval, it gives
\begin{equation}\label{eq:tail-gap}
 G_{\rm tail}(D)
 >0.00031743+0.127829\log(X/X_0)\qquad(D\geq D_0),
\end{equation}
where $G_{\rm tail}$ is the normalized difference between the arithmetic
lower bound and the radial upper bound. This difference is positive
throughout $D\geq D_0$, contradicting the determinant inequality.

The interval between $875$ and $D_0$ is covered as follows. The gap bounds in
Table~\ref{tab:coverage} are rounded downward from the smallest
lower endpoint in each family. Each comparison uses a closed interval, and
successive intervals share an endpoint.
\begin{table}[ht]
\centering
\begin{tabular}{lrr}
\toprule
Range & Number of intervals & Lower gap bound\\
\midrule
$[875,3\,968\,539]$ & 28 & $0.00397338$\\
$[3\,968\,539,4\,894\,699]$ & 37 & $0.00032629$\\
$[4\,894\,699,5\,760\,000]$ & 36 & $0.00040940$\\
\bottomrule
\end{tabular}
\caption{Interval comparisons covering $[875,D_0]$. The first
family uses norm-form estimates, the second exact finite prime sums, and the
third explicit prime-sum bounds of Rosser and Schoenfeld.}
\label{tab:coverage}
\end{table}

Thus \eqref{eq:tail-gap} and
Table~\ref{tab:coverage} exclude all $D\geq875$ under the class-number-one
hypothesis.

\subsection{Parameters and rank for the remaining comparisons}
By Lemma~\ref{lem:reduction}, the values still possible above $498$ are $499$
and the fifteen primes in Table~\ref{tab:candidates}. For these comparisons
we use $K=43$ and a prime cutoff $P$. The other parameters are
\[
 S=\left\lfloor\frac{6P}{5K}\right\rfloor,
 \qquad R=KS,\qquad J=\lfloor R/8\rfloor,
 \qquad M=R(S+1)/2.
\]
For $D=499$ we take $P=312\,000\,000$ and $n=35\,394\,400$.
Table~\ref{tab:parameters} gives all four parameter choices.

\begin{table}[ht]
\centering
\begin{tabular}{rrrrr}
\toprule
$P$ & $S$ & $R$ & $J$ & $n$\\
\midrule
$36\,000\,000$ & $1\,004\,651$ & $43\,199\,993$ & $5\,399\,999$ & $4\,822\,400$\\
$169\,000\,000$ & $4\,716\,279$ & $202\,799\,997$ & $25\,349\,999$ & $22\,638\,400$\\
$312\,000\,000$ & $8\,706\,976$ & $374\,399\,968$ & $46\,799\,996$ & $41\,793\,600$\\
$312\,000\,000$ & $8\,706\,976$ & $374\,399\,968$ & $46\,799\,996$ & $35\,394\,400$\\
\bottomrule
\end{tabular}
\caption{Parameters for all determinant comparisons below $875$. The last
line is used on $[499,503]$; the other lines at the individual discriminants
listed in Table~\ref{tab:candidates}.}
\label{tab:parameters}
\end{table}

We first prove nonvanishing. There are $814$ primes in
$(A,B]=(160\,000,170\,000]$. Lemma~\ref{lem:auxiliary-count} bounds the number of non-inert primes
by a sum of square roots. This sum is less than $765$, both on $[499,503]$
and at each value in Table~\ref{tab:candidates}. Hence at least one of the
$814$ primes in the interval is inert.

Lemma~\ref{lem:node-disk} gives $2\rho_n<28\,928$ for every comparison
in Table~\ref{tab:parameters}, and $2\rho_n<22\,492$ at $D=499$.
Both are below $A=160\,000$, so Corollary~\ref{cor:distinct-residues}
gives distinct residues modulo every prime in $(A,B]$. Moreover,
$K<A<B<R<A^2$ in every case. The remaining rank conditions are
$J\leq R$, $n\geq2S$, and
\begin{equation}\label{eq:uniform-rank-check}
 n\left(J-\frac{RJ}{A^2}-\frac RA\right)-\frac{(R+B)^2}{2K}>0.
\end{equation}
At the smaller node count the difference in \eqref{eq:uniform-rank-check}
exceeds $7.22\times10^{11}$.
The reason for narrowing the prime interval can now be seen quantitatively.
Writing $\theta=M/(nJ)$, inequality~\eqref{eq:uniform-rank-check} is equivalent to
\[
 \theta<\frac{(1+1/S)(1-R/A^2-R/(AJ))}{(1+B/R)^2}.
\]
At the parameters for $499$, the right side is approximately $0.984431$
when $A=160\,000$, just above the actual occupancy $0.983995$. With
$A=150\,000$ it is approximately $0.982414$, which is too small.
The narrower interval improves the sufficient rank estimate while still
containing more primes than the norm form can account for as non-inert.

At the other three parameter choices the respective differences exceed
$4.12\times10^{12}$, $9.02\times10^{13}$ and $2.95\times10^{14}$.
Theorem~\ref{thm:simple-rank} therefore proves full row rank in every case.

\subsection{The fifteen comparisons above 499}
We now evaluate Proposition~\ref{prop:comparison}, with the local arithmetic
estimates of Proposition~\ref{prop:general-arithmetic} and the spectral bound
\eqref{eq:metric-mesh}. The spectral parameters and mesh sizes are fixed as
in \eqref{eq:spectral-parameters}. Table~\ref{tab:candidates} gives the result;
each displayed lower bound is rounded downward. The calculation uses outward rounding.
\begin{table}[ht]
\centering
\begin{tabular}{rrr@{\qquad}rrr}
\toprule
$D$ & $P/10^6$ & Lower gap bound & $D$ & $P/10^6$ & Lower gap bound\\
\midrule
523&312&0.01913120 &683&36&0.24689369\\
547&312&0.06770128 &691&36&0.25866104\\
563&169&0.09816857 &739&36&0.33084083\\
571&169&0.10589325 &787&36&0.39774315\\
587&36&0.08275416 &811&36&0.42960657\\
619&36&0.13882635 &827&36&0.45047644\\
643&36&0.17968979 &859&36&0.48970182\\
659&36&0.20911703 &&&\\
\bottomrule
\end{tabular}
\caption{Positive lower bounds for $G=\arith-\metric$ at all fifteen primes
$D\equiv3\pmod8$ in $(499,875)$.}
\label{tab:candidates}
\end{table}
Each entry is positive and contradicts \eqref{eq:master}. Thus every value in
the table is excluded.

\subsection{The decisive comparison at 499}
At $D=499$ we retain $P=312\,000\,000$, $K=43$ and the same $R,S,J$ as in
the preceding comparison, but reduce the node count to $n=35\,394\,400$.
The arithmetic and size calculations are performed on the whole interval
$[499,503]$, using the same parameters throughout. Only $499$ needs to be excluded: $503\equiv7\pmod8$ is already
impossible by Lemma~\ref{lem:reduction}.

Table~\ref{tab:499} separates the three contributions to the comparison.
The common term is
\[
 U_{\rm common}=\frac{W_1+W_2}{W}\log\rho+
                     \frac{Mc+2M\log M}{W},
\]
so that $G=\arith-U_{\rm common}-U_{\rm var}$.
The displayed arithmetic numbers are lower bounds and the size numbers are
upper bounds, rounded in the safe direction. The first numerical column uses the larger node count, with all
other parameters fixed.
\begin{table}[ht]
\centering
\begin{tabular}{lrr}
\toprule
 & Larger node count & Smaller node count\\
\midrule
$n$ & $41\,793\,600$ & $35\,394\,400$\\
$\theta=M/(nJ)$, approximately & $0.83333113$ & $0.98399487$\\
Arithmetic lower bound $\arith$ & $9.99153787$ & $9.99493844$\\
Upper bound for $U_{\rm common}$ & $9.71143220$ & $9.62002864$\\
Upper bound for $U_{\rm var}$ & $0.40353803$ & $0.32371122$\\
\bottomrule
\end{tabular}
\caption{The comparison on $[499,503]$ with $M=1\,629\,945\,955\,088\,368$.
Only the occupancy row uses approximate decimals; the other decimal entries
are one-sided bounds.}
\label{tab:499}
\end{table}

The arithmetic decomposition \eqref{eq:arithmetic-parts} at the smaller
node count is, approximately,
\[
 \arith\approx15.49500-6.91595+1.37599+0.03990\approx9.99494.
\]
The last two contributions are the two- and three-factor corrections.
Together they restore about $1.416$ to the arithmetic lower bound, far more
than the final margin; even the three-factor term, about $0.040$, is
comparable to that margin.

The change in the gap is about $0.1746$. Of this, $0.0914$ comes from the
common term, $0.0798$ from the variable spectral term, and $0.0034$ from the
arithmetic lower bound. The common-term improvement is close to the radius
estimate $0.085$ in Section~\ref{subsec:central}. The variable term benefits
from the higher occupancy and the changed residual scale $\ell$. 

The complete directed calculation gives the following enclosure for the
computed gap bound at the smaller node count, rounded outward here:
\begin{equation}\label{eq:499-gap}
 G\in[0.05119859,\;0.05120060].
\end{equation}
In particular $G>0.05119859$, contradicting the determinant comparison.
For the larger node count the corresponding enclosure is
\[
 [-0.12343235,\;-0.12343034].
\]

\begin{proof}[Proof of Theorem~\ref{thm:main}]
Suppose that $h(-D)=1$ and $D>498$. Lemma~\ref{lem:reduction} makes $D$ a
prime congruent to $3$ modulo $8$. Equation~\eqref{eq:tail-gap} and
Table~\ref{tab:coverage} exclude $D\geq875$. Table~\ref{tab:candidates}
excludes all the remaining possibilities greater than $499$, and
\eqref{eq:499-gap} excludes $499$ itself. These cases exhaust all $D>498$.
\end{proof}

At $499$, the occupancy is already about $0.984$, so the number of nodes
is only about $1.6\%$ above the dimension-count minimum $M/J$. There is
therefore little room to reduce the radius further by improving rank alone:
in the radius approximation of Section~\ref{subsec:central}, reaching the
dimension count exactly would save only about $0.008$ in the common term.
Sharper arithmetic or spectral estimates are natural candidates for further
progress.

The class number one hypothesis is used to turn non-inert prime ideals into
representations by the norm form. A comparable estimate for general
Landau--Siegel zeros would require a substitute for this arithmetic input.

\appendix
\section{Local divisibility and finite row sums}\label{app:local}\label{supp:local}
The local estimates have two separate ingredients. A change of polynomial
basis makes divisibility visible in each row. A counting argument then bounds
how much of that divisibility can be lost to the derivative orders in a
selected minor. We give both arguments here, including the finite corrections
needed for the small primes.

Throughout this appendix, $K,S,J,n$ are positive integers, $R=KS$, and
\[
 \Lambda=\{(r,s)\in\NN^2:r+Ks<R\},\qquad
 M=\frac{R(S+1)}2,\qquad W=\frac{R^2S}3,\qquad
 \theta=\frac{M}{nJ}\leq1.
\]
The nodes $z_1,\ldots,z_n$ may be any elements of $\OO_D$. Columns are indexed
by $(i,j)$ with $1\leq i\leq n$ and $0\leq j<J$, and $C$ denotes any selection
of $M$ columns. Thus each derivative order occurs in $C$ at most $n$ times.
The determinant $\Delta_C$ is formed with the Hasse derivatives in
\eqref{eq:jet-matrix}; differentiation is in $z$, with $w$ held fixed.

\subsection{Passing from entries to a determinant}
The following elementary observation allows different derivative estimates
to be used on different sets of rows.

\begin{lemma}\label{lem:local-determinant}
Let $A=(a_{rc})$ be a square matrix over a discretely valued field, and let
$v$ be its valuation. Suppose that $E$ is a set of rows and
\[
 v(a_{rc})\geq a_r-b_c-\boldsymbol1_{r\in E}d_c,
 \qquad d_c\geq0.
\]
Then
\[
 v(\det A)\geq\sum_r a_r-\sum_c b_c
 -\max_{\substack{I\subseteq C\\|I|=|E|}}\sum_{c\in I}d_c.
\]
An integral unitriangular change of rows does not change the determinant.
If columns of order $j$ have cost $d(j)$, each order is available at most
$n$ times, and $d(j)$ is nondecreasing, the maximum cost of $m$ columns is
bounded by filling the orders $J-1,J-2,\ldots$ in that order, with capacity
$n$ at each order. The minimum of a nondecreasing cost is bounded from below
by filling orders $0,1,\ldots$ in the same way.
\end{lemma}
\begin{proof}
Each term of the determinant uses every row and every column once. The
columns paired with rows in $E$ form a set of size $|E|$, so the displayed
entry estimate gives the asserted lower bound for every determinant term.
The valuation of a sum is at least the minimum valuation of its terms.
A unitriangular row transformation has determinant one. Finally, if a
selected column has smaller cost than an available column, exchanging them
cannot decrease the total cost. Repeating this exchange proves the capacity
assertions, and reversing the exchanges proves the assertion about minima.
\end{proof}

\subsection{Two finite row sums}\label{supp:rows}
For integers $N\geq0$ and $m\geq1$, define
\[
 T(N,m)=\sum_{s=0}^{N-1}(N-s)\floor{s/m}.
\]
An empty sum is zero. Expressing $\floor{s/m}$ as the number of positive
multiples of $m$ not exceeding $s$, and interchanging the two finite sums,
gives
\begin{equation}\label{eq:T-floor}
 T(N,m)=\frac12\sum_{\substack{a\geq1\\am<N}}(N-am)(N-am+1).
\end{equation}
There are $K(S-s)$ rows with a given second coordinate $s$. Consequently
\[
 \sum_{(r,s)\in\Lambda}\sum_{a\geq1}\floor{s/q^{2a}}
 =K\sum_{a\geq1}T(S,q^{2a}).
\]
In particular, this sum is finite even when written with no upper limit.

The corresponding sum in the first coordinate is conveniently bounded by
\begin{equation}\label{eq:Sr}
 S_r(R,K,x)=\sum_{a=1}^{\floor{R/x}}(R-ax)(R+K-ax),\qquad x>0.
\end{equation}
For an integer $q\geq1$, put $Q=\floor{R/q}$ and $c_a=aq\bmod K$.
Then the exact identity is
\begin{equation}\label{eq:local-r-sum}
 2K\sum_{\Lambda}\floor{r/q}
 =S_r(R,K,q)+\sum_{a=1}^{Q}c_a(K-c_a)
 \geq S_r(R,K,q).
\end{equation}
To see this, count the rows with $r\geq aq$. Write $aq=Kb+c$, where
$0\leq c<K$, and put $u=S-b$. For each $s=0,\ldots,u-1$ there are
$K(u-s)-c$ such rows. Their number is $Ku(u+1)/2-cu$.
Multiplication by $2K$ gives
$(R-aq)(R+K-aq)+c(K-c)$. Summing over $a$ proves
\eqref{eq:local-r-sum}.

We will also need the shape of this weight as its argument varies. On an
interval on which $Q=\floor{R/x}$ is fixed, summing $a$ and $a^2$ gives
\begin{equation}\label{eq:local-Sr-polynomial}
 S_r(R,K,x)=QR(R+K)
 -\frac{Q(Q+1)(2R+K)}2x
 +\frac{Q(Q+1)(2Q+1)}6x^2.
\end{equation}
At an endpoint $x=R/Q$, the summand that disappears is zero, so $S_r$ is
continuous. If $u=R-Qx$, differentiation of this polynomial yields
\[
 -\frac{d}{dx}\bigl(xS_r(R,K,x)\bigr)
 =Q\left(\frac{Q-1}{2}x^2+u(2x-u)+K(x-u)\right).
\]
For $0<x<R$ we have $Q\geq1$ and $0\leq u<x$, and the right side is
positive. Since $\log x/x$ is positive and decreasing for $x\geq3$, it
follows that
\begin{equation}\label{eq:local-weight-decreasing}
 x\longmapsto\frac{S_r(R,K,x)\log x}{KW}
 \quad\hbox{is decreasing for }3\leq x<R.
\end{equation}
It is zero for $x\geq R$. This is the monotonicity used when replacing a
possible prime norm by a lower endpoint for that norm.

For the unbounded range, an explicit lower bound is more useful than a
piecewise polynomial. Set $\vartheta=\{R/q\}$ and substitute
$Q=R/q-\vartheta$ into \eqref{eq:local-Sr-polynomial}. Expanding gives
\begin{align*}
 S_r(R,K,q)={}&\frac{R^3}{3q}-\frac{R^2}2
 +\frac{KR^2}{2q}+\frac{qR}6-\frac{KR}2\\
 &-\frac{q^2}{6}\vartheta(1-\vartheta)(1-2\vartheta)
 +\frac{Kq}{2}\vartheta(1-\vartheta).
\end{align*}
The last term is nonnegative, and
$|\vartheta(1-\vartheta)(1-2\vartheta)|\leq1/(6\sqrt3)$.
Using $KW=R^3/3$, we obtain
\begin{equation}\label{eq:local-Bernoulli}
 \frac{S_r(R,K,q)}{KW}\geq
 \frac1q-\frac{3}{2R}+\frac{3K}{2Rq}
 +\frac{q}{2R^2}-\frac{3K}{2R^2}
 -\frac{q^2}{12\sqrt3 R^3}.
\end{equation}
The cubic expression in $\vartheta$ is the usual third Bernoulli-polynomial
remainder; the displayed calculation fixes both its sign and its constant.

\subsection{The basis for a small inert prime}
Fix an inert prime $q<K$, and write
\[
 h=h_q(r,s)=\min\left(\floor{r/q},
       \floor{(R-1-r-Ks)/(K-q)}\right),\qquad
 \tau_q(s)=\sum_{a\geq1}\floor{s/q^{2a}}.
\]
The first bound on $h$ permits $h$ copies of $z^q$ to be replaced by $w$.
The second ensures that the resulting monomials stay inside $\Lambda$.
For fixed $s$ put $t=S-s$. The inequalities $h\geq a$ say exactly
$qa\leq r<Kt-(K-q)a$, so there are $K(t-a)$ such rows when $a<t$,
and none otherwise. Summing these superlevel counts gives
\begin{equation}\label{eq:local-h-sum}
 2\sum_\Lambda h_q(r,s)=\frac{R(S^2-1)}3=W_2.
\end{equation}

We first describe the rows on which the two derivative estimates will be
used. In the part $0\leq r<q(S-s)$, one has $h=\floor{r/q}$, so
$0\leq r-qh<q$. In the remaining part, write
\[
 R-1-r-Ks=(K-q)h+e,\qquad 0\leq e<K-q.
\]
Here $h$ is the second member of the minimum. If
$d=S-s-1-h$ and $c=K-1-e$, then
\begin{equation}\label{eq:local-strip}
 r-qh=Kd+c,\qquad 0\leq d<S,\quad q\leq c<K.
\end{equation}
Conversely, for any such $d,c$, the values
\[
 0\leq s<S-d,\qquad h=S-d-1-s,\qquad r=qh+Kd+c
\]
give exactly these rows. Thus for each pair $(d,c)$ there is a strip of
length $t=S-d$, and there are $K-q$ strips of every length
$1,\ldots,S$. The number
\[
 m=\floor{(r-qh)/q^2}=\floor{(Kd+c)/q^2}
\]
is constant on a strip. This constancy is what makes the finite estimate
below possible.

Choose an integer $L$ with $0\leq L\leq S$, and let $E_L$ consist of the
strips of length at most $L$. In particular,
\begin{equation}\label{eq:local-cut-size}
 m_L:=|E_L|=\frac{(K-q)L(L+1)}2,\qquad \varpi_L:=\frac{m_L}{M}.
\end{equation}
On these rows we
will use all available factors $w^q-z$. On the other rows we keep the tower
in $w$, which has no derivative cost.

Let $T_a$ be the inert tower from \eqref{eq:tower-value}. Writing
$s=\sum_{a\geq0}d_aq^{2a}$, put
\[
 B_s(w)=\prod_{a\geq1}T_a(w)^{d_a},\qquad
 u(s)=\floor{s/q}\bmod q,\qquad r_0=r-qh-q^2m.
\]
Thus $B_s$ is monic of degree $q^2\floor{s/q^2}$ and
$v_q(B_s(\bar z_i))\geq\tau_q(s)$. Define
\begin{equation}\label{eq:local-mixed-basis}
 P_{r,s}(z,w)=z^{r_0}(z^{q^2}-z)^m(z^q-w)^h w^{s\bmod q}
 \begin{cases}
 (w^q-z)^{\floor{s/q}},&(r,s)\in E_L,\\
 (w^q-z)^{u(s)}B_s(w),&(r,s)\notin E_L.
 \end{cases}
\end{equation}
In both cases the leading choice from the factors is $z^rw^s$, with
coefficient one.

Here are the two checks needed to use these polynomials as row operations.
First, replacing $z^q$ by $w$ increases the weight $r+Ks$ by $K-q$, but
there are at most $h$ such replacements. Their total increase is at most
$R-1-r-Ks$. Every other replacement in
\eqref{eq:local-mixed-basis} decreases the weight. All monomials therefore
remain in $\Lambda$. Second, order monomials lexicographically by
$(qr+s,s)$. Replacing $z^{q^2}$ by $z$, or $z^q$ by $w$, strictly decreases
the first coordinate. Replacing $w^q$ by $z$ preserves the first coordinate
and decreases the second. A lower term of $B_s$ decreases the first
coordinate. Hence the change from $z^rw^s$ to $P_{r,s}$ is integral and
unitriangular. It preserves every selected minor.

\subsection{Valuations and derivative costs}
At a node, Frobenius and conjugation agree modulo an inert prime. Taylor
expansion in $z$ consequently has the form
\[
 (z_i+t)^q-\bar z_i=qU(t)+t^q,\qquad
 \bar z_i^q-(z_i+t)=qV-t,
\]
\[
 (z_i+t)^{q^2}-(z_i+t)=qH(t)-t+t^{q^2},
\]
with $U,H\in\OO_D[t]$ and $V\in\OO_D$. Let $a$ denote the combined number
of factors of the latter two types: $a=m+\floor{s/q}$ on $E_L$, and
$a=m+u(s)$ off $E_L$.

To lose a factor of $q$ from one of the $h$ factors of the first type, a
Taylor term must spend at least $q$ degrees in $t$. To lose one from any of
the other $a$ factors, it must spend at least one degree. If $a_1$ and $a_2$
are the respective numbers of such choices in a term of degree $j$, then
$qa_1+a_2\leq j$ and $a_2\leq a$. Its valuation is at least
$h+a-a_1-a_2$. These two inequalities imply, respectively,
\[
 h+a-a_1-a_2\geq h+a-j,
 \qquad
 h+a-a_1-a_2\geq h-\floor{(j-a)/q}.
\]
All factors not yet counted have integral coefficients. Thus the entry
valuation, before adding the tower contribution, is at least
\begin{equation}\label{eq:local-entry}
 \max\left(0,h+a-j,h-\floor{(j-a)/q}\right).
\end{equation}

Set $\delta_q(j)=j-\floor{j/q}$. For a row outside $E_L$, the last member
of \eqref{eq:local-entry}, together with the tower, gives
\begin{equation}\label{eq:local-entry-low}
 v_q\bigl(\partial_z^{[j]}P_{r,s}(z_i,\bar z_i)\bigr)
 \geq h+\tau_q(s)+\floor{(m+u(s))/q}-\floor{j/q}.
\end{equation}
For a row in $E_L$, the middle member gives instead
\begin{align}
 v_q\bigl(\partial_z^{[j]}P_{r,s}(z_i,\bar z_i)\bigr)
 \geq{}&h+\tau_q(s)
       +\bigl(m+\floor{s/q}-\tau_q(s)\bigr)\notag\\
       &-\floor{j/q}-\delta_q(j).
 \label{eq:local-entry-high}
\end{align}
Both estimates have the same baseline row contribution $h+\tau_q(s)$
and the same column cost $\floor{j/q}$. The extra row contribution is
nonnegative in both cases, since $\tau_q(s)\leq\floor{s/q}$.
Only rows in $E_L$ incur the additional cost $\delta_q(j)$.

\subsection{Summing the extra row contribution}
For the rows outside the cut, the following finite form of the floor-sum
identity controls an incomplete strip.

\begin{lemma}\label{lem:local-Hermite}
For integers $m,t\geq0$ and $q\geq2$,
\[
 \sum_{s=0}^{t-1}\floor{\bigl(m+(\floor{s/q}\bmod q)\bigr)/q}
 \geq\frac{tm}{q}-\frac{q^2}{4}.
\]
\end{lemma}
\begin{proof}
Write $m=aq+b$, with $0\leq b<q$. On each complete block of $q^2$
consecutive values of $s$, the summand is $a$ for $q(q-b)$ values and
$a+1$ for $qb$ values. Its average is exactly $m/q$.
It remains to examine a final block of length $kq+v$, with
$0\leq k,v<q$. If $k<q-b$, its deficit from this average is
$b(k+v/q)\leq b(q-b)$. If $k\geq q-b$, the deficit is
$(q-b)(q-k-v/q)\leq b(q-b)$. Both bounds are at most $q^2/4$.
The case $b=0$ has zero deficit.
\end{proof}

We now sum over the strips in \eqref{eq:local-strip}. Put
\[
 c_{K,q}=\frac{K+q-1}{2}-(q^2-1).
\]
For a fixed length $t$, summing
$\floor{(K(S-t)+c)/q^2}\geq(K(S-t)+c-(q^2-1))/q^2$
over $c=q,\ldots,K-1$ gives
\begin{equation}\label{eq:local-strip-mass}
 \sum_{c=q}^{K-1}\floor{(K(S-t)+c)/q^2}
 \geq\frac{K-q}{q^2}\bigl(K(S-t)+c_{K,q}\bigr).
\end{equation}
The right side is allowed to be negative; the exact extra contribution is
still nonnegative.

Define
\begin{align*}
 A_q(L)&=\frac{2(K-q)}{q^2}\sum_{t=1}^{L}
              t\bigl(K(S-t)+c_{K,q}\bigr),\\
 B_q(L)&=\frac{2(K-q)}{q^3}\sum_{t=L+1}^{S}
              t\bigl(K(S-t)+c_{K,q}\bigr),\\
 E_q(L)&=A_q(L)+B_q(L)-\frac{q^2}{2}(S-L)(K-q)\\
 &\hspace{12mm}+2(K-q)\left(T(L,q)-\sum_{a\geq1}T(L,q^{2a})\right).
\end{align*}
Twice the sum of the extra row contributions in
\eqref{eq:local-entry-low}--\eqref{eq:local-entry-high} is at least $E_q(L)$.
Indeed, on each strip of length $t\leq L$, the contribution of $m$ is $tm$;
\eqref{eq:local-strip-mass} sums these contributions to at least $A_q(L)/2$.
The remaining contribution on those strips is
\[
 (K-q)\sum_{t=1}^{L}\sum_{s=0}^{t-1}
       \left(\floor{s/q}-\tau_q(s)\right)
 =(K-q)\left(T(L,q)-\sum_{a\geq1}T(L,q^{2a})\right).
\]
Here a given $s$ occurs $L-s$ times, explaining the function $T$.

For a strip of length $t>L$, Lemma~\ref{lem:local-Hermite} bounds the extra
contribution from below by $tm/q-q^2/4$. Summing the first term with
\eqref{eq:local-strip-mass} gives $B_q(L)/2$; there are
$(S-L)(K-q)$ such strips. The rows in the remaining left part of the
triangle have nonnegative extra contribution and may be omitted. These
observations account for every term of $E_q(L)$. Since the exact total
extra contribution is itself nonnegative, we have the stronger lower bound
\begin{equation}\label{eq:local-extra-positive}
 2\sum_\Lambda(\text{extra row contribution})\geq\pos{E_q(L)}.
\end{equation}
The positive part is taken before subtracting the extra derivative cost.

\subsection{The extra column cost and the retained weight}
The cost $\delta_q(j)$ is nondecreasing. Its prefix sum has a simple exact
form:
\begin{equation}\label{eq:small-column-prefix}
 \Sigma_q(B):=\sum_{j=0}^{B-1}\delta_q(j)
 =\frac{q-1}{2q}B^2-\frac{r_B(q-r_B)}{2q},
 \qquad r_B=B\bmod q.
\end{equation}
This follows by writing $B=aq+r_B$ and summing the complete blocks of the
floor function. It also holds for $B=0$.

Write $m_L=nt+r$, where $0\leq r<n$, and put $A=J-t$. By
Lemma~\ref{lem:local-determinant}, the extra derivative cost is at most
\[
 n\bigl(\Sigma_q(J)-\Sigma_q(A)\bigr)+r\delta_q(A-1).
\]
The last term is omitted when $r=0$. In particular, $A-1$ is only used
when $A\geq1$. To estimate the remainder in this formula, rewrite it as
\[
 n\Sigma_q(J)-(n-r)\Sigma_q(A)-r\Sigma_q(A-1).
\]
For $r=0$ the term with $A-1$ is again omitted. The remainder in $\Sigma_q(J)$
is nonpositive, whereas each of the other two remainders is bounded above
in absolute value by $\floor{q^2/4}/(2q)$. The quadratic parts give
\[
 \frac{q-1}{q}\left(m_LJ-\frac{m_L^2}{2n}\right)
 -\frac{q-1}{2q}r\left(1-\frac rn\right).
\]
Dropping the final nonpositive term proves the upper bound
\begin{equation}\label{eq:local-extra-cost}
 \frac{q-1}{q}\left(m_LJ-\frac{m_L^2}{2n}\right)
 +\boldsymbol1_{m_L>0}\frac{n\floor{q^2/4}}{2q}.
\end{equation}
If $m_L=0$, the exact cost is zero, so no remainder is charged.

Normalize twice this cost by $W$, and recall $m_L=\varpi_LM$ and
$M/(nJ)=\theta$. Since $2M/W=3(S+1)/(RS)$, the resulting expression is
\[
 C_q(L)=3\left(1+\frac1S\right)\varpi_L
       \left(1-\frac{\varpi_L\theta}{2}\right)\frac{q-1}{q}\frac JR
       +\boldsymbol1_{L>0}\frac{n\floor{q^2/4}}{qW}.
\]
We have now obtained every term of the small-prime formula.

\begin{proposition}\label{prop:local-small-prime}
For an inert prime $q<K$, define
\begin{equation}\label{eq:wkeep}
 w_{\rm keep}(q)=
 \frac{W_2+2K\sum_{a\geq1}T(S,q^{2a})}{W}
 +\max_{0\leq L\leq S}
       \left(\frac{\pos{E_q(L)}}W-C_q(L)\right).
\end{equation}
Then every selected minor satisfies
\[
 v_q(\Delta_C)\geq\frac W2w_{\rm keep}(q)
                       -\sum_{(i,j)\in C}\floor{j/q}.
\]
The maximum is over integers $L$. Using any one admissible value of $L$
in place of the maximum also gives a valid lower bound.
\end{proposition}
\begin{proof}
For the basis associated with $L$, the sum of the baseline row contributions
is, by \eqref{eq:local-h-sum} and \eqref{eq:T-floor},
\[
 \sum_\Lambda(h+\tau_q(s))
 =\frac{W_2}{2}+K\sum_{a\geq1}T(S,q^{2a}).
\]
Equation~\eqref{eq:local-extra-positive} bounds the extra row contribution
from below, and \eqref{eq:local-extra-cost} bounds its extra column cost
from above. Apply Lemma~\ref{lem:local-determinant} to
\eqref{eq:local-entry-low} and \eqref{eq:local-entry-high}. Every choice
of $L$ gives a bound for the same determinant, because all the bases are
unitriangular. Taking the best of these bounds proves the proposition.
\end{proof}

In particular, the argument includes $q=2$ when $2$ is inert in the maximal
order. The factors $1+1/S$ and the last remainder in $C_q(L)$ are exact
finite-size corrections.

\subsection{Split and ramified primes}
For completeness we give the local contributions at the other primes.
They are needed when primes beyond the guaranteed inert range are retained.
At a split prime, a tower in $w$ gives a contribution with no derivative
cost. At a ramified prime, the inert estimates remain valid in the prime
ideal valuation, but this has half the norm contribution of the inert
valuation.

Suppose first that $q$ splits. Then
$\OO_D/q\OO_D\simeq\FF_q\times\FF_q$, so $b^q-b$ is divisible by $q$
for every $b\in\OO_D$. Define monic polynomials
\[
 G_1(w)=w^q-w,\qquad
 G_{a+1}(w)=G_a(w)^q-q^{q^a-1}G_a(w),\qquad
 \mu_a=\frac{q^a-1}{q-1}.
\]
They have degrees $q^a$, and $G_a(\bar z_i)\in q^{\mu_a}\OO_D$.
Indeed, if $G_a(\bar z_i)=q^{\mu_a}b$, the next value is
$q^{q\mu_a}(b^q-b)$, divisible by $q^{q\mu_a+1}=q^{\mu_{a+1}}$.
For the base-$q$ expansion $s=\sum_{a\geq0}d_aq^a$, the polynomial
\[
 \widetilde A_s(w)=w^{d_0}\prod_{a\geq1}G_a(w)^{d_a}
\]
is monic of degree $s$ and has values divisible by
$q^{\sum_{a\geq1}\floor{s/q^a}}$. The basis
$z^r\widetilde A_s(w)$ is unitriangular on $\Lambda$, and differentiating
in $z$ does not alter this divisibility. Hence
\begin{equation}\label{eq:local-split}
 \Delta_C\in q^{K\sum_{a\geq1}T(S,q^a)}\OO_D.
\end{equation}
Its contribution to $\log|\Delta_C|^2/W$ is at least
\[
 \bigl(e(q)+s(q)\bigr)\log q,\qquad
 e(q)=\frac{2K}{W}\sum_{a\geq1}T(S,q^{2a}),\quad
 s(q)=\frac{2K}{W}\sum_{a\geq1}T(S,q^{2a-1}).
\]
There is no first-power derivative debit in this estimate. When the
arithmetic lower bound is organized with a common debit
$\sum_C\floor{j/q}$ for all primes, a split prime therefore permits that
debit to be refunded. Its minimum over $C$ is given by the capacity
principle in Lemma~\ref{lem:local-determinant}.

Now suppose that $q\geq K$ and $(q)=\mathfrak q^2$ is ramified. In our
application the ramified prime is $q=D>K$. Normalize $v_{\mathfrak q}$
by $v_{\mathfrak q}(\mathfrak q)=1$. Conjugation acts trivially on
$\OO_D/\mathfrak q\simeq\FF_q$, so $z_i^q-\bar z_i$ and
$\bar z_i^q-\bar z_i$ belong to $\mathfrak q$.
A substitute for the inert tower is
\[
 \widehat B_s(w)=(w^q-w)^{q\floor{s/q^2}}.
\]
It is monic of degree $q^2\floor{s/q^2}$, and
\[
 v_{\mathfrak q}(\widehat B_s(\bar z_i))
 \geq q\floor{s/q^2}\geq\tau_q(s).
\]
For the last inequality, put $a=\floor{s/q^2}$; then
$\tau_q(s)=a+\floor{a/q^2}+\floor{a/q^4}+\cdots
\leq aq^2/(q^2-1)\leq aq$.

Use the unitriangular basis
\[
 z^{r\bmod q}(z^q-w)^{\floor{r/q}}
 w^{s\bmod q^2}\widehat B_s(w).
\]
As in \eqref{eq:tower-valuation}, its entries have valuation at least
$\floor{r/q}+\tau_q(s)-\floor{j/q}$. Therefore
\begin{equation}\label{eq:local-ramified}
 v_{\mathfrak q}(\Delta_C)\geq
 \frac{S_r(R,K,q)}{2K}+K\sum_{a\geq1}T(S,q^{2a})
 -\sum_{(i,j)\in C}\floor{j/q}.
\end{equation}
One may replace a negative right side by zero, since the entries are
integral. Finally, an inert prime ideal has norm $q^2$, whereas a ramified
prime ideal has norm $q$. The contribution of
$v_{\mathfrak q}(\Delta_C)$ to $\log|\Delta_C|^2$ is consequently
$v_{\mathfrak q}(\Delta_C)\log q$. This explains the factor of one half
in the normalized ramified contribution used in the arithmetic comparison.

\section{Norm representations and the arithmetic estimate}\label{app:arithmetic}
The local estimates of Appendix~\ref{app:local} treat each prime separately.
To sum them, we begin by assigning the inert-prime contribution to every
prime below a cutoff. We then subtract the excess assigned to split and
ramified primes. Under class number one, these primes are represented by
the norm form. Counting norm values gives an upper bound for the loss;
removing composite values improves it. This appendix makes that argument
precise and accounts for the derivative orders in the selected minor.

Throughout, $D>163$ is prime, $D\equiv3\pmod8$, and $h(-D)=1$.
Put $m_D=(D+1)/4$ and
\[
 N_D(x,y)=x^2+xy+m_Dy^2
        =(x+y/2)^2+Dy^2/4.
\]
An element $x+y\omega$ is primitive if $\gcd(x,y)=1$. This is equivalent
to not being divisible in $\OO_D$ by a rational integer greater than one.
Since the units are $\{\pm1\}$, every nonreal element has a representative
modulo sign and conjugation with $y>0$ and $2x+y\geq0$.

\subsection{Prime and composite norm values}\label{app:norm-products}
\begin{lemma}\label{lem:primitive-norms}
Every split prime has exactly one primitive representative with
$y>0$ and $2x+y>0$. The only primitive representative on the imaginary
axis is $(-1,2)$, of norm $D$. Every prime factor of a primitive norm
is at least $m_D$.

If a primitive norm has $r\geq1$ distinct split prime factors, then
it has $2^{r-1}$ representatives modulo sign and conjugation, whether
or not the ramified prime also divides it. The ramified prime can occur
only to the first power.
\end{lemma}
\begin{proof}
A prime ideal above a split prime $q$ is principal. A generator has norm
$q$, and its conjugate generates the other prime ideal. The four generators
are related by sign and conjugation, giving the stated representative.
They cannot lie on either axis. On the imaginary axis $2x+y=0$;
primitivity then gives $(x,y)=(-1,2)$.

An inert prime dividing a norm divides the element itself in $\OO_D$,
so cannot divide a primitive norm. Every remaining prime is the norm of
a nonreal element, and hence is at least $m_D$ by
Lemma~\ref{lem:small-inert}. If $q=\pi\bar\pi$ is split, a primitive
element cannot contain both $\pi$ and $\bar\pi$. For each distinct split
prime factor there is therefore a choice between its two conjugate prime
factors, with the exponent already fixed by the norm. Conjugation pairs
these $2^r$ choices. The ramified factor is associated to its conjugate,
and its square is a rational integer times a unit. It supplies no further
choice and cannot occur twice in a primitive element.
\end{proof}

The restriction on the prime factors gives a short way to make the
necessary finite lists. The bound below concerns the cutoff for retained
split-prime contributions; it need not hold at the larger cutoff used
for the row estimate.
\begin{lemma}\label{lem:complete-factor-list}
Let $Q<m_D^4$. Every composite primitive norm at most $Q$ is a product
of two or three prime factors, counted with multiplicity. If
$F\geq Q/m_D$ and $F<m_D^4$, all of its prime factors occur among the
primitive elements of norm at most $F$ after the reducible elements
have been removed.

To perform that removal, it suffices to form products of a primitive
element of norm at most $\sqrt F<m_D^2$ and a primitive element of norm
at most $F/m_D$.
\end{lemma}
\begin{proof}
Each prime factor is at least $m_D$, by Lemma~\ref{lem:primitive-norms}.
Thus four factors would have product at least $m_D^4$, and any factor
of a composite norm at most $Q$ is at most $Q/m_D$.

If a primitive element of norm at most $F$ is reducible, choose an
irreducible factor of smallest norm. Its norm is at most $\sqrt F$,
and the complementary factor has norm at most $F/m_D$.
Both factors are primitive: any rational divisor of either would divide
the original element. Conversely, a product of two nonunits is reducible.
The indicated products therefore remove exactly the reducible elements.
Under class number one the remaining nonunit elements have prime norm.
\end{proof}

Multiplication in these lists is ordinary multiplication in the order:
\begin{equation}\label{eq:norm-product}
 (x+y\omega)(x'+y'\omega)
 =(xx'-m_Dyy')+(xy'+x'y+yy')\omega.
\end{equation}
After multiplication one keeps the primitive products and identifies
representatives related by sign or conjugation. Lemma~\ref{lem:primitive-norms}
gives their multiplicities. This provides both a complete list of composite
values up to $Q$ and, by subtraction from all primitive norm values, an
exact list of the split primes up to $Q$.

\subsection{The loss at a non-inert prime}
Assume $2\leq K<D/4$ and choose $P\geq\max(S,J-1)$.
All primes below $K$ are inert. For $x\geq K$, define
\begin{equation}\label{eq:prime-loss-weights}
 \begin{split}
 w(x)&=\frac{S_r(R,K,x)\log x}{KW},\\
 e(x)&=\frac{2K}{W}\sum_{a\geq1}T(S,x^{2a}),\qquad
 s(x)=\frac{2K}{W}\sum_{a\geq1}T(S,x^{2a-1}).
 \end{split}
\end{equation}
The sums $T$ and $S_r$ are defined in Appendix~\ref{supp:rows}.
The provisional row contribution is
\begin{equation}\label{eq:all-prime-contribution}
 \arith_0=\sum_{q<K}w_{\rm keep}(q)\log q
       +\sum_{K\leq q\leq P}\bigl(w(q)+e(q)\log q\bigr),
\end{equation}
where all sums indexed by $q$ are over primes. Its common derivative
loss is $\sum_{q\leq j}\lfloor j/q\rfloor\log q$.
The cutoff includes every prime appearing in this loss and every
nonzero conjugate-variable term.

Choose a finite set $\mathcal S$ of normalized primitive representatives
other than $(-1,2)$, all of norm at most $J-1$, and let $\mathcal Q$ be
the set of prime norms represented in $\mathcal S$.
By Lemma~\ref{lem:primitive-norms}, these primes are split and each is
represented once. When $D$ varies over an interval, we keep the integer
coordinate pairs in $\mathcal S$ fixed and require the norm bound
throughout the interval; $\mathcal Q$ may then depend on $D$.
We retain the first-power derivative contributions of $\mathcal Q$
individually.
For other split primes, a bound independent of the selected columns is
more convenient. Write $M=nt+r$, with $0\leq r<n$, and put
\begin{equation}\label{eq:minimum-column-sum}
 C(x)=n\sum_{j=0}^{t-1}\lfloor j/x\rfloor+r\lfloor t/x\rfloor.
\end{equation}
There are $n$ columns of each order. Since $j\mapsto\lfloor j/x\rfloor$
is nondecreasing, every selection of $M$ columns satisfies
$\sum_C\lfloor j/x\rfloor\geq C(x)$.

At a split prime, the conjugate-variable estimate retains the even-power
contribution $e(q)\log q$ and adds $s(q)\log q$. There is also no
first-power derivative loss. Thus, with $\iota=1$ when $q\in\mathcal Q$
and $\iota=0$ otherwise, an upper bound for the loss from the provisional
row contribution is
\begin{equation}\label{eq:split-loss}
 L(q,\iota)=
 \left[w(q)-s(q)\log q-(1-\iota)\frac{2C(q)\log q}{W}\right]_+.
\end{equation}
The positive part permits us to discard any improvement beyond the
provisional contribution. When $\iota=1$, the derivative contribution
is instead removed from the common loss. The same contribution is
therefore never used twice.

If $D\leq P$, the ramified-prime estimate in Appendix~\ref{app:local}
gives the loss bound
\begin{equation}\label{eq:ramified-loss}
 L_{\rm ram}(D)=\frac12
 \left[w(D)+e(D)\log D-\frac{2C(D)\log D}{W}\right]_+.
\end{equation}
If $D>P$, this term is absent. In particular, a cutoff below $D/4$
requires no non-inert-prime correction at all.

\subsection{Summing the losses by the norm form}
Let $\mathcal N$ contain the normalized primitive representatives of
norm in $[m_D,P]$, excluding $(-1,2)$. For each representative $\alpha$,
put $\iota_\alpha=1$ if $\alpha\in\mathcal S$ and $\iota_\alpha=0$
otherwise. Define
\[
 L_{\rm norm}=\sum_{\alpha\in\mathcal N}L(N(\alpha),\iota_\alpha)
             +\boldsymbol1_{D\leq P}L_{\rm ram}(D).
\]
Every split prime at most $P$ occurs once in this sum. Composite values
add nonnegative terms, so $L_{\rm norm}$ is an upper bound for the
actual loss. We may subtract any lower bound for these extra terms.
Let $C_2$ and $C_3$ be such lower bounds, obtained from disjoint primitive
representatives whose norms have respectively two and three prime factors.
Their multiplicities are as in Lemma~\ref{lem:primitive-norms}.

\begin{proposition}[Arithmetic estimate]\label{prop:general-arithmetic}
Suppose $R=KS$, $1\leq J\leq R$, $M\leq nJ$,
$2\leq K<D/4$, and $P\geq\max(S,J-1)$. Set
\begin{equation}\label{eq:general-arithmetic-terms}
 \arith=\arith_0-L_{\rm norm}+C_2+C_3,\qquad
 b(j)=\sum_{q\leq j}\lfloor j/q\rfloor\log q
           -\sum_{q\in\mathcal Q}\lfloor j/q\rfloor\log q.
\end{equation}
Then every nonzero maximal minor of the integral interpolation matrix
satisfies
\[
 W\arith-2\sum_C b(j)\leq\log|\Delta_C|^2.
\]
The conclusion also holds if $L_{\rm norm}$ is replaced by an upper bound
and $C_2,C_3$ by lower bounds. Additional composite corrections may be omitted.
\end{proposition}
\begin{proof}
Apply the inert-prime estimates of Appendix~\ref{app:local} provisionally
to all primes up to $P$. For a split prime, replace its row contribution
and derivative loss by the actual split-prime estimate. Equations
\eqref{eq:minimum-column-sum} and \eqref{eq:split-loss} bound the cost
of that replacement uniformly in $C$. Equation~\eqref{eq:ramified-loss}
does the same for the ramified prime. The product formula for
$|\Delta_C|^2$ then gives the claimed inequality with the true prime losses.
The norm sum bounds those losses from above, and $C_2+C_3$ subtracts
only terms contributed by composite representatives. This proves the result.
\end{proof}

This proposition applies both to cutoffs below the guaranteed inert
range and to cutoffs beyond it. For the latter, the row sum may include
norms greater than $m_D^4$. Only the complete removal of composite values
in the set selecting $\mathcal Q$ requires the smaller cutoff in
Lemma~\ref{lem:complete-factor-list}.

\subsection{Bounds over an interval of discriminants}\label{app:norm-intervals}
The comparison on $[499,503]$ and the first family of intervals above
$875$ use one estimate for several values of $D$. The norm form makes
this possible because, for a fixed pair $(x,y)$,
\begin{equation}\label{eq:norm-interval-affine}
 N_D(x,y)=N_{D_-}(x,y)+\frac{D-D_-}{4}y^2.
\end{equation}
We may therefore enclose each norm by its endpoint values. The following
observation identifies many factors without a separate test at every
intermediate discriminant.
\begin{lemma}\label{lem:interval-irreducibility}
Put $m_-=(D_-+1)/4$. If a primitive pair with $y>0$ satisfies
$N_{D_-}(x,y)<m_-^2$, then $N_D(x,y)<m_D^2$ for every $D\geq D_-$.
Consequently it is irreducible whenever $\OO_D$ has class number one.
\end{lemma}
\begin{proof}
The norm bound gives $y^2<4m_-^2/(4m_--1)<m_-+1$, so $y^2\leq m_-$.
Writing $m_D=m_-+u$, equation~\eqref{eq:norm-interval-affine} yields
$N_D<m_-^2+um_-\leq(m_-+u)^2$. A reducible primitive element has
at least two prime factors, each of norm at least $m_D$.
\end{proof}

All possible primitive norm representatives are found by the finite
conditions $N_{D_-}(x,y)\leq P$, $y>0$ and $2x+y\geq0$.
Let $[a,b]$ enclose one such norm, restricting it to $[m_-,P]$ when
bounding a possible prime. By \eqref{eq:local-weight-decreasing}, $w$
is decreasing on $[3,R)$ and zero beyond; $s$ and $C$ are nonincreasing
by their floor-sum definitions. These functions are nonnegative, so
valid endpoint bounds are
\begin{align}
 L^+(a,b,\iota)&=
 \left[w(a)-s(b)\log a-(1-\iota)\frac{2C(b)\log a}{W}\right]_+,
 \label{eq:loss-upper}\\
 L^-(a,b,\iota)&=
 \left[w(b)-s(a)\log b-(1-\iota)\frac{2C(a)\log b}{W}\right]_+.
 \label{eq:loss-lower}
\end{align}
At the ramified representative the analogous upper bound is
\[
 \frac12\left[w(a)+e(a)\log b-\frac{2C(b)\log a}{W}\right]_+.
\]
Summing \eqref{eq:loss-upper} over all possible representatives bounds
the loss from above. Products known to occur throughout the interval,
with their multiplicities, can be subtracted using \eqref{eq:loss-lower}.

Products of norm intervals are handled as follows. Suppose the factors
have norm intervals $[a_i,b_i]$, and put
$A=\prod_i a_i$, $B=\prod_i b_i$. At each discriminant their product
has norm in $[A,B]$, even if the order of the candidate norm values changes
as $D$ varies. If $(X,Y)$ is its normalized representative at that
discriminant, then $Y^2\leq\lfloor4B/D_-\rfloor$. Writing
$\delta m=(D_+-D_-)/4$, we obtain
\begin{equation}\label{eq:product-window-bound}
 N_{D_+}(X,Y)\leq B+\delta m\lfloor4B/D_-\rfloor.
\end{equation}
This proves membership in an upper norm window whenever the right side
is below its cutoff; any lower endpoint condition is checked in the same
way using \eqref{eq:norm-interval-affine}. A product is subtracted only
when this inclusion holds at every discriminant in the interval.

The factor lists are complete by the argument of
Lemma~\ref{lem:complete-factor-list}, applied at each admissible integer
$D$ in the interval. To remove a reducible candidate of norm at most $F$,
one factor has norm at most $\sqrt F$ and the other at most $F/m_D$.
Their lower endpoint norms are no larger, so the finite endpoint lists
contain them. Equation~\eqref{eq:norm-product} detects their product
at that $D$. The initial factors below $m_-^2$ are covered by
Lemma~\ref{lem:interval-irreducibility}. On the broad intervals used above
$875$, $P<m_-^3$, so only two-factor products are needed. On $[499,503]$
the retained split-prime cutoff is below $m_-^4$, so two- and three-factor
products suffice. Corrections to the row loss use only products proved
to stay below $P$ throughout the interval.

\subsection{The derivative remainder}\label{app:column-remainder}
Legendre's formula applied to \eqref{eq:general-arithmetic-terms} gives
\begin{equation}\label{eq:factorial-split-remainder}
 F(j)=\log(j!)-b(j)
 =\sum_q\sum_{a\geq2}\lfloor j/q^a\rfloor\log q
       +\sum_{q\in\mathcal Q}\lfloor j/q\rfloor\log q.
\end{equation}
Both sums are nondecreasing in $j$. Choose a nonnegative lower bound
$\sigma$ for
\[
 \sigma_{\rm pp}+\sum_{q\in\mathcal Q}\frac{\log q}{q},\qquad
 \sigma_{\rm pp}=\sum_{\substack{q,\ a\geq2\\q^a\leq J-1}}
                                      \frac{\log q}{q^a}.
\]
The remaining task is to bound $H(j)=j\sigma-F(j)$ from above for every
order $0\leq j<J$.

Recall the set $\mathcal S$ defining $\mathcal Q$. Let
$q_{\rm cap}\leq J-1$ bound the upper endpoints of its norm intervals;
for a fixed discriminant these intervals may be singletons.
Let $\mathcal C_{\rm out}$ list every possible composite representative
in $\mathcal S$, and let $\mathcal C_{\rm in}$ list representatives known
to be composite and in $\mathcal S$ throughout the interval.
Each product has multiplicity $\mu$ and a norm interval $[A,B]$.
The preceding factor and inclusion arguments supply these two lists.
Possible extra products in $\mathcal C_{\rm out}$ only weaken the lower
bounds below. Then
\begin{align}
 \sigma_-&=\sigma_{\rm pp}+\sum_{[a,b]\in\mathcal S}\frac{\log a}{b}
                 -\sum_{\mathcal C_{\rm out}}\mu\frac{\log B}{A},
 \label{eq:slope-lower}\\
 \sigma_+&=\sigma_{\rm pp}+\sum_{[a,b]\in\mathcal S}\frac{\log b}{a}
                 -\sum_{\mathcal C_{\rm in}}\mu\frac{\log A}{B},
 \label{eq:slope-upper}\\
 T_-(j)&=\sum_{[a,b]\in\mathcal S}\lfloor j/b\rfloor\log a
                 -\sum_{\mathcal C_{\rm out}}\mu\lfloor j/A\rfloor\log B
 \label{eq:split-floor-lower}
\end{align}
bound the slope from below and above, and the split-prime part of
\eqref{eq:factorial-split-remainder} from below. These inequalities follow
by subtracting the composite contributions from the sum over all
representatives. The containing list is used for a lower bound, and the
contained list for an upper bound.

We use $\sigma=\sigma_-\geq0$. On a block of integer orders
$u\leq j<v$, monotonicity of the actual function $F$ gives
\begin{equation}\label{eq:remainder-block}
 H(j)\leq(v-1)\sigma_-
       -\sum_q\sum_{a\geq2}\lfloor u/q^a\rfloor\log q-T_-(u).
\end{equation}
Assign this upper bound to each order in the block, and sort the resulting
$J$ values as $h_{(1)}\geq\cdots\geq h_{(J)}$. With $M=nt+r$ as before,
\begin{equation}\label{eq:finite-column-constant}
 c=\frac2M\left(n\sum_{i=1}^t h_{(i)}+rh_{(t+1)}\right)
 \quad\Longrightarrow\quad 2\sum_C H(j)\leq Mc.
\end{equation}
When $r=0$, the last summand is omitted. This bound fills the largest
values first, with capacity $n$ per order. Blocks of at most $32$ orders
are used in the finite comparisons.

For the unbounded range it is useful to retain only prime powers
$q^a\leq Q$, where $Q$ is fixed. Omitting the others discards a favorable
factorial contribution. For $d=q^a\leq Q$, put
\[
 b_d=n\lceil J/d\rceil,\quad
 t_d=\min(d,\lfloor M/b_d\rfloor),\quad r_d=M-b_dt_d.
\]
Each residue modulo $d$ occurs at most $b_d$ times among the columns.
Selecting residues in decreasing order therefore gives
\begin{equation}\label{eq:power-capacity}
 \frac2M\sum_C\sum_{\substack{q,a\geq2\\q^a\leq Q}}
       \{j/q^a\}\log q
 \leq2\sum_{d=q^a\leq Q}
 \frac{b_d\bigl(t_d(d-1)-\binom{t_d}{2}\bigr)
              +r_d\max(0,d-1-t_d)}{Md}\log q.
\end{equation}
The estimate treats each prime power separately, so does not require
simultaneously realizing the largest residues for different moduli.

\subsection{An inert prime in an auxiliary interval}\label{app:auxiliary-prime}
The following count supplies the prime used in Section~\ref{sec:completion}.
It uses the same norm representation, but requires only an upper bound
for the number of split primes, so composite values need not be removed.
\begin{lemma}\label{lem:auxiliary-count}
Let $0<A<B$, $16<D_-\leq D\leq D_+$, and suppose $D\notin(A,B]$.
Under class number one, the number of non-inert primes in $(A,B]$ is at most
\begin{equation}\label{eq:auxiliary-count}
 \sum_{y=1}^{\lfloor\sqrt{4B/D_-}\rfloor}
 \left(1+\sqrt{B-D_-y^2/4}
       -\sqrt{\max(0,A-D_+y^2/4)}\right)_+.
\end{equation}
\end{lemma}
\begin{proof}
Such a prime is split, so has a representative
$q=X^2+Dy^2/4$ with $X\geq0$ and $y\geq1$.
For fixed $y$, the possible $X=x+y/2$ lie in a translate of $\ZZ$,
and in the interval with endpoints
$\sqrt{\max(0,A-D_+y^2/4)}$ and $\sqrt{B-D_-y^2/4}$.
There are at most one plus its length. Also $D_-y^2/4\leq B$.
Summing over $y$ proves the bound.
\end{proof}

\section{Additional spectral estimates}\label{app:spectral}\label{app:metric}
The estimates in Section~\ref{sec:metric} reduce the size of a selected
minor to three bounds for a node Gram matrix. We first prove the
complementary singular-product bound and the Schur bound. We then give
finite estimates for the row weights and the three node bounds. All
parameters remain free, so the same estimates apply when the polynomial
space or the number of evaluation points changes.

\subsection{The complementary singular product}\label{supp:complementary}
The forward estimate uses an affine map taking the disk into itself.
For a finite polynomial space, an expanding affine map is also available.
To estimate its largest singular values, we pass to the inverse map and
estimate the complementary product.

\begin{proof}[Proof of Lemma~\ref{lem:complementary}]
The matrix $T_J$ is triangular, with diagonal entries $1,d,\ldots,d^{J-1}$.
Consequently
\[
 \det(T_J^*T_J)=d^{J(J-1)}.
\]
The product of its largest $m$ squared singular values is this determinant
times the product of the largest $k=J-m$ squared singular values of
$T_J^{-1}$. The inverse is composition by $(z-c_0)/d$.

Put $D_{s_0}=\operatorname{diag}(s_0^j)_{0\leq j<J}$, with $G_0,s_0$
as in the lemma. Then $T_J^{-1}D_{s_0}$ is composition by
\[
 \psi(z)=-\frac{c_0r_0}{G_0}+\frac{r_0}{G_0}z.
\]
This map is a strict disk contraction. Indeed, $G_0$ is the larger root
of
\[
 X^2-(1+r_0^2)X+(1-c_0^2)r_0^2=0.
\]
The left side at $X=(1+c_0)r_0$ equals
$-(1+c_0)r_0(1-r_0)^2<0$, so $G_0>(1+c_0)r_0$.
Changing the sign of the translation does not change singular values,
by the unitary reflection $f(z)\mapsto f(-z)$.
Lemma~\ref{lem:composition}, applied to $\psi$, gives squared singular
values $G_0^b r_0^{2l}$, $l\geq0$, for the infinite composition operator.
Compression to polynomials of degree less than $J$ cannot increase them.

For matrices $X,Y$, the norm inequality
$\|\bigwedge^k(XY)\|\leq
\|\bigwedge^kX\|\,\|\bigwedge^kY\|$ bounds the product of the largest
$k$ singular values of a product. Apply it to
\[
 T_J^{-1}=(T_J^{-1}D_{s_0})D_{s_0}^{-1}.
\]
The first factor contributes at most $G_0^{bk}r_0^{k(k-1)}$ to the
squared singular product. If $s_0\leq1$, the largest $k$ diagonal
entries of $D_{s_0}^{-1}$ have indices $J-k,\ldots,J-1$; if $s_0\geq1$,
they have indices $0,\ldots,k-1$. Their squared products are respectively
$s_0^{-k(2J-k-1)}$ and $s_0^{-k(k-1)}$. Multiplying these bounds and
taking logarithms gives \eqref{eq:complementary-cost}.
\end{proof}

\subsection{The Schur estimate}\label{app:schur-proof}
The Schur bound sums all principal minors of a given size. Its proof has
two parts. The Meixner factorization first compares each minor with a
gamma moment determinant. Schur polynomials describe how that determinant
depends on the selected orders, and the Cauchy identity sums this
dependence over all selections.

We use the following standard notation. A partition of length at most
$m$ is an integer sequence $\lambda_1\geq\cdots\geq\lambda_m\geq0$;
write $|\lambda|=\sum_i\lambda_i$. Its Schur polynomial is
\[
 s_\lambda(x_1,\ldots,x_m)
 =\frac{\det[x_j^{\lambda_i+m-i}]_{i,j=1}^m}
        {\det[x_j^{m-i}]_{i,j=1}^m}.
\]
The quotient extends to a symmetric polynomial when some variables
coincide. At $1^m=(1,\ldots,1)$ it has value
\begin{equation}\label{eq:schur-dimension}
 s_\lambda(1^m)
 =\prod_{1\leq i<j\leq m}
       \frac{\lambda_i-\lambda_j+j-i}{j-i}.
\end{equation}
We will use the Littlewood--Richardson product rule and the Cauchy
identity in their usual Schur specializations; see~\cite{fulton} and
\cite[Section~2]{forresterwarnaar}.

\begin{lemma}[Gamma moments]\label{lem:gamma-schur}
Let $b>0$ and $m\geq1$, let $0\leq j_0<\cdots<j_{m-1}$, and put
$D_T=\det[(b)_{j_i+j_k}]_{i,k=0}^{m-1}$. Then
\begin{equation}\label{eq:gamma-minor-bound}
 D_T\leq
 \left(\prod_{0\leq i<k<m}\frac{j_k-j_i}{k-i}\right)^2
 \prod_{i=0}^{m-1}i!(b)_{2j_i-i}.
\end{equation}
\end{lemma}
\begin{proof}
Let $\mathrm d\mu_b(x)=x^{b-1}e^{-x}\mathrm dx/\Gamma(b)$ on
$(0,\infty)$. Its $r$th moment is $(b)_r$. Expanding both determinants
and integrating term by term gives Andreief's identity in the form
\[
 D_T=\frac1{m!}\int_{(0,\infty)^m}
          \bigl(\det[x_k^{j_i}]_{i,k=0}^{m-1}\bigr)^2
          \prod_{k=0}^{m-1}\mathrm d\mu_b(x_k).
\]
Set $\lambda_{m-i}=j_i-i$. The alternant formula shows that
$\det[x_k^{j_i}]$ is the Vandermonde determinant times $s_\lambda(x)$,
up to a sign. For the consecutive orders $0,\ldots,m-1$, the moment
determinant is
\[
 D_0=\det[(b)_{i+k}]_{i,k=0}^{m-1}
     =\prod_{i=0}^{m-1}i!(b)_i.
\]
To see the last equality, factor $(b)_i$ from row $i$. The remaining
entry $(b+i)_k$ is monic of degree $k$ in $b+i$, so its determinant is
the Vandermonde product $\prod_{i=0}^{m-1}i!$.

Write $\Delta(x)=\prod_{i<k}(x_k-x_i)$. Normalize the density
$\Delta(x)^2\prod_k\mathrm d\mu_b(x_k)$ to have mass one, and denote
expectation for it by $\mathbb E$. The preceding calculation gives
$D_T=D_0\mathbb E[s_\lambda(X)^2]$. For any partition $\nu$ of length
at most $m$, a second use of Andreief gives
\begin{equation}\label{eq:gamma-schur-average}
 \mathbb E[s_\nu(X)]
 =s_\nu(1^m)\prod_{i=1}^m(b+m-i)_{\nu_i}.
\end{equation}
To evaluate it, Andreief writes the expectation as
\[
 \mathbb E[s_\nu(X)]
 =\frac{\det[(b)_{\nu_i+m-i+m-j}]_{i,j=1}^m}{D_0}.
\]
Factor $(b)_{\nu_i+m-i}$ from row $i$. The remaining entries are the
monic polynomials $(b+\nu_i+m-i)_{m-j}$. Dividing their Vandermonde
determinant by the one for $\nu=0$ gives
$s_\nu(1^m)$ by \eqref{eq:schur-dimension}; the ratios of the row
factors give the rising factorials in \eqref{eq:gamma-schur-average}.

The Littlewood--Richardson rule writes
$s_\lambda^2=\sum_\nu c_{\lambda\lambda}^{\nu}s_\nu$, with
nonnegative coefficients. Every occurring $\nu$ has size $2|\lambda|$
and is dominated by $2\lambda$: for every $r$,
$\sum_{i\leq r}\nu_i\leq2\sum_{i\leq r}\lambda_i$.
Terms with more than $m$ parts vanish in $m$ variables.
The decreasing vector $(b+m-i+2\lambda_i)_{i=1}^m$ therefore
majorizes $(b+m-i+\nu_i)_{i=1}^m$. Convexity of $\log\Gamma$ yields
\[
 \prod_{i=1}^m(b+m-i)_{\nu_i}
 \leq\prod_{i=1}^m(b+m-i)_{2\lambda_i}.
\]
Apply \eqref{eq:gamma-schur-average} to the product expansion and
evaluate that expansion at $1^m$. We obtain
\[
 D_T\leq D_0s_\lambda(1^m)^2
                    \prod_{i=1}^m(b+m-i)_{2\lambda_i}.
\]
Substituting $\lambda_{m-i}=j_i-i$ and using
\eqref{eq:schur-dimension} gives \eqref{eq:gamma-minor-bound}.
\end{proof}

\begin{proof}[Proof of \eqref{eq:schur}]
The case $m=0$ is immediate, so assume $m\geq1$. Use the notation
$M_{jk}(b,w)$, $\mathcal B$, $\sigma_0$, $c_*$ and $z_*$ from
Section~\ref{sec:metric}. The Meixner identities
\eqref{eq:meixner-norm}--\eqref{eq:meixner-expansion} give
\begin{equation}\label{eq:positive-meixner}
 M_{jk}(b,w)
 =(1-w)^{-b-j-k}(b)_j(b)_k
 \sum_{l=0}^{\min(j,k)}
 \frac{\fall jl\fall kl}{l!(b)_l}w^{j+k-l}.
\end{equation}
These are finite consequences of the generating identity
\eqref{eq:nb-generating}; the standard Meixner formulas are also recorded
in~\cite[Sections~18.19 and~18.23]{dlmf}.

Fix $T=\{j_0<\cdots<j_{m-1}\}$, write $s=\sum_{j\in T}j$, and put
$V_{lj}=\fall jl/\sqrt{l!(b)_l}$ for $0\leq l,j<J$.
Cauchy--Binet applied to \eqref{eq:positive-meixner} gives
\begin{equation}\label{eq:meixner-minor-positive}
 (1-w)^{bm+2s}\det M[T,T]
 =\prod_{j\in T}(b)_j^2
   \sum_{\substack{L\subseteq\{0,\ldots,J-1\}\\|L|=m}}
      \det(V_{L,T})^2w^{2s-\sum_{l\in L}l}.
\end{equation}
A nonzero term matches each $l\in L$ to an order $j\geq l$.
Thus $\sum_{l\in L}l\leq s$. After taking out $w^s$, the remaining
polynomial has nonnegative coefficients and is at most its value at
$w=1$. The finite Vandermonde identity is
\[
 (b)_j(b)_k\sum_l\frac{\fall jl\fall kl}{l!(b)_l}=(b)_{j+k}.
\]
For example, it follows by comparing the coefficient of $x^j$ in
$(1-x)^{-b-k}=(1-x)^{-b}(1+x/(1-x))^k$.
Consequently \eqref{eq:meixner-minor-positive}, with the column and
mass factors restored, gives
\begin{equation}\label{eq:schur-gamma-reduction}
 \det G[T,T]\leq\mathcal B^m(1-w)^{-bm}c_*^{2s}D_T.
\end{equation}

Write $a_*=J-m$ and $s_*=ma_*+m(m-1)/2$, and set
$d_i=a_*+i-j_i$ for $0\leq i<m$. The inequalities on the $j_i$ imply
\[
 a_*\geq d_0\geq\cdots\geq d_{m-1}\geq0.
\]
Thus $d$ is a partition, $s=s_*-|d|$, and the dimension factor in
\eqref{eq:gamma-minor-bound} is $s_d(1^m)$.
The ratio $(b)_{2j_i-i}/(b)_{2a_*+i}$ removes $2d_i$ factors, all
at least $b$. Combining Lemma~\ref{lem:gamma-schur} with
\eqref{eq:schur-gamma-reduction} therefore gives
\[
 \frac{\det G[T,T]}
 {\mathcal B^m(1-w)^{-bm}c_*^{2s_*}
                         \prod_{i=0}^{m-1}i!(b)_{2a_*+i}}
 \leq s_d(1^m)^2 z_*^{|d|}.
\]
The assumption $bc_*>1$ says $0<z_*<1$. Sum over the order sets $T$
and enlarge the sum to all partitions of length at most $m$.
The Schur Cauchy identity, specialized to two lists of $m$ variables
all equal to $\sqrt{z_*}$, gives
\[
 \sum_d s_d(1^m)^2z_*^{|d|}=(1-z_*)^{-m^2}.
\]
This is precisely the remaining factor in \eqref{eq:schur}.
\end{proof}

\subsection{The signed row sum}
We now estimate the exact formulas at finite $R$ and $S$. Throughout,
$R=KS$, $1\leq J\leq R$, and the weights are
\eqref{eq:row-weights}, with $a>0$, $0<v<1$, and $\beta<0$.
The estimates below keep the row sum separate, since it can be negative.

Write $L(x)=x\log x$, with $L(0)=0$, and define
\[
 f(x)=L(x)-L(a+x)+L(a)-x\log v,
 \qquad C_v=-a\log(1-v).
\]
\begin{lemma}\label{lem:spectral-row-finite}
Suppose $aR>1$ and $f(1)<0$. Then
\begin{equation}\label{eq:row-finite}
 \frac1W\sum_\Lambda\lambda_{r,s}\leq
 3\int_0^1(1-x)f(x)\,\mathrm dx+\frac{9C_v}{R}
 +\frac\beta{2K}(1-S^{-2})
 +\frac{3(1+S^{-1})}{2R}\delta_R,
\end{equation}
where
\[
 \delta_R=\log(R+1)+1+\log\frac{a+1}{a-1/R}.
\]
\end{lemma}
\begin{proof}
Integral comparison gives, for $r\geq0$ and $b>1$,
\[
 \log(r!)\leq r\log r-r+\log(r+1)+1,
 \qquad
 \log((b)_r)\geq\int_0^r\log(b+t)\,\mathrm dt
                        -\log\frac{b+r}{b-1}.
\]
For the second inequality, compare the sum with the integral over
$[-1,r-1]$ and bound the two endpoint intervals. Substituting $b=aR$
in the row weights yields
$\lambda_{r,s}\leq Rf(r/R)+\beta s+\delta_R$.

The function $f$ is convex, $f(0)=0$, and its minimum on $[0,\infty)$
is $-C_v$. Thus $f\leq0$ on $[0,1]$ and its total variation there is
at most $2C_v$. The function $(1-x)f(x)$ has total variation at most
$3C_v$: apply the product bound for variation, using $|f|\leq C_v$.
There are $S-\lfloor r/K\rfloor\geq S(1-r/R)$ rows with first
exponent $r$. Since $f$ is nonpositive, their contribution is at most
$RS(1-r/R)f(r/R)$. The left Riemann sum differs from the integral
by at most $3C_v/R$, giving the first two terms of
\eqref{eq:row-finite} after division by $W=R^2S/3$.

Finally,
$\sum_\Lambda s=KS(S^2-1)/6$ and
$|\Lambda|=R(S+1)/2$. Summing $\beta s+\delta_R$ gives the last two
terms of \eqref{eq:row-finite}.
\end{proof}
The integral in \eqref{eq:row-finite} is elementary; it involves only
integrals of $x\log x$ and $x^2\log x$, after translation by $a$.

\subsection{Column capacities and the forward estimate}
Put $\varepsilon=J/R$. For a fixed upper squared radius $0<q\leq1$
and a column scale $\tau>0$, define
\[
 g_\tau(x)=L(x)+L(a+x)-L(a)
                  +2x(\ell-\log\tau-1),
 \qquad E_{\rm col}=\frac{\log(R+1)+1}{R}.
\]
The price \eqref{eq:column-prices} satisfies
\begin{equation}\label{eq:column-finite-error}
 \frac{\nu_j}{R}\leq g_\tau(j/R)+E_{\rm col}.
\end{equation}
Indeed, use the factorial upper bound above and
$\log((aR)_j)\leq\int_0^j\log(aR+x)\,\mathrm dx$.
The terms containing $\log R$ cancel those in the column scale, and
$j+1\leq J\leq R$ bounds the factorial remainder.

The sum of the largest $m$ column prices can be bounded without deciding
which orders occur. Partition $[0,\varepsilon]$ into $N_z$ equal bins.
For a bin $I=[u,v]$, give it the upper cost
$\gamma_I=\max\{0,g_\tau(u),g_\tau(v)\}$ and capacity
\[
 c_z=\frac{\varepsilon}{N_z}+\frac1R.
\]
The function $g_\tau$ is convex, so $\gamma_I$ bounds it throughout
$I$. Each bin contains at most $Rc_z$ of the possible points $j/R$.
Assign boundary points to just one neighboring bin, and define
\[
 C_{q,\tau}(t)
 =\max\left\{\sum_I\gamma_I u_I:
                0\leq u_I\leq c_z,\quad\sum_Iu_I=t\right\}.
\]
This function is nondecreasing and piecewise linear. Filling bins in
decreasing order of $\gamma_I$ evaluates it. Every selection of $m$
orders gives a feasible choice at $t=m/R$; hence
\[
 \frac{\nu_m^*(\tau)}{R^2}\leq C_{q,\tau}(t)+tE_{\rm col}.
\]
The replacement of negative bin costs by zero ensures monotonicity and
only enlarges the upper bound.

Retain $\delta_q,v_q,w,\mathcal B$ from
\eqref{eq:moment-matrix}. For the forward choice $0<\xi<1$, take
$\Gamma,\tau$ as in Section~\ref{sec:metric}, and set
\[
 d_{q,\xi}=a\log\frac\Gamma{1-w}+\frac{\delta_q}{K}
             +\frac{\log R}{R}+E_{\rm col}-\frac{\log\xi}{R}.
\]
Since $\log\mathcal B/R\leq\log R/R+\delta_q/K$, the forward
bound \eqref{eq:forward-cost} gives
\begin{equation}\label{eq:forward-finite}
 \frac{F_q^{\rm f}(m)}{R^2}
 \leq C_{q,\tau}(t)+d_{q,\xi}t+t^2\log\xi.
\end{equation}
The last term in $d_{q,\xi}$ retains the difference between
$m(m-1)$ and $m^2$.

To use \eqref{eq:H-occupancy}, we need a bound throughout an interval
of possible values of $t$, rather than only at its endpoints. On a cell
$[l,h]\subseteq[0,\varepsilon]$, a bound for
$F_q^{\rm f}(m)/R^2-a_0t$ is
\begin{equation}\label{eq:forward-cell}
 C_{q,\tau}(h)
 +\max\{(d_{q,\xi}-a_0)l,(d_{q,\xi}-a_0)h\}
 +l^2\log\xi.
\end{equation}
The three terms follow respectively from the monotonicity of
$C_{q,\tau}$, the sign of the linear coefficient, and $\log\xi<0$.

\subsection{Finite form of the complementary estimate}
For $\kappa>0$, put
\[
 d=\frac{\kappa}{1+\sqrt w},\qquad c_0=d\sqrt w,
 \qquad \tau=\frac{d(1-w)}{\sqrt{v_q}}.
\]
These are compatible with the affine map \eqref{eq:affine-map}.
Choose $0<r_0<1$ and use $G_0,s_0$ from
Lemma~\ref{lem:complementary}. Define
\begin{align*}
 c_1&=-a\log G_0-\varepsilon
                  (2\log r_0-\log s_0-|\log s_0|),\\
 c_2&=\log r_0-|\log s_0|<0,\\
 e_q&=-a\log(1-w)+\frac{\delta_q}{K}
          +\frac{\log R}{R}+E_{\rm col}
          +\frac{\max(0,-\log d)}{R}.
\end{align*}
The complementary bound for the node minor, divided by $R^2$, is at most
\begin{equation}\label{eq:complement-finite}
 C_{q,\tau}(t)+\varepsilon(a+\varepsilon)\log G_0
                          +(c_1+e_q)t+c_2t^2.
\end{equation}
To verify the finite term, substitute $J=\varepsilon R$, $m=tR$ and
$k=(\varepsilon-t)R$ into \eqref{eq:complementary-cost}. Using
$\log s_0=\log d+\log r_0-\log G_0$, both branches become
\[
 \varepsilon(a+\varepsilon)\log G_0+c_1t+c_2t^2
       -\frac{t\log d+(\varepsilon-t)\log G_0}{R}.
\]
Here $G_0\geq1$ and $0\leq t\leq\varepsilon$, so the last term is at
most $t\max(0,-\log d)/R$. Restoring the mass and column prices gives
\eqref{eq:complement-finite}.

On a cell $[l,h]$, subtract $a_0t$ and use
\begin{equation}\label{eq:complement-cell}
 \begin{split}
 C_{q,\tau}(h)+\varepsilon(a+\varepsilon)\log G_0
 &+\max\{(c_1+e_q-a_0)l,(c_1+e_q-a_0)h\}\\
 &+c_2l^2.
 \end{split}
\end{equation}
As in the forward estimate, every term bounds the whole cell.

\subsection{The finite error in the Schur estimate}
At a positive upper squared radius $q$, put
\[
 \zeta=\ell+\tfrac12\log v_q-\log(1-w),\qquad
 z_*=a^{-2}e^{-2\zeta},\qquad
 I(x)=\tfrac12x^2\log x-\tfrac14x^2,
\]
with $I(0)=0$. In \eqref{eq:schur}, the parameters of the completed
Gram matrix satisfy $c_*=e^\zeta/R$ and $bc_*=ae^\zeta$.
Define
\begin{align}
 g_\varepsilon(t)={}&t\bigl(\delta_q/K-a\log(1-w)\bigr)
          +(2\varepsilon t-t^2)(\zeta-1)+I(t)\notag\\
 &+I(a+2\varepsilon-t)-I(a+2\varepsilon-2t)
               -ta\log a-t^2\log(1-z_*).
 \label{eq:schur-finite}
\end{align}

\begin{proposition}\label{prop:schur-finite-error}
Suppose $R>3$, $0<a\leq1$, $0<\varepsilon=J/R\leq1$,
$\varepsilon(a+2\varepsilon)\leq1$, and $ae^\zeta>1$.
For $0\leq m\leq J$ and $t=m/R$, the logarithm of the right side of
\eqref{eq:schur}, divided by $R^2$, is at most
\begin{equation}\label{eq:schur-error}
 g_\varepsilon(t)+E_0,\qquad
 E_0=\frac{\varepsilon}{R}
           (2\log R+2-\log\varepsilon-\log a).
\end{equation}
\end{proposition}
\begin{proof}
The assertion is immediate for $m=0$. For $m\geq1$, put $a_*=J-m$,
as in \eqref{eq:schur}. In the product there, use
\[
 \log(i!)\leq i\log i-i+\log(i+1)+1,\qquad
 \log((aR)_n)\leq\int_0^n\log(aR+x)\,\mathrm dx.
\]
For $i=0,\ldots,m-1$, set $x_i=i/R$ and take $n=2a_*+i$.
The terms containing $\log R$ have total coefficient
$\sum_i(2a_*+2i)=2s_*$, and cancel those in
$c_*^{2s_*}=\exp(2s_*(\zeta-\log R))$.
The remaining sum, divided by $R^2$, is the left Riemann sum
$R^{-1}\sum_{i=0}^{m-1}f_t(x_i)$, where
\[
 \begin{split}
 f_t(x)={}&x\log x+(a+2\varepsilon-2t+x)
                         \log(a+2\varepsilon-2t+x)\\
          &-a\log a-2\varepsilon+2t-2x.
 \end{split}
\]

The other finite terms have simple upper bounds. Since $a\leq1$ and
$ae^\zeta>1$, we have $\zeta>0$. Thus replacing
$2s_*/R^2=2\varepsilon t-t^2-t/R$ by $2\varepsilon t-t^2$ drops the
nonpositive term $-t\zeta/R$. The factor $\mathcal B^m$ contributes
$t\delta_q/K+t\log S/R$, whose second term is at most
$\varepsilon\log R/R$. Since $i+1\leq J\leq R$, the factorial
remainders contribute at most $\varepsilon(\log R+1)/R$.

It remains to estimate the Riemann sum. For $0<x\leq t$,
\[
 f_t'(x)=\log\bigl(x(a+2\varepsilon-2t+x)\bigr)\leq0,
\]
because $x\leq t\leq\varepsilon$ and
$\varepsilon(a+2\varepsilon)\leq1$. Its total variation is at most
\[
 -\int_0^t(\log x+\log a)\,\mathrm dx
 =t(1-\log t-\log a)
 \leq\varepsilon(1-\log\varepsilon-\log a).
\]
The last inequality follows by differentiation, since $at\leq1$.
For a decreasing function, the left Riemann sum exceeds the integral
by at most its variation divided by $R$. Finally, direct integration
gives
\[
 \begin{split}
 \int_0^t f_t(x)\,\mathrm dx
 ={}&I(t)+I(a+2\varepsilon-t)-I(a+2\varepsilon-2t)\\
    &-ta\log a-2\varepsilon t+t^2.
 \end{split}
\]
Combining this integral with the mass, $c_*^{2s_*}$, and
$(1-z_*)^{-m^2}$ gives \eqref{eq:schur-finite}. The three error bounds
sum to $E_0$, proving \eqref{eq:schur-error}.
\end{proof}

To obtain a bound on a whole cell, we use concavity. Differentiating
\eqref{eq:schur-finite} twice gives, for $t>0$,
\[
 g_\varepsilon''(t)=
 \log\frac{t(a+2\varepsilon-t)}
 {e^{2\zeta}(a+2\varepsilon-2t)^4(1-z_*)^2}.
\]
In particular, the sufficient condition
\begin{equation}\label{eq:schur-concavity}
 \varepsilon(a+\varepsilon)z_*<a^2(1-z_*)^2
\end{equation}
makes $g_\varepsilon$ concave on $[0,\varepsilon]$. Indeed, the
numerator is at most $\varepsilon(a+\varepsilon)$, the fourth power
in the denominator is at least $a^4$, and
$e^{2\zeta}=1/(a^2z_*)$.
For a cell with midpoint $t_0>0$ and half-width $h_0$, its tangent
therefore gives the upper bound
\begin{equation}\label{eq:schur-cell}
 g_\varepsilon(t_0)-a_0t_0
        +h_0|g_\varepsilon'(t_0)-a_0|+E_0
\end{equation}
for $g_\varepsilon(t)-a_0t+E_0$ throughout the cell. This also applies
to the first cell, whose left endpoint is zero, by continuity.

\subsection{Combining the finite estimates}
Partition $[0,\varepsilon]$ into $N_m$ cells. At each positive upper
radius $q$, the forward and complementary formulas give bounds on every
cell. Include \eqref{eq:schur-cell} whenever the hypotheses of
Proposition~\ref{prop:schur-finite-error} and
\eqref{eq:schur-concavity} hold. For each cell, take the minimum over
the available bounds and their parameter choices. Let $\Phi(q)$ be
the maximum of these cell bounds and zero. Each chosen bound holds
throughout its cell. Taking $F_q(m)$ to be the minimum of the
corresponding exact node bounds, we obtain
\[
 F_q(m)\leq a_0Rm+R^2\Phi(q)\qquad(0\leq m\leq J).
\]
Thus $\Phi$ satisfies \eqref{eq:H-occupancy}. Radial monotonicity
(Lemma~\ref{lem:radial}) and the node count in
\eqref{eq:metric-mesh} now give the upper bound for the full minor.
The signed row sum is supplied by \eqref{eq:row-finite}.

For the finite comparisons, the additional parameter choices are
\[
 \kappa\in\{5/2,3,4,6,10\},
\]
\begin{multline*}
 \xi,r_0\in\{1/100,3/100,3/50,1/10,3/20,1/5,3/10,2/5,\\
                 1/2,3/5,7/10,4/5,9/10,19/20,99/100,9999/10000\}.
\end{multline*}
The estimates themselves apply to any finite collections satisfying
$\kappa>0$ and $0<\xi,r_0<1$. They may also use an upper bound for
$\ell$: each fixed principal minor is multiplied by
$e^{2\ell\sum_{j\in T}j}$, so increasing $\ell$ enlarges it.
All radial upper endpoints are positive; the original node minors at
$q=0$ are covered by Lemma~\ref{lem:radial}.

\section{Resultants and full rank}\label{app:resultants}
For the range $D\geq875$ we use a rank criterion that does not require
an auxiliary inert prime. The nodes lie on a small number of horizontal
lattice lines. A polynomial of small weighted degree can vanish at few
points of those lines unless it contains one of them as a factor;
a resultant bounds the additional multiplicity at the remaining points.

For a polynomial $F(z,w)$, write
\[
 \deg_{S,R}F=\max\{Si+Rj:z^iw^j\text{ occurs in }F\}.
\]
This degree is additive under multiplication, since the highest weighted
homogeneous parts multiply to a nonzero polynomial. A polynomial supported
on $\Lambda$ has weighted degree at most $S(R-1)<RS$.

\begin{lemma}\label{lem:weighted-resultant}
Let $F\in\CC[z,w]$ be irreducible, of $z$-degree $d\geq1$ and weighted
degree $H=\deg_{S,R}F$. Then $\mathcal R(w)=\operatorname{Res}_z(F,\partial_zF)$
is nonzero and
\[
 \deg\mathcal R\leq\frac{(2d-1)H-Sd^2}{R}
                  \leq\frac{H^2}{RS}.
\]
If $F(z,w_0)$ has a root of multiplicity $m$ at $z_0$, then
$\mathcal R$ vanishes to order at least $m-1$ at $w_0$.
\end{lemma}
\begin{proof}
Write $F=\sum_{i=0}^d c_i(w)z^i$. Irreducibility and characteristic zero
imply that $F$ and $\partial_zF$ are relatively prime over $\CC(w)$.
Their resultant is therefore nonzero. Also
$R\deg c_i\leq H-Si$.

Each monomial in the Sylvester determinant has coefficient degree
$2d-1$, and the sum of its coefficient indices is $d^2$.
For example, the latter follows by replacing $z$ by $tz$: the resultant
of $F(tz,w)$ and its $z$-derivative is $t^{d^2}\mathcal R(w)$.
Thus each term has $w$-degree at most $((2d-1)H-Sd^2)/R$.
The second inequality follows from
\[
 H^2-S\bigl((2d-1)H-Sd^2\bigr)=(H-Sd)^2+SH\geq0.
\]

It remains to justify the multiplicity statement, including when the
leading coefficient vanishes at $w_0$. The specialization $F(z,w_0)$
is not identically zero: otherwise $w-w_0$ would divide the irreducible
polynomial $F$, contrary to $d\geq1$. Every row of its specialized
Sylvester matrix is a polynomial divisible by $(z-z_0)^{m-1}$.
The $m-1$ functionals given by evaluation of the first $m-1$ Taylor
coefficients at $z_0$ therefore annihilate its row space. They are
independent, so the matrix has corank at least $m-1$.
For a polynomial matrix, corank at least $k$ at $w_0$ forces its
determinant to be divisible by $(w-w_0)^k$: constant row and column
operations put its value at $w_0$ in rank normal form, and each of the
remaining rows supplies a factor $w-w_0$ in the determinant expansion.
This proves the assertion without a condition on $c_d(w_0)$.
\end{proof}

Let the nodes lie in $|z|\leq\rho$. Write $M_y$ for an upper bound
on the number of horizontal lattice lines meeting this disk, and $M_x$
for an upper bound on the number of nodes on any one such line. We may take
\[
 M_y=2\lfloor2\rho/\sqrt D\rfloor+1,
 \qquad M_x=\lfloor2\rho\rfloor+1.
\]
On a line of imaginary coordinate $y\sqrt D/2$ we have
$w=\bar z=z-y\sqrt{-D}$.

\begin{lemma}\label{lem:lattice-line-multiplicity}
Let $F$ be as in Lemma~\ref{lem:weighted-resultant}. At a node $z_i$,
let $m_i$ be the multiplicity of $z_i$ as a root of $F(z,\bar z_i)$.
If $F$ is a constant multiple of $z-w-y\sqrt{-D}$ for some integer $y$,
then $H=R$ and $\sum_i m_i\leq M_x$. Otherwise
\[
 \sum_i m_i\leq\frac{M_yH}{S}+\frac{H^2}{RS}.
\]
\end{lemma}
\begin{proof}
A horizontal-line factor has simple $z$-roots and vanishes at nodes of
one line only. In the other case, irreducibility implies that
$F(z,z-y\sqrt{-D})$ is nonzero for every integer $y$.
Its ordinary degree is at most $H/S$, since $S\leq R$.
There are consequently at most $M_yH/S$ nodes at which $F$ vanishes.

For the excess multiplicities, apply Lemma~\ref{lem:weighted-resultant}
at $w=\bar z_i$. These values of $w$ are distinct, because conjugation
is injective. Hence $\sum_i(m_i-1)_+\leq\deg\mathcal R\leq H^2/(RS)$.
Adding the two bounds proves the result.
\end{proof}

\begin{proposition}[Resultant rank criterion]\label{prop:resultant-rank}
The interpolation matrix $V$ has full row rank if
\begin{equation}\label{eq:resultant-rank}
 nJ>R(S+M_y),\qquad \max(J,M_x)\leq R+KM_y.
\end{equation}
\end{proposition}
\begin{proof}
Otherwise there is a nonzero polynomial $F$ supported on $\Lambda$
such that $F(z,\bar z_i)$ vanishes to order at least $J$ at $z_i$ for
every node. Factor it over $\CC$ as
\[
 F(z,w)=F_0(w)\prod_\ell F_\ell(z,w),
\]
listing repeated irreducible factors separately. Put $d_0=\deg F_0$.
Let $k$ count the horizontal-line factors and let $H_*$ be the sum of
the weighted degrees of all other factors involving $z$. Additivity gives
\[
 R(d_0+k)+H_*<RS.
\]
A root of $F_0$ affects at most one node. Truncating each node's
multiplicity at $J$, these roots contribute at most $Jd_0$.
Lemma~\ref{lem:lattice-line-multiplicity} and
$\sum H_\ell^2\leq(\sum H_\ell)^2$ give
\[
 nJ\leq Jd_0+M_xk+\frac{M_yH_*}{S}+\frac{H_*^2}{RS}.
\]
Write $\beta=H_*/(RS)$, so $0\leq\beta<1$ and
$d_0+k<(1-\beta)S$. The right side is strictly less than
\[
 (1-\beta)S\max(J,M_x)+\beta M_yR+\beta^2RS
 \leq(1-\beta)R(S+M_y)+\beta R(S+M_y).
\]
This contradicts the first inequality in \eqref{eq:resultant-rank}.
\end{proof}

\section{A radial estimate for the larger discriminants}\label{app:radial}
For the range $D\geq875$ we use an entrywise estimate in place of the
spectral estimate of Section~\ref{sec:metric}. The two estimates start from
the same normalized matrix. Here the additional saving comes from keeping
track of which nodes and derivative orders can occur together in a
determinant term. A node used many times must supply smaller derivative
orders, while nodes near the origin give stronger radial decay. We express
these restrictions by a convex function of the numbers of columns selected
at the different radii.

Throughout this appendix assume $K\geq2$ and $S\geq1$. Put $R=KS$, $M=R(S+1)/2$, $W=R^2S/3$, and
$\theta=M/(nJ)\leq1$. The derivative scale is denoted by $\ell$, as in
\eqref{eq:residual-matrix}. We assume $\ell\geq0$ and put
$\varepsilon=J/R$. The parameters in Appendix~\ref{app:global} satisfy
these assumptions.

\subsection{Counting the ordered exponents}
Put
\[
 u(G)=1-\sqrt{1-G},\qquad
 v(G)=\begin{cases}
 \sqrt{G/K},&0\leq G\leq1/K,\\
 1-\sqrt{(1-1/K)(1-G)},&1/K\leq G\leq1.
 \end{cases}
\]
The second function is the quantile of $r+s$, after scaling by $R$, on
the continuous triangle $r+Ks\leq R$. The following bounds account for
the discrete rows of the interpolation matrix.

\begin{lemma}\label{lem:radial-quantiles}
List the values of $r$ and $r+s$ on $\Lambda$ in nondecreasing order as
$r_{(k)}$ and $z_{(k)}$, respectively. For $G_k=k/M$,
\[
 Ru(G_k)-1\leq r_{(k)}\leq(R+K)u(G_k),\qquad
 z_{(k)}\geq Rv(G_k)-2.
\]
\end{lemma}
\begin{proof}
For $0\leq a\leq R$, the number of rows with $r\geq a$, with $a$ an
integer, is
\[
 T_r(a)=\sum_{b=1}^S(Kb-a)_+
       =\frac K2\{y^2+y+\alpha(1-\alpha)\},
 \quad y=S-a/K,\quad\alpha=\{a/K\}.
\]
This follows by summing the arithmetic progression of positive terms.
It implies
\[
 M(1-a/R)^2\leq T_r(a)
       \leq M(1-a/(R+K))^2.
\]
For the first inequality, subtracting its right-hand quadratic from
$y^2+y$ leaves $y(1-y/S)\geq0$.
For the second, write $a/K=q+\alpha$. After multiplication by $S+1$,
the difference of the two sides, apart from the positive factor $K/2$,
is $q(S+1-q-2\alpha)+S\alpha^2\geq0$.
The endpoints $a=0,R$ satisfy the same inequalities.
Applying the lower tail bound at $a=r_{(k)}+1$ and the upper tail bound
at $a=r_{(k)}$ proves the two assertions about $r_{(k)}$.

For the joint exponent, let $N(t)$ count rows with $r+s\leq t$, for an
integer $0\leq t<R$. Let $F$ be the inverse of $v$, extended by $F(x)=1$
for $x\geq1$. We show that $N(t)\leq MF((t+2)/R)$.
If $t\leq S-2$, the whole triangle $r+s\leq t$ is contained in
$\Lambda$, so
\[
 N(t)=\frac{(t+1)(t+2)}2
       \leq\frac{S+1}{2S}(t+2)^2=MF((t+2)/R).
\]
If $t\geq S-1$, put $h=K-1$ and $A=R-t-1$.
For $A\geq1$ the complementary count is
\[
 M-N(t)=\sum_{s=0}^{S-1}(A-hs)_+
       \geq\frac{A^2}{2h}+\frac A2.
\]
Here $1\leq A\leq hS$. On this interval
\[
 A^2+hA\geq(1+1/S)(A-1)^2:
\]
the difference is a concave quadratic in $A$ and is nonnegative at both
endpoints. The last lower bound is therefore at least
$M\{1-F((t+2)/R)\}$. When $A=0$ the desired assertion is immediate.
Taking $t=z_{(k)}$ and using $k\leq N(t)$ proves the remaining bound.
\end{proof}

\subsection{The falling factorial and the derivative orders}
For $0\leq t\leq x\leq1$, define
\[
 f(x,t)=x\log x-(x-t)\log(x-t)-t,
 \qquad
 g(x,t)=\max_{0\leq a\leq\min(x,t)}\{a\ell+f(x,a)\},
\]
with $0\log0=0$. We also use $g(x,t)$ for $t>x$ with this same
definition. Its maximum occurs at
$a=\max(0,\min(t,x-e^{-\ell}))$.
For $0\leq j\leq r<R$, integral comparison gives
\begin{equation}\label{eq:radial-factorial}
 \log\frac{\fall rj}{R^j}
 \leq Rf\left(\frac{r+1}{R},\frac jR\right).
\end{equation}
Indeed, on each unit interval the decreasing function
$\log((r+1-t)/R)$ is at least its value at the right endpoint.
This also proves the inequality when $j=r$.

Fix a set $C$ of $M$ columns, and write $c_i$ for the number selected at
node $i$. Replace its selected orders by $J-c_i,\ldots,J-1$, and list
the resulting multiset increasingly as $j_k^+$. If the original orders
are listed as $j_{(k)}$, then
\begin{equation}\label{eq:radial-order-envelope}
 j_{(k)}\leq j_k^+\leq t_k,
 \qquad t_k=J-\frac{M-k}{n}.
\end{equation}
The first assertion follows by ordering the columns at each node; the
second follows because at most $n(J-a)$ orders can be at least $a$.

Consider a nonzero term in the expansion of a selected determinant.
Every assigned derivative satisfies $j\leq r$. The mixed derivative
$f_{xt}=1/(x-t)$ is positive. Exchanging two oppositely ordered pairs
therefore increases the sum of $f$; the exchange preserves feasibility,
since the larger derivative was already assigned to the smaller row degree.
Successive exchanges pair the increasingly ordered row degrees and
derivatives. Increasing the row bound and then using that $g$ is
nondecreasing in its second argument gives
\begin{equation}\label{eq:radial-rearrangement}
 \sum_C\left(j\ell+
       \log\frac{\fall rj}{R^j}\right)
 \leq R\sum_{k=1}^M g(x_k,j_k^+/R),
 \qquad x_k=\min\left(1,(1+1/S)u(G_k)+1/R\right).
\end{equation}
The exchange argument extends to the boundary $x=t$ by continuity.

Put $h(x,t)=t\ell-g(x,t)$. For $\ell\geq0$, this function is
nonnegative, increases with $t$ and $\ell$, and decreases with $x$.
For example $g_t=(\ell+\log(x-t))_+$ before the maximum becomes
constant, so $0\leq g_t\leq\ell$. The other assertions follow either
by differentiation at the maximizing argument or directly from the
definition. Thus an admissible normalized factorial saving is any
number $S_{\rm der}$ satisfying
\begin{equation}\label{eq:radial-derivative-saving}
 0\leq S_{\rm der}\leq\frac3M\sum_{k=1}^M h(x_k,t_k/R).
\end{equation}
It can be bounded by grouping consecutive indices: on a group with
$x_k\leq b$ and $t_k/R\geq a$, use its relative size times $3h(b,a)$.
This requires only a finite sum, rather than a list of the $M$ rows.

\subsection{An additional saving from the number of columns at a node}
Suppose $\theta\varepsilon<1/2$. Choose integers
$0<m_1<\cdots<m_L<M/n$, put $\tau_r=m_r/J$, and suppose
\[
 a_r=1-\sqrt{\tau_r/\theta}-(1-\tau_r)\varepsilon>0,
 \qquad \nu_r=\ell+\log a_r>0.
\]
These $\nu_r$ decrease with $r$: the derivative of
$1-\sqrt{\tau/\theta}-(1-\tau)\varepsilon$ is
$\varepsilon-1/(2\sqrt{\theta\tau})<0$ for $\tau<\theta$.
Set $\nu_{L+1}=0$ and $\lambda_r=\nu_r-\nu_{r+1}$.
Smaller nonnegative increments may be used throughout.

If $t_k\geq J-m_r$, then $G_k\geq1-\tau_r/\theta$.
The function
$u(G)-\varepsilon(1-\theta(1-G))$ is increasing, since its derivative
is at least $1/2-\theta\varepsilon>0$.
It follows that $x_k-t_k/R\geq a_r$. On the interval between
$j_k^+$ and $t_k$, the derivative of $Rg(x_k,t/R)$ is therefore at
least the staircase
$\sum_r\lambda_r\boldsymbol1_{t\geq J-m_r}$.
Integration gives an improvement to \eqref{eq:radial-rearrangement}.
The improvement can be summed explicitly, because
\begin{align*}
 \sum_k(t_k-J+m)_+&=\frac{nm^2+m}{2},\\
 \sum_k(j_k^+-J+m)_+
 &=\frac{nm(m-1)}2
       -\frac12\sum_i(m-c_i)_+(m-c_i-1).
\end{align*}
Subtracting these identities, and using
$\sum_i(m-c_i)_+\leq nm$, gives a difference at least
$\frac12\sum_i(m-c_i)_+^2$. Consequently
\begin{equation}\label{eq:radial-derivative-coupled}
 \frac2W\sum_C\left(j\ell+\log\frac{\fall rj}{R^j}\right)
 \leq\gamma\ell-S_{\rm der}
       -\frac1W\sum_{i,r}\lambda_r(m_r-c_i)_+^2,
 \quad
 \gamma=\frac2W\left(MJ(1-\theta/2)+\theta J/2\right).
\end{equation}
In obtaining the factorial saving we used $2RM/W=3(1+1/S)\geq3$.

\subsection{Radial weights and the convex estimate}
Order the nodes by modulus. Suppose $N$ divides $n$, and divide the
nodes into $N$ consecutive groups of $n/N$ nodes. Choose numbers
$L_0\geq\cdots\geq L_{N-1}\geq L_N=0$ such that
\[
 \log(\rho/|z_i|)\geq L_b\qquad
       (bn/N<i\leq(b+1)n/N).
\]
At a zero node the assertion has its limiting meaning. The disk estimate
\eqref{eq:node-disk} permits the choice
\begin{equation}\label{eq:radial-bin-weights}
 L_b=\log\frac{1+\eta}
 {\sqrt{\min(1,(b+1)/N+1/n)}+\eta},
 \qquad \eta=\frac{\sqrt{D+4}/4}{\sqrt{n\sqrt D/(2\pi)}}.
\end{equation}
The extra $1/n$ weakens the bound and is convenient when a common bound
is used for several values of $n$.
Choose an index $b_0$ and a number $\overline L\geq L_{b_0}$.
Write
\[
 w_b=\frac{N}{nJ}\sum_{i\text{ in group }b}c_i,
 \qquad G_b(w)=\frac1{N\theta}\sum_{a=0}^bw_a.
\]
Thus $0\leq w_b\leq1$ and $\sum_bw_b=N\theta$.
Define
\begin{align}
 \Phi(G)&=3\int_0^G\left(v(t)-\frac{J+2}{R}\right)_+\,dt,
                                                        \label{eq:radial-Phi}\\
 \Delta&=\frac2W\left\{\binom{J+2}{3}
                         -\binom{(J-S)_++2}{3}\right\}.
                                                        \label{eq:radial-Delta}
\end{align}
In particular, when $J\geq S$,
\[
 \Delta=3\left(\frac{J+1}{R}\right)^2
       -\frac3K\frac{J+1}{R}+\frac1{K^2}-\frac1{R^2}.
\]
On the polytope $\mathcal P_\theta=\{w\in[0,1]^N:\sum_bw_b=N\theta\}$,
put
\begin{align}
 \mathcal F(w)={}&\sum_{b=0}^{N-1}(L_b-L_{b+1})\Phi(G_b(w))
                                                     \notag\\
 &+\frac{3\varepsilon}{2N\theta}\sum_{b=0}^{N-1}
   \left\{L_{\max(b,b_0)}w_b^2
             +\sum_{r=1}^L\lambda_r(\tau_r-w_b)_+^2\right\}
       -\overline L\Delta.                 \label{eq:radial-functional}
\end{align}

\begin{proposition}[Radial determinant estimate]\label{prop:radial-size}
With the preceding choices, let $F_*\leq\min_{\mathcal P_\theta}\mathcal F$.
Then every maximal minor satisfies
\begin{equation}\label{eq:radial-minor-bound}
 \frac1W\log|\det B_C(\ell)|^2
 \leq\gamma\ell-S_{\rm der}-F_*+\frac{2M\log M}{W}.
\end{equation}
\end{proposition}
\begin{proof}
Consider a nonzero determinant term, and put $z=r+s$ for its assigned
row. Since $z\geq r\geq j$, for a radial weight $l\geq0$ and
$\widetilde l=\min(l,\overline L)$ one has
\[
 (z-j)l\geq(z-J)_+l+(J-j)\widetilde l
                               -(J-z)_+\widetilde l.
\]
Indeed, subtracting the first term leaves
$(\min(z,J)-j)l$, whose coefficient is nonnegative.
This observation allows us to retain the full radial weight on the
first term while bounding the possible negative correction by
$\overline L$.

Among all assignments of rows to columns, pairing the smallest values
of $(r+s-J)_+$ with the largest radial weights minimizes their sum.
For any integer $h\leq M$, Lemma~\ref{lem:radial-quantiles} and the
right-endpoint integral inequality for an increasing function give
\[
 \frac2W\sum_{k=1}^h(z_{(k)}-J)_+
 \geq\frac{2RM}{3W}\Phi(h/M)\geq\Phi(h/M).
\]
Summation by parts over the node groups therefore gives the first sum
in \eqref{eq:radial-functional}.
At node $i$, the distinct derivative orders satisfy
$\sum(J-j)\geq c_i(c_i+1)/2\geq c_i^2/2$.
Also
\[
 \sum_{\Lambda}(J-r-s)_+
 \leq\sum_{s=0}^{S-1}\binom{(J-s)_++1}{2}
 =\binom{J+2}{3}-\binom{(J-S)_++2}{3}.
\]
Thus the clipped radial terms contribute at least
$W^{-1}\sum_i c_i^2L_{\max(b(i),b_0)}-\overline L\Delta$.

Combine this with \eqref{eq:radial-derivative-coupled}.
Convexity on each node group gives
\[
 \sum_{i\text{ in group }b}c_i^2\geq\frac nN J^2w_b^2,
 \qquad
 \sum_{i\text{ in group }b}(m_r-c_i)_+^2
       \geq\frac nN J^2(\tau_r-w_b)_+^2.
\]
Since $nJ^2/W=3\varepsilon(1+1/S)/(2\theta)$, dropping the factor
$1+1/S$ proves that the total normalized saving is at least
$S_{\rm der}+\mathcal F(w)$. The modulus of each nonzero determinant
term is bounded accordingly. There are at most $M!\leq M^M$ terms;
taking the square of their sum proves \eqref{eq:radial-minor-bound}.
\end{proof}

The function $\Phi$ is convex, because its derivative is increasing.
Every coefficient of the convex terms in \eqref{eq:radial-functional}
is nonnegative. Hence $\mathcal F$ is convex. In particular, any
$w^*\in\mathcal P_\theta$ gives the explicit lower bound
\begin{equation}\label{eq:radial-supporting-plane}
 F_*=\mathcal F(w^*)-\nabla\mathcal F(w^*)\cdot w^*
                 +\min_{w\in\mathcal P_\theta}
                         \nabla\mathcal F(w^*)\cdot w.
\end{equation}
The last minimum is found by filling the coordinates with the smallest
coefficients first, subject to the upper bound one and the prescribed
sum. This follows by transferring mass from a larger coefficient to a
smaller one whenever the latter coordinate is not full. No assertion
that $w^*$ minimizes $\mathcal F$ is required.
For numerical evaluation its gradient is obtained directly from
\eqref{eq:radial-functional}; the positive-part squares are continuously
differentiable, including at their endpoints.

\section{The range \texorpdfstring{$D\geq875$}{D >= 875}}\label{app:global}
We now apply the radial estimate to the part of the proof preceding the
finite comparisons of Section~\ref{sec:completion}. The argument has two
parts. On an unbounded interval the arithmetic estimate grows faster
than the size estimate, with an explicit positive difference at the left
endpoint. Below that endpoint we keep the integer parameters fixed on
each of finitely many adjoining intervals.

\subsection{Parameters and geometry on the unbounded interval}
Set
\[
 D_0=5\,760\,000,\quad X=(D-3)/4,\quad X_0=(D_0-3)/4,
 \quad K=31,\quad c_R=23/10,\quad e=29/500.
\]
For $D\geq D_0$ choose
\begin{gather*}
 S=\lfloor c_RX/K\rfloor,\quad R=KS,\quad J=\lfloor eR\rfloor,\\
 T=106\,838,\quad R_0=KT=3\,311\,978,\quad J_0=192\,094.
\end{gather*}
Thus $S\geq T$, $R\geq R_0$ and $J\geq J_0$.
Write $J_*=J_0+1$, $M_0=R_0(T+1)/2$, and put
\begin{align*}
 c_-&=c_R-K/X_0,& e_-&=e-1/R_0,\\
 r_0&=\frac{4\sqrt{1.02K/e_-}}{\sqrt{2\pi T\sqrt{D_0}}}
            +\frac{1+\sqrt{1+4/D_0}}T,&
 \gamma_0&=1+r_0+800e/T.
\end{align*}
These give $\gamma_0<1.002779879<1.02$.
Choose $n$ by rounding $(1+r_0)RS/J$ strictly upward to a multiple of
$800$. Then
\begin{equation}\label{eq:global-node-bounds}
 \frac{RS}{J}<n\leq\gamma_0\frac{RS}{J}.
\end{equation}
The disk bound gives $M_y\leq4\rho/\sqrt D+1$, and substitution of
$n\leq1.02RS/J$ shows $M_y/S\leq r_0$. The selected $n$ therefore
satisfies $nJ>R(S+M_y)$. Indeed, the definition of $n$ implies $n\leq\gamma_0RS/J<1.02RS/J$, which
justifies the radius estimate used to bound $M_y$.

For later use put
\begin{align*}
 \theta_-&=398891129/800000000,&\theta_+&=(1+1/T)/2,\\
 \nu_-&=c_-/(Ke),&\nu_+&=\gamma_0c_R/(Ke_-),\\
 \eta_+&=\frac{\sqrt{D_0+4}/4}
                  {\sqrt{\nu_-X_0\sqrt{D_0}/(2\pi)}},\\
 C_-&=\tfrac12\log(\nu_-/\pi),&
 C_+&=\tfrac12\log\left(\nu_+\sqrt{1+3/(4X_0)}/\pi\right)
                                      +\log(1+\eta_+).
\end{align*}
Equation~\eqref{eq:global-node-bounds} gives
\begin{equation}\label{eq:global-geometry}
 \theta_-\leq\theta\leq\theta_+,\quad
 \nu_-X\leq n\leq\nu_+X,\quad \eta\leq\eta_+,
 \quad \tfrac34\log X+C_-\leq\log\rho\leq\tfrac34\log X+C_+.
\end{equation}
For the first inequality use $\theta_-<1/(2\gamma_0)$.
For the third, the square of the bound for $\eta$ is a constant times
$(D+4)/((D-3)\sqrt D)$, which decreases for $D>4$.
The last inequality follows by writing $D=4X+3$ in the disk formula.
In particular $\eta_+<0.022619737$.
The remaining resultant hypothesis holds as well: $J<R$, and
\[
 \frac{2\rho+1}{R}
 \leq\frac{2e^{C_+}}{c_-}X^{-1/4}+\frac1{c_-X}<1.
\]
The right side decreases and is less than one at $X_0$.
Thus $\max(J,M_x)\leq R+KM_y$, and the resultant criterion supplies
a nonzero maximal minor throughout this interval.

\subsection{A uniform arithmetic lower bound}
Only primes at most $X$ are used in this part of the argument. They are
all inert by Lemma~\ref{lem:small-inert}. We give the lower bound
explicitly, including its dependence on the anchor $T$, so that its
validity as $D$ increases is visible.

For an integer $q\geq4$, put
\[
 E_T(q)=\left[\frac1{Kq}-\frac3{2R_0}+
             \frac3{2R_0q}-\frac{q}{2\sqrt3KT^3}\right]_+.
\]
The floor-sum estimate for $T(S,q)$ gives
$2KT(S,q)/W\geq E_T(q)$ when $S\geq T$.
One can obtain it by expanding the quadratic sum
\eqref{eq:T-floor}; the only signed periodic term is bounded using
$|t(1-t)(1-2t)|\leq1/(6\sqrt3)$, $0\leq t\leq1$.
For $4\leq q\leq S$, the positive term $q/(2R_0^2)$ absorbs the
Bernoulli remainder in the displayed expression; for $q>S$ the sum is zero.
Dropping the remaining nonnegative terms gives the displayed expression.
It increases when $T$ increases, since its terms depending on $T$ are
$-3(1-1/q)/(2KT)-q/(2\sqrt3KT^3)$.

For a prime $p<K$, write $a_p=K-p$ and define
\[
 W_p^{(0)}=\frac{1-T^{-2}}K+
                      \sum_{p^{2l}\leq T}E_T(p^{2l}),
 \qquad C(v)=v^2(3-2v),\quad e_* =J_*/R_0.
\]
For each $v\in\{0,1/200,\ldots,1\}$, put
\[
 v_-=(v-1/T)_+,\qquad q_v=\frac{a_p}{K}v\min(1,v+1/T).
\]
Define the corresponding row gains and column loss by
\begin{align*}
 A_p(v)&=\frac{a_p}{Kp^2}[C(v_-)-3/T]_+
          +\frac{a_p}{Kp^3}[1-C(v)-6/T]_+\\
 &\quad+\min\left(0,\frac{3a_p(K+p-1-2(p^2-1))}{2K^2p^2T}\right)
                  -\frac{3a_pp^2}{2K^2T^2},\\
 B_p(v)&=\left[\frac{a_p}{K^2}
     \left\{\left(\frac1p-\frac1{p^2-1}\right)v_-^3-\frac3{2T}\right\}
                                                        \right]_+,\\
 C_p(v)&=3(1+1/T)q_v(1-q_v\theta_-/2)\frac{p-1}{p}e_*\\
 &\quad+\boldsymbol1_{v>0}
       \frac{3(1+1/T)\lfloor p^2/4\rfloor}
                  {2p\theta_-(J_0/R_0)R_0^2}.
\end{align*}
The resulting small-prime weight is
\[
 W_p=W_p^{(0)}+
       \max\left(0,\max_v\{[A_p(v)+B_p(v)]_+-C_p(v)\}\right).
\]
These are lower bounds for the small-prime row weights in
\eqref{eq:wkeep}. To see the direction of each estimate, take
$L=\lfloor vS\rfloor$ there, so that $v-1/S\leq L/S\leq v$.
The sums of $t(S-t)$ in the two parts of the triangle, after
normalization, have limiting values $C(v)$ and $1-C(v)$; summing
$t$ and $t^2$ bounds their errors by $3/S$ on either part.
The constant term of the strip estimate contributes at least its
negative part. The incomplete-strip loss is at most
$3a_pp^2/(2K^2S^2)$. For the extra odd-power sum, integral comparison
of the floors gives the lower bound in $B_p$.
Finally, the proportion of selected rows is at most $q_v$ and
$q(1-q\theta/2)$ increases for $0\leq q\leq1$; the column estimate
therefore gives $C_p$. Replacing $S$ by $T$ and $\theta$ by $\theta_-$ only weakens these bounds.
For the final incomplete-column remainder, the denominator is bounded using
$JR\geq J_0R_0$, which follows directly from
$J\geq J_0$ and $R\geq R_0$.

Before replacing $S$ by $T$, the corresponding lower gains increase with
$S$, while the losses decrease: $(v-1/S)_+$ increases,
$(a_p/K)v\min(1,v+1/S)$ decreases, and $C$ is increasing on $[0,1]$.
Thus each fixed choice of $v$ gives a uniform lower bound. Taking the
maximum over the grid preserves this bound; any additional tower terms
are nonnegative.

Write $L_0=\log X_0$, $u_-=1/c_R$, $u_+=1/c_-$, and
\[
 E_M=-\gamma_{\rm Euler}-\sum_p\frac{\log p}{p(p-1)}.
\]
The explicit estimates of Rosser and Schoenfeld
\cite[Theorems 4 and 6]{rosserschoenfeld}, followed by partial summation,
give, for $x\geq X_0$,
\begin{align}
 \sum_{p\leq x}\frac{\log p}{p}
       &\geq\log x+E_M-\frac1{2\log x},\notag\\
 \vartheta(x)&\leq x\left(1+\frac1{2\log x}\right),\notag\\
 \sum_{K\leq p\leq x}p\log p
       &\geq\frac{x^2}{2}
         \left(1-\frac3{2\log x}-\frac1{3\log^2x}\right)-160000.
                                                    \label{eq:global-prime-sums}
\end{align}
For the last inequality, integrate the bound for $\vartheta$ from $563$
and use
\[
 \frac{d}{dt}\left(\frac{t^2}{2\log t}
                         +\frac{t^2}{3\log^2t}\right)
 \geq\frac{t}{\log t}\qquad(t\geq563).
\]
The constant $160000$ also covers the omitted initial interval.

The large-prime row weight has the lower bound
\[
 \frac{S_r(R,K,p)}{KW}
 \geq\frac1p-\frac3{2R}+\frac{3K}{2Rp}
            +\frac{p}{2R^2}-\frac{3K}{2R^2}
            -\frac{p^2}{12\sqrt3R^3}.
\]
The two terms containing $3K$ have nonnegative sum when $p\leq X<R$.
Also $p^2\leq Xp$. Applying \eqref{eq:global-prime-sums} therefore
shows that the normalized arithmetic term is at least
\begin{equation}\label{eq:global-arithmetic-growth}
 A_0+\log(X/X_0),
\end{equation}
where
\begin{align}
 A_0={}&\sum_{p<K}W_p\log p
    +\sum_{\substack{p\geq K, l\geq1\\p^{2l}\leq T}}
                                     E_T(p^{2l})\log p\notag\\
 &+L_0+E_M-\frac1{2L_0}-\sum_{p<K}\frac{\log p}{p}
       -\frac3{2c_-}\left(1+\frac1{2L_0}\right)\notag\\
 &+\left(\frac{u_-^2}{2}-\frac{u_+^3}{12\sqrt3}\right)
       \left\{\frac12\left(1-\frac3{2L_0}-\frac1{3L_0^2}\right)
                       -\frac{160000}{X_0^2}\right\}.
                                                    \label{eq:global-A0}
\end{align}
The coefficient on the final line is positive. As $X$ increases, its
braced lower bound increases; the negative logarithmic errors decrease
in modulus. This proves the asserted uniformity in
\eqref{eq:global-arithmetic-growth}. Evaluation gives
\begin{equation}\label{eq:global-A0-value}
 A_0>10.8131594547983.
\end{equation}

\subsection{Uniform bounds for the derivative and radial terms}
Retain only prime powers $q=p^a\leq1000$, $a\geq2$, and put
\[
 \sigma=\sum_{p^a\leq1000,\ a\geq2}\frac{\log p}{p^a},
 \quad\ell_- =\tfrac14\log X_0+\log c_- -C_+-\sigma,
 \quad\ell_+ =\tfrac14\log X_0+\log c_R-C_- -\sigma.
\]
By \eqref{eq:global-geometry}, the actual slope is between
$\ell_-+\frac14\log(X/X_0)$ and
$\ell_++\frac14\log(X/X_0)$. Discarding further factorial savings
gives a valid upper bound for the determinant.

Each residue modulo $q$ can occur in at most $n\lceil J/q\rceil$
selected columns. Dividing by $M$ bounds this capacity by
$a_q=(1+q/J_0)/(\theta_-q)$.
Put $b_q=\min(q,\lfloor1/a_q\rfloor)$ and
$r_q=(1-b_qa_q)_+$. Filling the largest remainders first proves that
\begin{equation}\label{eq:global-column-remainder}
 c_j=2\sum_{q=p^a\leq1000,\ a\geq2}\frac{\log p}{q}
 \left\{a_q\left(b_q(q-1)-\frac{b_q(b_q-1)}2\right)
                       +r_q(q-1-b_q)_+\right\}
       <60.016864219225
\end{equation}
is an admissible bound for the column remainder.

For the factorial saving, use $3200$ intervals $[u,v]$ of equal length
in $[0,1]$. On each interval take
\[
 a=(1-\theta_+(1-u)^2)e_-,\quad
 b=\min(1,(1+1/T)v+1/R_0),\quad
 \alpha=[(1-u)^2-(1-v)^2-1/M_0]_+.
\]
The number of quantiles in this interval is at least $\alpha M$.
The monotonicities following \eqref{eq:radial-rearrangement}, with
$\ell=\ell_-$, therefore give
\begin{equation}\label{eq:global-factorial-saving}
 S_{\rm der}=3\sum_{[u,v]}\alpha h(b,a)
                       >0.1803971924811329636.
\end{equation}
All the inequalities used here improve with $R,S,J$ at their stated
lower bounds; no new derivative orders or prime powers are assumed
to have the same residues as those at the endpoint.

We next put the convex estimate on a fixed domain. Set $N=800$,
$b_0=480$, and $n_0=1\,842\,400$, which is a lower bound for $n$ by
\eqref{eq:global-node-bounds}. In \eqref{eq:radial-bin-weights} use
$\eta_+$ and $n_0$, and denote the resulting numbers by $L_b^0$.
Take
\[
 \overline L_0=\tfrac12\log(N/(b_0+1)),\qquad
 u_0=(J_*+2)/R_0,
\]
and let $\Delta_0$ be \eqref{eq:radial-Delta} evaluated at $(R_0,T,J_*)$.
For every $D\geq D_0$ the actual threshold $(J+2)/R$ is at most
$u_0$. Also $\Delta\leq\Delta_0$: substitute
$(J+1)/R\leq e+1/R_0$ in its quadratic expression, whose derivative
is positive above $1/(2K)$, and compare with $J_*/R_0$.
The resulting strict bound absorbs the omitted negative term $-1/R^2$.

Choose $120$ integers
\[
 m_r=\lfloor.02J_0\rfloor+
 \left\lfloor\frac{r(\lfloor.46J_0\rfloor-\lfloor.02J_0\rfloor)}{119}
                                                      \right\rfloor,
 \quad 0\leq r<120,
\]
and put $\tau_r=m_r/J_0$. Retain the initial segment on which
\[
 G_r=1-\frac{\tau_r+1/J_0}{\theta_-}>0,
 \quad a_r=1-\sqrt{1-G_r}-(1-\tau_r)J_*/R_0>0,
 \quad\nu_r=\ell_-+\log a_r>0.
\]
Use successive differences of these decreasing $\nu_r$, with the next
value after the last retained cut equal to zero, and reduce the cuts
to $\widetilde\tau_r=(\theta_-/\theta_+)\tau_r$.
These choices give lower staircase levels for the cuts
$\lceil\tau_rJ(D)\rceil$ at every $D\geq D_0$; the extra $1/J_0$
accounts for rounding the cuts upward.

In \eqref{eq:radial-functional} replace $\theta$ by $\theta_-$,
the radial coefficients by $L_b^0$, the threshold in $\Phi$ by $u_0$,
and the quadratic coefficient $3\varepsilon/(2\theta)$ by
$3(J_0-1)/(2R_0\theta_-)$. Use the cuts
$\widetilde\tau_r$ and the preceding staircase levels, and replace
the negative correction by $-\overline L_0\Delta_0$.
Call the resulting function $\mathcal F_0$.
For any actual occupancy vector $w$, put
$\widetilde w=(\theta_-/\theta)w$. Then
$\widetilde w\in\mathcal P_{\theta_-}$ and its cumulative fractions
$G_b$ are unchanged. Moreover,
\[
 \frac{\varepsilon}{\theta}w_b^2
 =\frac{\varepsilon\theta}{\theta_-^2}\widetilde w_b^2
 \geq\frac{J_0-1}{R_0\theta_-}\widetilde w_b^2.
\]
The same calculation applies to the positive-part squares, since their
rescaled cuts are at least $\widetilde\tau_r$.
The radial weights and their successive differences decrease when either
$\eta$ or the shift $1/n$ increases; this follows by differentiating
the logarithm in \eqref{eq:radial-bin-weights}.
Finally $\overline L_0$ is an upper bound for the clipping weight.
Each of these replacements therefore weakens the lower bound, and
\begin{equation}\label{eq:global-functional-comparison}
 \mathcal F_D(w)\geq\mathcal F_0(\widetilde w).
\end{equation}

The supporting-plane formula \eqref{eq:radial-supporting-plane}, applied
to $\mathcal F_0$, gives
\begin{equation}\label{eq:global-convex-value}
 \min_{\mathcal P_{\theta_-}}\mathcal F_0
                          >0.0825081841978627281=:F_*.
\end{equation}
The supporting-plane evaluation uses only the explicit function above and
the finite linear minimization in~\eqref{eq:radial-supporting-plane}; the
repository code records the parameter values and evaluates it with
outward-rounded arithmetic.

\subsection{The resulting contradiction}
Put
\[
 A=1+\frac1K+\frac{3(K-1)}{2KT}+\frac{K-5}{2KT^2},\qquad
 B=\frac{3(T+1)J_*}{R_0T}(1-\theta_-/2)
                     +\frac{3\theta_+J_*}{R_0^2T}.
\]
These bound $(W_1+W_2)/W$ and $\gamma$ in
\eqref{eq:radial-derivative-coupled} from above.
Combining the radial determinant bound with the arithmetic column
remainder gives the Archimedean upper bound
\begin{equation}\label{eq:global-size-growth}
 U_0+\tfrac34(A+B)\log(X/X_0),
\end{equation}
where
\begin{align}
 U_0={}&A(\tfrac34\log X_0+C_+)+B\ell_+-F_*-S_{\rm der}\notag\\
 &+\frac{3(T+1)c_j}{2R_0T}
                         +\frac{3(T+1)\log M_0}{R_0T}
       <10.8128420196003.                  \label{eq:global-U0}
\end{align}
The coefficient of $c_j$ decreases with $S$. So does the last term:
$((S+1)/S^2)\log(KS(S+1)/2)$ has negative derivative for $S\geq T$,
as follows by logarithmic differentiation and $\log(KS(S+1)/2)>2$.
The geometry grows by at most $3A\log(X/X_0)/4$, while the derivative
slope grows by at most $B\log(X/X_0)/4$.
The larger coefficient in \eqref{eq:global-size-growth} is thus valid.

Subtracting \eqref{eq:global-size-growth} from
\eqref{eq:global-arithmetic-growth} yields
\begin{equation}\label{eq:tail-positive}
 G_{\rm tail}(D)>0.00031743+0.127829\log(X/X_0)>0.
\end{equation}
Here $A_0-U_0>0.0003174351980$ and
$1-3(A+B)/4>0.1278297612987827$.
The nonzero minor supplied above must satisfy the opposite determinant
inequality. This excludes every $D\geq D_0$.

\subsection{The finite interval below \texorpdfstring{$D_0$}{D0}}
For $875\leq D\leq D_0$, keep $P,K,S,R,J,n$ fixed on each closed
interval $[d_-,d_+]$. The lattice radius is enclosed using its endpoint
values. Choose $n$ to be a multiple of $800$ satisfying the resultant
conditions with the upper radius and the lower value of $D$.
The arithmetic estimates use norm representations when $P>(d_--3)/4$;
otherwise every retained prime is inert.

In all cases use $N=800$ radial groups, $3200$ groups for the factorial
saving, and $120$ cuts in \eqref{eq:radial-derivative-coupled}.
The latter run from $\max(1,\lfloor.02J\rfloor)$ to
\[
 \max\left(\max(1,\lfloor.02J\rfloor)+1,
       \min(\lfloor M/n\rfloor-1,\lfloor.495J\rfloor)\right),
\]
equally spaced with integer rounding, retaining the initial segment
where the stated positivity conditions hold. Take $b_0=480$.
The prime-power cutoff is $J-1$.

The intervals using norm representations have the parameter families
in Table~\ref{tab:global-hybrid-parameters}; there
$S=\lfloor c_RP/K\rfloor$, $R=KS$, $J=\lfloor eR\rfloor$.
Their slope windows have cutoff at most $750000$ in the first row and
$140000$ thereafter. The prime-product bounds in Appendix~\ref{app:arithmetic}
apply on each constituent interval.
\begin{table}[ht]
\centering\small
\begin{tabular}{lrrrr}
\toprule
Range & $P$ & $K$ & $c_R$ & $e$\\
\midrule
$[875,1167]$ & $8\,999\,999$ & $23$ & $7/5$ & $7/125$\\
$[1167,1359]$ & $1\,699\,999$ & $23$ & $8/5$ & $27/500$\\
$[1359,1871]$ & $1\,699\,999$ & $23$ & $9/5$ & $27/500$\\
$[1871,6959]$ & $1\,699\,999$ & $25$ & $2$ & $7/125$\\
$[6959,319591]$ & $1\,699\,999$ & $29$ & $11/5$ & $29/500$\\
$[319591,3968539]$ & $1\,699\,999$ & $31$ & $23/10$ & $29/500$\\
\bottomrule
\end{tabular}
\caption{Parameter families for the norm-form comparisons. Each range is
subdivided into closed intervals on which all integer parameters are fixed.}
\label{tab:global-hybrid-parameters}
\end{table}

For the remaining intervals take $K=31$, $c_R=23/10$, $e=29/500$, and
$P=\lfloor(d_--3)/4\rfloor$. Between $3968539$ and $4894699$ the prime
sums are evaluated directly. Above $4894699$ use the explicit bounds of
\cite[Theorems 4, 6 and 18]{rosserschoenfeld}; the additional estimate
$\vartheta(x)<x$ is used only for $x\leq10^8$, within its stated range.
On each interval the radial function is bounded below with the upper
value of $\eta$, while the negative clipping term uses the lower value.
The derivative scale likewise uses the appropriate endpoint in each
term. These substitutions bound the entire interval.

The adjoining intervals cover $[875,D_0]$, and their lower bounds are
summarized in Table~\ref{tab:coverage}. Each is positive. Together with
\eqref{eq:tail-positive}, this excludes every $D\geq875$.
The interval endpoints and the numerical calculations are available with the
\href{https://github.com/dualverse-ai/tenth_discriminant/tree/main/bound-498/code}{accompanying code}.

\Needspace{18\baselineskip}
\section*{AI use}\label{sec:ai-use}
The research for this paper was carried out with autonomous rStation agents
using Gemini 4 Argon and GPT-6 Astra. The author launched these agents, which
conducted the research without human intervention. The initial arguments
were recorded as
\href{https://rstation.ai/problems/tenth_discriminant/papers/2/v12}{paper 2,
version 12} (under the agent name Severi Jet), giving the bound $499$, and
\href{https://rstation.ai/problems/tenth_discriminant/papers/22/v1}{paper 22,
version 1} (under the agent name Mira Cairn), giving $498$.
The local arithmetic, resultant and global estimates used here are in
Sections 3--6 of the former; the spectral estimates are developed in
\href{https://rstation.ai/problems/tenth_discriminant/papers/16/v2}{paper 16,
version 2}, Sections 2--5 and 6.4. The necessary derivations are included in the appendices. The proof was subsequently checked externally by separate
AI agents using GPT-6 Astra and Claude Opus 5.5, run independently of the
agents that made the discoveries.

GPT-6 Astra, Claude Opus 5.5 and Gemini 4 Argon were then used together to
simplify the proof and prepare this exposition. This included revisions of
the outline and individual sections, mathematical and editorial review,
and recomputation of the numerical bounds. In particular, narrowing the
auxiliary-prime interval gives the shorter rank argument proved in
Section~\ref{sec:rank}.

\begin{thebibliography}{99}
\raggedright

\bibitem{baker}
Alan Baker, \emph{Linear forms in the logarithms of algebraic numbers},
Mathematika \textbf{13} (1966), no.~2, 204--216.
\href{https://doi.org/10.1112/S0025579300003971}{doi:10.1112/S0025579300003971}.

\bibitem{green2026}
Ben Green, \emph{On Landau--Siegel zeros},
Hexagon preprint 2610.00311, version 2, 18 October 2026.
\url{https://hexagonmath.org/pdf/2610.00311v2}.

\bibitem{dlmf}
NIST Digital Library of Mathematical Functions, §§18.19 and 18.23,
\url{https://dlmf.nist.gov/18.19} and \url{https://dlmf.nist.gov/18.23}.

\bibitem{forresterwarnaar}
P. J. Forrester and S. O. Warnaar, \emph{The importance of the Selberg integral},
Bull. Amer. Math. Soc. (N.S.) \textbf{45} (2008), 489--534.
\href{https://doi.org/10.1090/S0273-0979-08-01221-4}{doi:10.1090/S0273-0979-08-01221-4}.

\bibitem{fulton}
William Fulton, \emph{Young Tableaux}, London Mathematical Society Student Texts,
vol.~35, Cambridge University Press, 1997.
\href{https://doi.org/10.1017/CBO9780511626241}{doi:10.1017/CBO9780511626241}.

\bibitem{heegner}
Kurt Heegner, \emph{Diophantische Analysis und Modulfunktionen},
Math. Z. \textbf{56} (1952), 227--253.
\href{https://doi.org/10.1007/BF01174749}{doi:10.1007/BF01174749}.

\bibitem{rosserschoenfeld}
J. Barkley Rosser and Lowell Schoenfeld,
\emph{Approximate formulas for some functions of prime numbers},
Illinois J. Math. \textbf{6} (1962), 64--94.
\href{https://doi.org/10.1215/ijm/1028998657}{doi:10.1215/ijm/1028998657}.

\bibitem{stark}
Harold M. Stark, \emph{A complete determination of the complex quadratic fields of class-number one},
Michigan Math. J. \textbf{14} (1967), 1--27.
\href{https://doi.org/10.1307/mmj/1028999653}{doi:10.1307/mmj/1028999653}.

\bibitem{taosketch}
Terence Tao, \href{https://terrytao.wordpress.com/2026/10/19/a-sketch-of-the-proof-of-non-existence-of-siegel-zeroes/}{\emph{A sketch of the proof of non-existence of Siegel zeroes}},
\emph{What's New}, 19 October 2026.

\bibitem{tao2026}
Terence Tao, \href{https://terrytao.wordpress.com/2026/10/20/an-experimental-ai-math-challenge-using-tenth-discriminant-bounds-to-explore-arguments-ruling-out-siegel-zeroes/}{\emph{An experimental AI math challenge: using tenth discriminant bounds to explore arguments ruling out Siegel zeroes}},
\emph{What's New}, 20 October 2026.

\end{thebibliography}

\end{document}
